% Les transformations de l’élevage et de la pêche traditionnels — Géographie Tle C — Cours signé OnBuch+
% Programme officiel MINESEC. Compiler avec : tectonic N05-transformations-elevage-peche-traditionnels.tex
\newcommand{\DOCMATIERE}{Géographie}
\newcommand{\DOCNIVEAU}{Tle C}
\newcommand{\DOCMODULE}{Module 1 : Le Cameroun}
\newcommand{\DOCLECON}{N05}
\newcommand{\DOCTITRE}{Les transformations de l’élevage et de la pêche traditionnels}
% OnBuch+ — préambule « cours signés » ENRICHI (compilé avec tectonic / XeLaTeX)
% Système visuel « Pop » d'OnBuch+ : fond crème (jamais blanc), bordures encre
% épaisses, ombres « dures » décalées, titres Archivo Black en capitales.
% Version MATHÉMATIQUES : graphes (pgfplots/tikz), boîtes pédagogiques
% (définition, propriété, méthode, attention, le savais-tu, exercice résolu,
% exercice, TP, bilan), page de garde et signature OnBuch+ ancrée en bas.
%
% Macros attendues dans main.tex AVANT \input{preamble} :
%   \DOCMATIERE  \DOCNIVEAU  \DOCMODULE  \DOCLECON (numéro)  \DOCTITRE
\documentclass[11pt]{article}

\usepackage{fontspec}
\usepackage[french]{babel}
\usepackage[a4paper,margin=2cm,top=3.4cm,bottom=2.8cm,headheight=1.4cm,footskip=1.3cm]{geometry}
\usepackage{amsmath,amssymb,amsfonts}
\usepackage{mathtools} % \xmapsto, \coloneqq... souvent utilisés par les modèles
\usepackage{cancel} % simplifications barrées (bilans de forces)
% \abs{z} (module, valeur absolue) et \norm{u} (norme), utilisés par les modèles
\providecommand{\abs}{}\let\abs\relax\DeclarePairedDelimiter{\abs}{\lvert}{\rvert}
\providecommand{\norm}{}\let\norm\relax\DeclarePairedDelimiter{\norm}{\lVert}{\rVert}
\usepackage[table]{xcolor} % [table] : fournit aussi \rowcolors (alternance de lignes)
\usepackage{pagecolor}
\usepackage{tikz}
\usetikzlibrary{arrows.meta,positioning,calc,shapes.geometric,shapes.misc,shapes.arrows,decorations.pathmorphing,decorations.pathreplacing,decorations.markings,patterns,shadows,mindmap,trees,fit,backgrounds,matrix,intersections,angles,quotes}
\usepackage{pgfplots}
\usepackage[version=4]{mhchem} % équations nucléaires \ce{^{235}_{92}U}
\usepackage[european,siunitx]{circuitikz} % schémas électriques
\pgfplotsset{compat=1.18}
\usepgfplotslibrary{fillbetween,groupplots} % aires entre courbes (calcul intégral)
\usepackage{enumitem}
\usepackage{fancyhdr}
% [most] charge toutes les bibliothèques tcolorbox (listings, minted, raster,
% documentation...) et gonfle inutilement la mémoire TeX sur un document long
% (~300 boîtes) : on ne charge que ce qui est réellement utilisé.
\usepackage[skins,breakable]{tcolorbox}
\usepackage{graphicx}
\graphicspath{{/home/user/t-t/geo-onbuch/images/}}
\usepackage{colortbl}
\usepackage{booktabs}
\usepackage{tabularx}
\usepackage{multirow}
\usepackage{diagbox} % case d'en-tête diagonale des tableaux à double entrée
% Colonnes extensibles centrée / gauche / droite (C, L, R), souvent utilisées par les modèles
\newcolumntype{C}{>{\centering\arraybackslash}X}
\newcolumntype{L}{>{\raggedright\arraybackslash}X}
\newcolumntype{R}{>{\raggedleft\arraybackslash}X}
\usepackage{array}
\usepackage{multicol}
\usepackage{siunitx}
% parse-numbers=false : accepte aussi \SI{2,5 \times 10^{-3}}{...}
\sisetup{locale=FR,output-decimal-marker={,},group-minimum-digits=4,parse-numbers=false}
\DeclareSIUnit\ohm{\ensuremath{\symup{\Omega}}} % U+2126 absent de la police
\DeclareSIUnit\division{div} % réglages d'oscilloscope (ms/div, V/div)
\DeclareSIUnit\tr{tr} % tours (vitesse de rotation, ex. \SI{1500}{\tr\per\minute})
\DeclareSIUnit\molar{M} % concentration molaire (ex. \SI{3}{\molar}, \SI{1}{\micro\molar})
\DeclareSIUnit\an{an} % année (le modèle confond parfois avec \year, un compteur TeX) — ex. \SI{5}{\kilo\meter\per\an}
\DeclareSIUnit\FCFA{FCFA} % franc CFA (ex. \SI{500}{\FCFA\per\kilo\gram})
\DeclareSIUnit\degreeBrix{\textdegree Brix} % teneur en sucre d'un sirop/jus (réfractométrie)
\DeclareSIUnit\month{mois} % le modèle utilise parfois \month, un compteur TeX, comme unité de durée
\DeclareSIUnit\microliter{\micro\liter} % le modèle écrit parfois \microliter en un seul token
\DeclareSIUnit\million{millions} % dénombrement (ex. \SI{15}{\million\per\mL} spermatozoïdes)
\usepackage{qrcode}
% \degree : commande fréquemment utilisée par les modèles pour un angle ou une
% température en dehors de \SI{}{} (non fournie par les paquets chargés).
\providecommand{\degree}{^{\circ}}
% \qed (fin de démonstration), utilisé par les modèles sans amsthm
\providecommand{\qed}{\hfill$\square$}
% \barre : double barre des valeurs interdites dans un tableau de variations (tkz-tab)
\providecommand{\barre}{\ensuremath{\|}}
% environnement proof (amsthm non chargé) utilisé par les modèles
\ifdefined\proof\else\newenvironment{proof}[1][Démonstration]{\par\noindent\textit{#1.}\ }{\qed\par}\fi
\usepackage{lastpage}
\usepackage{hyperref}
\hypersetup{hidelinks}

% ============================================================
% Palette « Pop » OnBuch+ — mêmes hex que la classe P (plus_ui.dart)
% ============================================================
\definecolor{poporange}{HTML}{F59321}
\definecolor{popdark}{HTML}{E07A0C}
\definecolor{popcream}{HTML}{FFF6E8}
\definecolor{popink}{HTML}{1C1714}
\definecolor{poppurple}{HTML}{7A5AE0}
\definecolor{popgreen}{HTML}{2E9E6A}
\definecolor{poppink}{HTML}{E0457B}
\definecolor{popblue}{HTML}{2F7DE1}
\definecolor{popgold}{HTML}{C98A12}
\definecolor{popmuted}{HTML}{8A7F76}
\definecolor{popline}{HTML}{EADFD2}
\definecolor{popgray}{HTML}{8A7F76}
\colorlet{popgrey}{popgray}
% Alias tolérants (le modèle nomme parfois les couleurs différemment)
\colorlet{popwhite}{white}
\colorlet{popblack}{popink}
\colorlet{popred}{poppink}
\colorlet{popyellow}{popgold}
% Teintes claires (fonds de tableaux/encadrés) — le modèle nomme parfois
% les couleurs avec un suffixe L (ex. popblueL) sans les avoir définies.
\colorlet{poporangeL}{poporange!20}
\colorlet{popdarkL}{popdark!20}
\colorlet{poppurpleL}{poppurple!20}
\colorlet{popgreenL}{popgreen!20}
\colorlet{poppinkL}{poppink!20}
\colorlet{popblueL}{popblue!20}
\colorlet{popgoldL}{popgold!20}
\colorlet{popredL}{popred!20}
\colorlet{popgrayL}{popgray!20}
% Teintes claires pour fonds de tableaux / zones
\colorlet{popblueL}{popblue!10!white}
\colorlet{poporangeL}{poporange!14!white}
\colorlet{popgreenL}{popgreen!12!white}
\colorlet{poppurpleL}{poppurple!12!white}
\colorlet{poppinkL}{poppink!12!white}
\colorlet{popgoldL}{popgold!14!white}

\pagecolor{popcream}

% ============================================================
% Typographie
% ============================================================
\defaultfontfeatures{Ligatures=TeX}
\setmainfont{PlusJakartaSans-Medium.ttf}[Path=../../../fonts/,BoldFont=PlusJakartaSans-Bold.ttf]
% Maths en sans-serif (Fira Math) avec chiffres et lettres droites de Plus
% Jakarta, pour que les formules se fondent dans le texte.
\usepackage[math-style=ISO,bold-style=ISO]{unicode-math}
\setmathfont{FiraMath-Regular.otf}[Path=../../../fonts/]
\setmathfont{PlusJakartaSans-Medium.ttf}[Path=../../../fonts/,range={up/num,up/{Latin,latin},\mathup}]
\usepackage{icomma} % virgule décimale sans espace en mode math
% Caractères mathématiques Unicode absents des polices de texte (Plus Jakarta,
% Archivo Black) : rendus par leur équivalent mathématique, en texte comme en
% formule — sinon ils s'affichent en carré vide (ex. « Arithmétique dans ℤ »).
\usepackage{newunicodechar}
\newunicodechar{ℝ}{\ensuremath{\mathbb{R}}}
\newunicodechar{ℤ}{\ensuremath{\mathbb{Z}}}
\newunicodechar{ℂ}{\ensuremath{\mathbb{C}}}
\newunicodechar{ℕ}{\ensuremath{\mathbb{N}}}
\newunicodechar{ℚ}{\ensuremath{\mathbb{Q}}}
\newunicodechar{𝔻}{\ensuremath{\mathbb{D}}}
\newunicodechar{α}{\ensuremath{\alpha}}
\newunicodechar{β}{\ensuremath{\beta}}
\newunicodechar{γ}{\ensuremath{\gamma}}
\newunicodechar{δ}{\ensuremath{\delta}}
\newunicodechar{ε}{\ensuremath{\varepsilon}}
\newunicodechar{θ}{\ensuremath{\theta}}
\newunicodechar{λ}{\ensuremath{\lambda}}
\newunicodechar{μ}{\ensuremath{\mu}}
\newunicodechar{σ}{\ensuremath{\sigma}}
\newunicodechar{τ}{\ensuremath{\tau}}
\newunicodechar{φ}{\ensuremath{\varphi}}
\newunicodechar{ψ}{\ensuremath{\psi}}
\newunicodechar{ω}{\ensuremath{\omega}}
\newunicodechar{Σ}{\ensuremath{\Sigma}}
% majuscules produites par \MakeUppercase dans les titres (α-aminés -> Α)
\newunicodechar{Α}{\ensuremath{\alpha}}
\newunicodechar{Β}{\ensuremath{\beta}}
\newunicodechar{Γ}{\ensuremath{\Gamma}}
\newunicodechar{Δ}{\ensuremath{\Delta}}
\newunicodechar{Ε}{\ensuremath{\varepsilon}}
\newunicodechar{Θ}{\ensuremath{\Theta}}
\newunicodechar{Λ}{\ensuremath{\Lambda}}
\newunicodechar{Μ}{\ensuremath{\mu}}
\newunicodechar{Τ}{\ensuremath{\tau}}
\newunicodechar{Φ}{\ensuremath{\Phi}}
\newunicodechar{Ψ}{\ensuremath{\Psi}}
\newunicodechar{Ω}{\ensuremath{\Omega}}
\newunicodechar{↦}{\ensuremath{\mapsto}}
\newunicodechar{∈}{\ensuremath{\in}}
\newunicodechar{∉}{\ensuremath{\notin}}
\newunicodechar{∧}{\ensuremath{\wedge}}
\newunicodechar{⇔}{\ensuremath{\Leftrightarrow}}
\newunicodechar{⟺}{\ensuremath{\Longleftrightarrow}}
\newunicodechar{⇒}{\ensuremath{\Rightarrow}}
\newunicodechar{⟹}{\ensuremath{\Longrightarrow}}
\newunicodechar{∘}{\ensuremath{\circ}}
\newunicodechar{∪}{\ensuremath{\cup}}
\newunicodechar{∩}{\ensuremath{\cap}}
\newunicodechar{≡}{\ensuremath{\equiv}}
\newunicodechar{⊂}{\ensuremath{\subset}}
\newunicodechar{⊥}{\ensuremath{\perp}}
\newunicodechar{∃}{\ensuremath{\exists}}
\newunicodechar{∀}{\ensuremath{\forall}}
\newunicodechar{∛}{\ensuremath{\sqrt[3]{\,}}}
\newunicodechar{∜}{\ensuremath{\sqrt[4]{\,}}}
\newunicodechar{⃗}{\ensuremath{{}^{\to}}}
\newunicodechar{ⁿ}{\ifmmode^{n}\else\textsuperscript{\ensuremath{n}}\fi}
\newunicodechar{ˣ}{\ifmmode^{x}\else\textsuperscript{\ensuremath{x}}\fi}
\newunicodechar{ᵇ}{\ifmmode^{b}\else\textsuperscript{\ensuremath{b}}\fi}
\newunicodechar{ʳ}{\ifmmode^{r}\else\textsuperscript{\ensuremath{r}}\fi}
\newunicodechar{ᵘ}{\ifmmode^{u}\else\textsuperscript{\ensuremath{u}}\fi}
\newunicodechar{ʸ}{\ifmmode^{y}\else\textsuperscript{\ensuremath{y}}\fi}
\newunicodechar{ᵃ}{\ifmmode^{a}\else\textsuperscript{\ensuremath{a}}\fi}
\newunicodechar{ᵉ}{\ifmmode^{e}\else\textsuperscript{\ensuremath{e}}\fi}
\newunicodechar{ᵏ}{\ifmmode^{k}\else\textsuperscript{\ensuremath{k}}\fi}
\newunicodechar{ᵖ}{\ifmmode^{p}\else\textsuperscript{\ensuremath{p}}\fi}
\newunicodechar{ᴺ}{\ifmmode^{N}\else\textsuperscript{\ensuremath{N}}\fi}
\newunicodechar{ᶜ}{\ifmmode^{c}\else\textsuperscript{\ensuremath{c}}\fi}
\newunicodechar{ʰ}{\ifmmode^{h}\else\textsuperscript{\ensuremath{h}}\fi}
\newunicodechar{ᵠ}{\ifmmode^{q}\else\textsuperscript{\ensuremath{q}}\fi}
\newunicodechar{ₙ}{\ifmmode_{n}\else\textsubscript{\ensuremath{n}}\fi}
\newunicodechar{ₐ}{\ifmmode_{a}\else\textsubscript{\ensuremath{a}}\fi}
\newunicodechar{ᵢ}{\ifmmode_{i}\else\textsubscript{\ensuremath{i}}\fi}
\newunicodechar{ᵣ}{\ifmmode_{r}\else\textsubscript{\ensuremath{r}}\fi}
\newunicodechar{ₖ}{\ifmmode_{k}\else\textsubscript{\ensuremath{k}}\fi}
\newunicodechar{ᵦ}{\ifmmode_{\beta}\else\textsubscript{\ensuremath{\beta}}\fi}
\newunicodechar{ₓ}{\ifmmode_{x}\else\textsubscript{\ensuremath{x}}\fi}
\newfontfamily\popdisplay{ArchivoBlack-Regular.ttf}[Path=../../../fonts/]
\newfontfamily\poplogo{SpaceGrotesk-Bold.ttf}[Path=../../../fonts/]
\newfontfamily\popsemi{PlusJakartaSans-SemiBold.ttf}[Path=../../../fonts/]
\newfontfamily\popxbold{PlusJakartaSans-ExtraBold.ttf}[Path=../../../fonts/]
\newfontfamily\popxboldls{PlusJakartaSans-ExtraBold.ttf}[Path=../../../fonts/,LetterSpace=3.0]

\setlist[enumerate]{leftmargin=1.7em,itemsep=5pt,topsep=4pt}
\setlist[itemize]{leftmargin=1.7em,itemsep=4pt,topsep=4pt,label={\color{poporange}\textbullet}}
\setlength{\parindent}{0pt}
\setlength{\parskip}{5pt}
\renewcommand{\arraystretch}{1.25}

% pgfplots : style Pop par défaut
\pgfplotsset{
  every axis/.append style={axis line style={popink,line width=1pt},tick style={popink},
    grid style={popline,line width=0.5pt},label style={font=\small},tick label style={font=\footnotesize}},
  popcurve/.style={poporange,line width=1.8pt,no markers,smooth},
}
% Même style disponible dans un \draw TikZ simple (hors pgfplots)
\tikzset{popcurve/.style={poporange,line width=1.8pt}}
\tikzset{popgrid/.style={popline,very thin}}
\colorlet{popcurve}{poporange} % le nom du style est parfois utilisé comme couleur

% ============================================================
% Pill / squiggle
% ============================================================
\newcommand{\poppill}[3]{%
  \tikz[baseline=-0.55ex]\node[draw=popink,line width=1.2pt,rounded corners=2.6pt,%
    fill=#2,inner xsep=6.5pt,inner ysep=3pt]{%
    \popxboldls\fontsize{7}{7}\selectfont\color{#3}\MakeUppercase{#1}};%
}
\newcommand{\popsquiggle}[1]{%
  \tikz[baseline=-0.4ex]{
    \draw[#1,line width=2.6pt,line cap=round]
      plot [smooth] coordinates {(0,0.09) (0.34,0.01) (0.68,0.21) (1.02,0.01) (1.36,0.10)};
  }%
}
% Mot-clé surligné (marqueur orange sous le texte)
\newcommand{\cle}[1]{{\popxbold\color{popdark}#1}}

% ============================================================
% En-tête / pied
% ============================================================
\pagestyle{fancy}
\fancyhf{}
\renewcommand{\headrulewidth}{0pt}
\renewcommand{\footrulewidth}{0pt}
\lhead{\raisebox{-0.15ex}{\poplogo\fontsize{13}{13}\selectfont\color{popink}OnBuch\color{poporange}+}}
\rhead{\poppill{\DOCMATIERE\ \textbullet\ \DOCNIVEAU\ \textbullet\ Leçon \DOCLECON}{white}{popink}}
\lfoot{\popxbold\fontsize{7.6}{7.6}\selectfont\color{popmuted}\MakeUppercase{Cours signé OnBuch+}\ \textbullet\ {\popsemi\DOCTITRE}}
\rfoot{\tikz[baseline=-0.6ex]\node[rounded corners=8.5pt,fill=popink,minimum height=17pt,minimum width=17pt,inner xsep=5pt,inner sep=0]{\popxbold\fontsize{8}{8}\selectfont\color{popcream}\thepage\,/\,\pageref*{LastPage}};}

% ============================================================
% Titres
% ============================================================
\newcommand{\chaptitre}[2]{%
  \begin{tcolorbox}[enhanced,colback=poporange,colframe=popink,boxrule=2.6pt,arc=11pt,
      drop shadow=popink,left=15pt,right=15pt,top=12pt,bottom=11pt,before skip=3pt,after skip=18pt]
    \poppill{#2}{popink}{popcream}\par\vspace{9pt}
    {\raggedright\hyphenpenalty=10000\exhyphenpenalty=10000\popdisplay\fontsize{22}{24}\selectfont\color{popcream}\MakeUppercase{#1}\par}
  \end{tcolorbox}
}
% Section numérotée : pastille orange avec le numéro + titre Archivo
\newcounter{popsec}
\newcommand{\coursec}[1]{%
  \stepcounter{popsec}\setcounter{popsub}{0}%
  \par\vspace{16pt}\noindent
  \begin{minipage}{\linewidth}
  \tikz[baseline=-0.7ex]\node[draw=popink,line width=1.6pt,fill=poporange,rounded corners=4pt,minimum size=22pt,inner sep=0]{\popdisplay\fontsize{12}{12}\selectfont\color{popcream}\thepopsec};%
  \hspace{8pt}{\popdisplay\fontsize{15}{17}\selectfont\color{popink}\MakeUppercase{#1}}\par
  \vspace{4pt}\noindent{\color{poporange}\rule{36pt}{3.6pt}}
  \end{minipage}\par\nopagebreak\vspace{8pt}
}
\newcounter{popsub}
\newcommand{\courssub}[1]{%
  \stepcounter{popsub}%
  \par\vspace{9pt}\noindent{\popxbold\fontsize{11.5}{13}\selectfont\color{popink}{\color{poporange}\thepopsec.\thepopsub}\hspace{6pt}#1}\par\nopagebreak\vspace{4pt}
}

% ============================================================
% Boîtes pédagogiques — même recette : fond blanc, bordure épaisse colorée,
% ombre dure encre, badge plein. Titre optionnel [#1].
% ============================================================
\tcbset{popbase/.style={breakable,enhanced,colback=white,boxrule=2.3pt,arc=9pt,
  shadow={3pt}{-3pt}{0pt}{popink},left=14pt,right=14pt,top=11pt,bottom=11pt,before skip=11pt,after skip=15pt,
  fontupper=\popsemi\fontsize{10.6}{14.5}\selectfont\color{popink},
  fontlower=\popsemi\fontsize{10.6}{14.5}\selectfont\color{popink}}}
% Coche dessinée (le glyphe ✓ n'existe pas dans les polices du document)
\newcommand{\popcheck}[1]{\tikz[baseline=-0.5ex]\draw[#1,line width=1.6pt,line cap=round,line join=round](-0.13,0.0)--(-0.04,-0.09)--(0.13,0.1);}
\providecommand{\coche}{\popcheck{popgreen}}
\providecommand{\resultat}[1]{\par\smallskip\noindent\fcolorbox{popgreen}{popgreenL}{\parbox{\dimexpr\linewidth-2\fboxsep-2\fboxrule\relax}{\textbf{Résultat.} #1}}\par\smallskip}
\newcommand{\popboxhead}[3]{\poppill{#1}{#2}{white}\ifx\relax#3\relax\else\hspace{6pt}{\popxbold\fontsize{10.6}{12}\selectfont\color{popink}#3}\fi\par\vspace{7pt}}

\newtcolorbox{aretenir}[1][]{popbase,colframe=popgreen,before upper={\popboxhead{À retenir}{popgreen}{#1}}}
\newtcolorbox{exemplebox}[1][]{popbase,colframe=poporange,before upper={\popboxhead{Exemple}{poporange}{#1}}}
\newtcolorbox{definition}[1][]{popbase,colframe=popblue,before upper={\popboxhead{Définition}{popblue}{#1}}}
\newtcolorbox{propriete}[1][]{popbase,colframe=poppurple,before upper={\popboxhead{Propriété}{poppurple}{#1}}}
\newtcolorbox{methode}[1][]{popbase,colframe=popgold,before upper={\popboxhead{Méthode}{popgold}{#1}}}
\newtcolorbox{attention}[1][]{popbase,colframe=poppink,before upper={\popboxhead{Attention}{poppink}{#1}}}
\newtcolorbox{experience}[1][]{popbase,colframe=popblue,colback=popblueL,before upper={\popboxhead{Expérience}{popblue}{#1}}}
% « Le savais-tu ? » : carte violette pleine (contexte, culture, Cameroun)
\newtcolorbox{savaistu}[1][]{popbase,colback=poppurple,colframe=popink,
  fontupper=\popsemi\fontsize{10.4}{14.2}\selectfont\color{white},
  before upper={\poppill{Le savais-tu\,?}{white}{poppurple}\ifx\relax#1\relax\else\hspace{6pt}{\popxbold\color{white}#1}\fi\par\vspace{7pt}}}
% Exercice résolu : énoncé en haut, \tcblower puis la solution sur fond crème
\newcounter{popexo}\newcounter{popexr}
\newtcolorbox{exoresolu}[1][]{popbase,colframe=poporange,colbacklower=popcream,
  segmentation style={popink,line width=1.2pt,dashed},
  before upper={\stepcounter{popexr}\popboxhead{Exercice résolu \thepopexr}{poporange}{#1}},
  before lower={\poppill{Solution}{popink}{popcream}\par\vspace{6pt}}}
% Exercice (non résolu en place) — la correction est dans la partie Corrigés
\newtcolorbox{exercice}[1][]{popbase,colframe=popink,
  before upper={\stepcounter{popexo}\popboxhead{Exercice \thepopexo}{popink}{#1}}}
% Corrigé
\newtcolorbox{corrige}[1][]{popbase,colframe=popgreen,colback=popgreenL,
  before upper={\popboxhead{Corrigé}{popgreen}{#1}}}
% Objectifs / prérequis / bilan
\newtcolorbox{objectifs}{popbase,colframe=popink,colback=poporangeL,
  before upper={\popboxhead{Objectifs de la leçon}{popink}{}}}
\newtcolorbox{prerequis}{popbase,colframe=popink,colback=popblueL,
  before upper={\popboxhead{Prérequis}{popblue}{}}}
\newtcolorbox{bilan}{popbase,colframe=popink,colback=white,boxrule=2.8pt,
  before upper={\popboxhead{Fiche bilan}{poporange}{L'essentiel à connaître pour l'examen}}}
% Figure encadrée avec légende
\newtcolorbox{popfigure}[1][]{enhanced,breakable=false,colback=white,colframe=popline,boxrule=1.4pt,arc=8pt,
  left=10pt,right=10pt,top=10pt,bottom=8pt,before skip=10pt,after skip=12pt,halign=center,
  lower separated=false,halign lower=center,
  fontlower=\popsemi\fontsize{9}{11.5}\selectfont\color{popmuted},
  #1}
\newcommand{\legende}[1]{\par\vspace{6pt}{\popsemi\fontsize{9}{11.5}\selectfont\color{popmuted}#1}}

% ============================================================
% Signature OnBuch+
% ============================================================
\newcommand{\poplogoinline}[1]{{\poplogo\fontsize{#1}{#1}\selectfont\color{popink}OnBuch\color{poporange}+}}
\newcommand{\signatureonbuch}{%
  \begin{tcolorbox}[enhanced,colback=popink,colframe=popink,boxrule=2.2pt,arc=10pt,
      drop shadow=poporange,left=14pt,right=14pt,top=10pt,bottom=10pt,before skip=0pt,after skip=0pt]
    \tikz[baseline=-0.7ex]\node[circle,fill=popgreen,draw=popcream,line width=1.3pt,minimum size=22pt,inner sep=0]{\popcheck{white}};%
    \hspace{9pt}\begin{minipage}[c]{0.62\linewidth}
      {\popxbold\fontsize{12}{13}\selectfont\color{popcream}Signé\ \textbullet\ {\poplogo OnBuch\color{poporange}+}}\par\vspace{2pt}
      {\raggedright\popsemi\fontsize{8}{10.5}\selectfont\color{popline}\MakeUppercase{Équipe pédagogique OnBuch+}\\\MakeUppercase{Conforme au programme officiel MINESEC}\par}
    \end{minipage}\hfill
    {\poplogo\fontsize{22}{22}\selectfont\color{popcream}OnBuch\color{poporange}+}
  \end{tcolorbox}%
}

% ============================================================
% Page de garde
% #1 titre · #2 matière · #3 niveau · #4 module · #5 n° leçon · #6 durée ·
% #7 description / objectifs (texte ou itemize)
% ============================================================
\newcommand{\pagegarde}[7]{%
  \thispagestyle{empty}
  \newgeometry{a4paper,margin=1.7cm,top=1.6cm,bottom=1.4cm}%
  \begin{tikzpicture}[remember picture,overlay]
    \fill[poporange!18] ([xshift=-3.2cm,yshift=-3.6cm]current page.north east) circle (4.6cm);
    \fill[poppurple!14] ([xshift=2.4cm,yshift=7.6cm]current page.south west) circle (3.8cm);
  \end{tikzpicture}%
  {\poplogo\fontsize{30}{30}\selectfont\color{popink}OnBuch\color{poporange}+}\hfill
  \raisebox{0.6ex}{\poppill{Cours signé \textbullet\ Premium}{popink}{popcream}}\par
  \vspace{1pt}\popsquiggle{poppurple}\par
  \vspace{8pt}
  \begin{tcolorbox}[enhanced,colback=poporange,colframe=popink,boxrule=2.8pt,arc=13pt,
      drop shadow=popink,left=18pt,right=18pt,top=12pt,bottom=13pt,after skip=10pt]
    \poppill{#2 \textbullet\ #3}{popink}{popcream}\hspace{5pt}\poppill{#4}{white}{popink}\par\vspace{8pt}
    {\popdisplay\fontsize{14}{14}\selectfont\color{popink}LEÇON #5}\par\vspace{4pt}
    {\raggedright\hyphenpenalty=10000\exhyphenpenalty=10000\popdisplay\fontsize{25}{27}\selectfont\color{popcream}\MakeUppercase{#1}\par}
    \vspace{8pt}
    {\popxbold\fontsize{9.5}{9.5}\selectfont\color{popink}\MakeUppercase{Durée conseillée : #6}}
  \end{tcolorbox}
  \begin{tcolorbox}[enhanced,colback=white,colframe=popink,boxrule=2.2pt,arc=10pt,
      drop shadow=popink,left=16pt,right=16pt,top=10pt,bottom=10pt,after skip=8pt]
    \poppill{Dans cette leçon}{poppurple}{white}\par\vspace{5pt}
    \popsemi\fontsize{9.3}{12.6}\selectfont\color{popink}#7
  \end{tcolorbox}
  \noindent\begin{minipage}[t]{0.487\linewidth}
    \begin{tcolorbox}[enhanced,colback=popgreen,colframe=popink,boxrule=2.2pt,arc=10pt,
        drop shadow=popink,left=11pt,right=11pt,top=10pt,bottom=10pt,height=3.1cm,valign=center]
      \tcbox[colback=white,colframe=popink,boxrule=1.3pt,arc=5pt,boxsep=0pt,nobeforeafter,
        left=3pt,right=3pt,top=3pt,bottom=3pt,valign=center]{\qrcode[height=1.9cm]{https://play.google.com/store/apps/details?id=com.luvvix.onbuch&hl=fr}}%
      \hspace{8pt}\begin{minipage}[c]{0.5\linewidth}
        {\popdisplay\fontsize{11}{12.5}\selectfont\color{white}\MakeUppercase{Télécharge OnBuch}}\par\vspace{4pt}
        {\raggedright\popsemi\fontsize{8.4}{11}\selectfont\color{white}Gratuit \textbullet\ Google Play\\Tuteur IA, annales, concours, résultats\par}
      \end{minipage}
    \end{tcolorbox}
  \end{minipage}\hfill
  \begin{minipage}[t]{0.487\linewidth}
    \begin{tcolorbox}[enhanced,colback=poppurple,colframe=popink,boxrule=2.2pt,arc=10pt,
        drop shadow=popink,left=11pt,right=11pt,top=10pt,bottom=10pt,height=3.1cm,valign=center]
      \tikz[baseline=-0.7ex]\node[circle,fill=white,draw=popink,line width=1.3pt,minimum size=1.6cm,inner sep=0]{\popdisplay\fontsize{24}{24}\selectfont\color{poppurple}+};%
      \hspace{8pt}\begin{minipage}[c]{0.6\linewidth}
        {\popdisplay\fontsize{11}{12.5}\selectfont\color{white}\MakeUppercase{Passe à OnBuch+}}\par\vspace{4pt}
        {\raggedright\popsemi\fontsize{8.4}{11}\selectfont\color{white}Tous les cours signés, Tuteur IA illimité, fiches PDF — onglet OnBuch+ de l'app\par}
      \end{minipage}
    \end{tcolorbox}
  \end{minipage}
  \vspace{8pt}
  \popcontenu
  \vfill
  \signatureonbuch
  \restoregeometry
  \newpage
}

% Rangée « Ce cours contient » (page de garde)
\newcommand{\poptile}[3]{% #1 couleur #2 gros texte #3 légende
  \begin{tcolorbox}[enhanced,colback=#1,colframe=popink,boxrule=2pt,arc=9pt,shadow={2.5pt}{-2.5pt}{0pt}{popink},
      left=6pt,right=6pt,top=9pt,bottom=9pt,halign=center,valign=center,height=1.75cm,width=\linewidth,nobeforeafter]
    {\popdisplay\fontsize{12.5}{13}\selectfont\color{white}\MakeUppercase{#2}}\par\vspace{3pt}
    {\centering\popsemi\fontsize{7.8}{9.5}\selectfont\color{white}#3\par}
  \end{tcolorbox}}
\newcommand{\popcontenu}{%
  \par\noindent{\popxboldls\fontsize{7.5}{7.5}\selectfont\color{popmuted}CE COURS SIGNÉ CONTIENT}\par\vspace{7pt}
  \noindent\begin{minipage}[t]{0.235\linewidth}\poptile{popblue}{Cours}{complet, illustré, pas à pas}\end{minipage}\hfill
  \begin{minipage}[t]{0.235\linewidth}\poptile{poporange}{Méthodes}{et exercices résolus}\end{minipage}\hfill
  \begin{minipage}[t]{0.235\linewidth}\poptile{poppink}{Exercices}{type examen + corrigés}\end{minipage}\hfill
  \begin{minipage}[t]{0.235\linewidth}\poptile{popgreen}{Fiche bilan}{l'essentiel à retenir}\end{minipage}\par}

% Sommaire
\newtcolorbox{sommaire}{enhanced,colback=white,colframe=popink,boxrule=2.2pt,arc=10pt,
  drop shadow=popink,left=18pt,right=18pt,top=13pt,bottom=13pt,before skip=6pt,after skip=20pt,
  before upper={\poppill{Sommaire}{poppurple}{white}\par\vspace{9pt}\popsemi\fontsize{10.6}{14.5}\selectfont\color{popink}}}

% Signature de fin de cours — poussée tout en bas de la dernière page
\newcommand{\findecours}{%
  \par\vfill\nopagebreak
  \begin{center}\popsquiggle{poporange}\end{center}\vspace{4pt}
  \signatureonbuch
}
\AtBeginDocument{\def\llbracket{\mathopen{[\mkern-3.5mu[}}\def\rrbracket{\mathclose{]\mkern-3.5mu]}}} % intervalles d'entiers (glyphe ⟦ absent de la police)
% Glyphes absents de Fira Math : redéfinis à partir de glyphes disponibles
\AtBeginDocument{%
  \def\setminus{\mathbin{\backslash}}\let\smallsetminus\setminus
  \def\vdots{\mathord{\vbox{\baselineskip3.2pt\lineskiplimit0pt\kern4pt\hbox{.}\hbox{.}\hbox{.}}}}%
  \def\ddots{\mathinner{\mkern1mu\raise6.5pt\hbox{.}\mkern2mu\raise3.6pt\hbox{.}\mkern2mu\raise0.7pt\hbox{.}\mkern1mu}}%
  \def\triangle{\mathord{\scalebox{0.9}{$\symup{\Delta}$}}}%
  \def\nabla{\mathord{\raisebox{\depth}{\scalebox{1}[-1]{$\symup{\Delta}$}}}}%
  \def\checkmark{\popcheck{popdark}}%
  \def\star{\mathbin{\ast}}\def\bigstar{\mathord{\ast}}%
  \def\diamond{\mathbin{\scalebox{0.6}{$\Diamond$}}}%
  \def\bigcup{\mathop{\mathchoice{\raisebox{-0.3em}{\scalebox{1.7}{$\cup$}}}{\scalebox{1.25}{$\cup$}}{\cup}{\cup}}\displaylimits}%
  \def\bigcap{\mathop{\mathchoice{\raisebox{-0.3em}{\scalebox{1.7}{$\cap$}}}{\scalebox{1.25}{$\cap$}}{\cap}{\cap}}\displaylimits}%
  \def\lesseqqgtr{\mathrel{\lessgtr}}\def\bigoplus{\mathop{\oplus}\displaylimits}%
}
\providecommand{\ssi}{si et seulement si} % abréviation fréquente
\AtBeginDocument{\providecommand{\wideparen}{\overparen}} % arc de cercle

\begin{document}

\pagegarde{Les transformations de l’élevage et de la pêche traditionnels}{Géographie}{Tle C}{Module 1 : Le Cameroun}{N05}{2 h + dossier 1 (2 h)}{Cette leçon étudie les mutations de l'élevage et de la pêche traditionnels au Cameroun : amélioration génétique des troupeaux, modernisation des équipements de pêche, essor de la pisciculture et nouvelles formes d'élevage (ranches, bergeries). Elle montre comment ces transformations, soutenues par l'État, contribuent à l'intégration nationale et à la valorisation des produits locaux.
\vspace{4pt}
\begin{multicols}{2}\small
\begin{itemize}
\item À la fin de cette leçon, je sais définir les notions de transhumance, bergerie, ranch et étang piscicole.
\item À la fin de cette leçon, je sais décrire les transformations de l'élevage traditionnel (sélection, croisements, alimentation, soins, infrastructures).
\item À la fin de cette leçon, je sais identifier et localiser sur une carte les grandes zones d'élevage et de pêche du Cameroun.
\item À la fin de cette leçon, je sais expliquer la modernisation des équipements de pêche artisanale et ses effets sur les productions.
\item À la fin de cette leçon, je sais analyser l'essor de la pisciculture et son rôle dans la sécurité alimentaire.
\item À la fin de cette leçon, je sais construire un croquis, légender une carte et construire un diagramme à partir de données statistiques sur l'élevage et la pêche.
\end{itemize}\end{multicols}}

\begin{sommaire}
\begin{multicols}{2}
\begin{enumerate}
\item Introduction : un élevage et une pêche traditionnels en mutation
\item Les transformations de l'élevage traditionnel
\item La modernisation de la pêche artisanale
\item L'essor de la pisciculture
\item Dossier 1 : le rôle de l'État dans le développement des activités agro-pastorales et piscicoles
\item TP guidé : cartographier et représenter les transformations
\item Activité d'intégration
\item Méthodes et exercices résolus
\item Exercices
\item Corrigés des exercices
\item Fiche bilan
\end{enumerate}
\end{multicols}
\end{sommaire}

\begin{objectifs}
\begin{itemize}
\item À la fin de cette leçon, je sais définir les notions de transhumance, bergerie, ranch et étang piscicole.
\item À la fin de cette leçon, je sais décrire les transformations de l'élevage traditionnel (sélection, croisements, alimentation, soins, infrastructures).
\item À la fin de cette leçon, je sais identifier et localiser sur une carte les grandes zones d'élevage et de pêche du Cameroun.
\item À la fin de cette leçon, je sais expliquer la modernisation des équipements de pêche artisanale et ses effets sur les productions.
\item À la fin de cette leçon, je sais analyser l'essor de la pisciculture et son rôle dans la sécurité alimentaire.
\item À la fin de cette leçon, je sais construire un croquis, légender une carte et construire un diagramme à partir de données statistiques sur l'élevage et la pêche.
\item À la fin de cette leçon, je sais interpréter des documents (cartes, photographies, tableaux statistiques) relatifs aux activités pastorales et piscicoles.
\item À la fin de cette leçon, je sais présenter le rôle de l'État dans le développement des activités agro-pastorales et piscicoles (programmes, financement, formation, régulation).
\item À la fin de cette leçon, je sais argumenter en faveur de la consommation des produits locaux dans une perspective d'intégration nationale.
\end{itemize}
\end{objectifs}

\begin{prerequis}
\begin{itemize}
\item Les pratiques agro-pastorales traditionnelles au Cameroun (leçons précédentes du chapitre 2) : élevage extensif, pêche artisanale, leurs limites.
\item La localisation des grandes régions agro-écologiques du Cameroun (Adamaoua, Nord, littoral, Sud forestier).
\item Les notions de production, productivité et de filière vues en classe de Première.
\item Les techniques de lecture et de légendage d'une carte géographique.
\item La notion d'intégration nationale et les enjeux de l'autosuffisance alimentaire.
\end{itemize}
\end{prerequis}

\newpage

\coursec{Introduction : un élevage et une pêche traditionnels en mutation}

\courssub{Situation d'accroche : deux portraits, un même défi}

À Maroua, dans l'Extrême-Nord, Ibrahim, éleveur peul, a fait un choix qui aurait surpris son grand-père : il a remplacé une partie de son troupeau de zébus Goudali par des croisés plus productifs, construit une bergerie moderne où ses animaux passent la nuit à l'abri, et vend désormais son lait à une laiterie de Ngaoundéré, à plus de 300 kilomètres de là. À Kribi, sur la côte atlantique, des pêcheurs artisanaux équipés de pirogues motorisées et de glacières remplies de glace approvisionnent chaque jour les marchés de Douala en bar, en doré, en thon et en crevettes. Pourtant, dans les rayons des supermarchés de Yaoundé et de Douala, le poisson congelé importé — maquereau, sardine, chinchard venus d'Europe ou d'Asie — concurrence directement les produits locaux, souvent à des prix inférieurs.

Ces deux scènes résument le paradoxe camerounais : un pays doté d'un immense potentiel pastoral et halieutique, mais qui dépend encore lourdement de l'extérieur pour nourrir sa population. Comment l'élevage et la pêche traditionnels camerounais se transforment-ils pour répondre à la demande nationale et renforcer l'intégration économique du pays ? Quel rôle l'État joue-t-il dans cette modernisation ? C'est à ces questions que cette leçon va répondre.

\begin{attention}[Problématique et annonce du plan]
La problématique de la leçon peut se formuler ainsi : \emph{comment les activités d'élevage et de pêche, longtemps traditionnelles, se modernisent-elles pour nourrir le pays et réduire la dépendance aux importations ?} Nous étudierons d'abord les transformations de l'élevage traditionnel (sélection, croisements, alimentation, soins, bâtiments), puis la modernisation de la pêche artisanale et l'essor de la pisciculture, enfin, dans le dossier 1, le rôle de l'État dans l'encadrement et le financement de ces mutations.
\end{attention}

\courssub{Rappel : les caractéristiques de l'élevage traditionnel camerounais}

Avant de comprendre ce qui change, il faut se souvenir de ce qui existait. Imagine un troupeau de cent zébus conduits par un berger à travers les savanes de l'Adamaoua : les bêtes paissent librement, sans clôture, sans complément alimentaire, sans suivi vétérinaire régulier. C'est l'image de l'\cle{élevage traditionnel}, dominant dans les régions du Nord, de l'Extrême-Nord et de l'Adamaoua.

\begin{definition}[Élevage extensif]
L'\cle{élevage extensif} est un mode d'élevage dans lequel les animaux paissent librement sur de vastes espaces naturels (parcours, savanes), avec peu d'investissements en bâtiments, en alimentation et en soins. Il se caractérise par une \textbf{faible productivité par tête} (peu de lait, croissance lente, reproduction non contrôlée), même si les troupeaux peuvent être nombreux.
\end{definition}

Les traits essentiels de cet élevage traditionnel sont les suivants :

\begin{itemize}
\item \textbf{Des troupeaux nombreux mais peu rentables.} Le nombre de têtes est un signe de richesse sociale et de prestige chez les éleveurs peuls (Foulbé) ou arabes choa, plus qu'un capital économique à rentabiliser. On vend peu, on consomme peu de lait transformé, et les pertes dues aux maladies (trypanosomiase, péripneumonie) et à la sécheresse sont élevées.
\item \textbf{La mobilité pastorale.} L'élevage repose sur la \cle{transhumance} : le déplacement saisonnier des troupeaux entre les pâturages de saison sèche et ceux de saison des pluies, parfois sur plusieurs centaines de kilomètres (de l'Extrême-Nord vers l'Adamaoua, voire vers le Nord-Ouest ou le Tchad voisin).
\item \textbf{Une faible productivité.} Une vache zébu de race locale produit souvent moins de 2 à 3 litres de lait par jour, contre 15 à 25 litres pour une race améliorée bien nourrie. Les carcasses sont légères et l'engraissement long.
\end{itemize}

\begin{exemplebox}[Un cheptel important mais mal valorisé]
Le cheptel camerounais compte environ \textbf{7 à 8 millions de bovins} (ordre de grandeur), auxquels s'ajoutent plusieurs millions de petits ruminants (ovins et caprins), de porcs et de volailles. Ce cheptel est concentré pour \textbf{plus des deux tiers dans les régions du Nord, de l'Extrême-Nord et de l'Adamaoua}, qui constituent le véritable « grenier à viande » du pays. Pourtant, la production nationale ne couvre pas entièrement la demande des villes du Sud, et le Cameroun importe encore de la viande et des produits laitiers (poudre de lait notamment).
\end{exemplebox}

\courssub{Rappel : les caractéristiques de la pêche traditionnelle}

Sur les berges de la Sanaga ou le long des côtes de Kribi et de Limbé, la scène traditionnelle est celle d'une pirogue creusée dans un tronc, propulsée à la pagaie ou à la voile, avec des filets à grandes mailles réparés à la main. C'est la \cle{pêche traditionnelle}, pratiquée sur la côte atlantique (de Campo à la baie d'Ambas), sur les fleuves (Sanaga, Wouri, Nyong, Logone), et sur les lacs (Tchad, Maga, Ossa, Barombi).

\begin{definition}[Pêche artisanale traditionnelle]
La \cle{pêche artisanale traditionnelle} est une pêche de petite échelle, pratiquée avec des embarcations non motorisées (pirogues à pagaie ou à voile), des engins rudimentaires (filets maillants, nasses, lignes, harpons) et des techniques de conservation sommaires (fumage, séchage au soleil). Elle vise l'autoconsommation et la vente sur les marchés locaux, avec des captures limitées et périssables.
\end{definition}

Ses caractéristiques principales :

\begin{itemize}
\item \textbf{Des équipements rudimentaires} : pirogues non motorisées limitant le rayon d'action à quelques kilomètres des côtes ou des berges ; filets peu adaptés qui laissent échapper une partie du poisson ou, à l'inverse, capturent des juvéniles.
\item \textbf{Une faible conservation des captures} : en l'absence de glace et de chaîne du froid, le poisson doit être vendu dans les heures qui suivent, fumé ou séché. Les pertes après capture sont considérables (on estime qu'une part importante des prises se dégrade avant la vente).
\item \textbf{Une production insuffisante} face à une demande urbaine en forte croissance, portée par l'augmentation de la population (plus de 28 millions d'habitants, ordre de grandeur).
\end{itemize}

\begin{attention}[Erreur fréquente : pêche artisanale $\neq$ pêche figée]
Ne confonds pas \emph{pêche artisanale} et \emph{pêche traditionnelle figée dans le passé}. La pêche artisanale désigne une pêche de \textbf{petite échelle} (petites unités, capitaux modestes, main-d'œuvre familiale), par opposition à la pêche industrielle (grands chalutiers). Or une pêche artisanale peut être \textbf{modernisée} tout en restant artisanale : pirogue motorisée, filets adaptés, glacière à bord. À l'inverse, « traditionnel » qualifie les techniques anciennes. Dans ta copie, emploie chaque terme avec précision : c'est toute la logique de cette leçon.
\end{attention}

\courssub{Le constat chiffré : un paradoxe camerounais}

Voici le fait qui dérange et qui justifie toute cette leçon : malgré ses 400 kilomètres de côtes, ses fleuves, ses lacs et ses millions de bovins, \textbf{le Cameroun importe chaque année environ 200\,000 à 300\,000 tonnes de poisson} (ordre de grandeur ; les volumes varient selon les années et les sources), essentiellement du poisson congelé bon marché destiné à la consommation populaire.

\begin{popfigure}
\begin{center}
\begin{tikzpicture}
\begin{axis}[
    ybar, bar width=9pt,
    width=0.92\linewidth, height=6.2cm,
    xlabel={Année}, ylabel={Importations (milliers de tonnes)},
    symbolic x coords={2010,2013,2016,2019,2021,2023},
    xtick=data, x tick label style={font=\small},
    ymin=0, ymax=350,
    ytick={0,50,100,150,200,250,300,350},
    nodes near coords, nodes near coords style={font=\scriptsize, color=popdark},
    every node near coord/.append style={/pgf/number format/fixed},
    axis lines=left,
    ymajorgrids=true, grid style={popline!40},
]
\addplot[fill=popblue, draw=popdark] coordinates {(2010,150) (2013,180) (2016,200) (2019,230) (2021,250) (2023,280)};
\end{axis}
\end{tikzpicture}
\end{center}
\legende{Évolution indicative des importations de poisson du Cameroun, 2010-2023 (en milliers de tonnes). Valeurs approchées, données à titre pédagogique : la tendance générale à la hausse traduit l'écart croissant entre la demande nationale et la production locale.}
\end{popfigure}

Le graphique montre une tendance claire : les importations augmentent avec la croissance démographique et urbaine, car la production nationale (pêche maritime artisanale, pêche continentale, aquaculture naissante) stagne ou progresse trop lentement. Le même constat vaut, dans une moindre mesure, pour la viande et surtout pour les produits laitiers.

\begin{savaistu}[Le coût de la dépendance]
Le Cameroun dépense chaque année \textbf{plus de 100 milliards de FCFA} (ordre de grandeur) pour importer du poisson, alors que son potentiel piscicole — plans d'eau naturels, zones favorables à la pisciculture dans l'Ouest, le Centre et le Sud — est largement sous-exploité. Cet argent qui sort du pays pourrait financer des milliers de fermes piscicoles et d'emplois ruraux. C'est l'un des enjeux majeurs de la souveraineté alimentaire nationale.
\end{savaistu}

\courssub{Localiser : les grandes zones d'élevage et de pêche}

Pour raisonner juste, il faut d'abord situer. L'élevage est avant tout une affaire du \textbf{Nord} (savanes sèches, peu de glossines, longue tradition pastorale), tandis que la pêche se partage entre la \textbf{côte atlantique}, les \textbf{grands fleuves} et les \textbf{lacs} du Nord et du Littoral.

\begin{popfigure}
\begin{center}
\includegraphics[width=\linewidth,height=11cm,keepaspectratio]{N05/cmr_regions.jpg}
\end{center}
\legende{Fond de carte : le Cameroun en grands ensembles (carte en anglais : Northern Cameroon = Grand Nord, Adamaoua, Coastal Cameroon = littoral), avec villes, routes, voies ferrées, fleuves (Logone, Chari, Bénoué, Sanaga, Nyong) et plans d'eau (lac Tchad, retenues de Lagdo et de Mbakaou). Les élèves y repèrent les grandes zones d'élevage (Adamaoua, Nord et Extrême-Nord) et les zones de pêche (côte atlantique de Limbé à Kribi, lac Tchad, Logone, Sanaga). \textbf{À compléter} sur ton propre fond de carte : le lac de Maga, le lac Ossa, le Wouri et les ports de pêche. \\ Source : Wikimedia Commons, Peter Fitzgerald (modifié par Burmesedays), CC BY 3.0.}
\end{popfigure}

\begin{methode}[Lire et compléter un croquis de localisation]
\begin{enumerate}
\item \textbf{Observer} la répartition : où se situent les grandes savanes d'élevage (Adamaoua, Nord, Extrême-Nord) par rapport à la côte et aux plans d'eau de pêche ? Quel contraste Nord/Sud observes-tu ?
\item \textbf{Expliquer} la localisation : savanes sèches et tradition pastorale au Nord ; côte, fleuves et lacs pour la pêche.
\item \textbf{Compléter} le croquis en ajoutant les éléments manquants demandés (villes, plans d'eau, axes d'écoulement des produits vers les villes du Sud).
\item \textbf{Légender} proprement : chaque symbole utilisé doit figurer dans la légende, avec un titre général au croquis.
\end{enumerate}
\end{methode}

\courssub{De la tradition à la modernisation : la logique de la leçon}

Face à ce paradoxe — un potentiel abondant mais une dépendance coûteuse —, l'élevage et la pêche camerounais connaissent depuis quelques décennies de profondes \cle{transformations}. L'intuition est simple : pour produire plus et mieux, il faut agir sur trois leviers — les \textbf{animaux et les techniques} (races, alimentation, soins, engins de pêche), les \textbf{équipements et infrastructures} (bergeries, poulaillers, moteurs, glacières, étangs), et l'\textbf{organisation} (encadrement, financement, commercialisation). Chacun de ces leviers fera l'objet d'une étude détaillée.

\begin{exoresolu}[Analyser un constat chiffré]
\textbf{Énoncé.} Le Cameroun importe environ 250\,000 tonnes de poisson par an pour une dépense d'environ 100 milliards de FCFA. (1) Calcule le prix moyen d'un kilogramme de poisson importé. (2) Explique en cinq lignes pourquoi ce constat justifie la modernisation de la pêche et de l'élevage nationaux.
\tcblower
\textbf{Solution détaillée.}
(1) Conversion : 250\,000 tonnes $= 250\,000\,000$ kg, soit $2,5 \times 10^{8}$ kg. Prix moyen :
\[ \frac{100 \times 10^{9}\ \text{FCFA}}{2,5 \times 10^{8}\ \text{kg}} = 400\ \text{FCFA/kg}. \]
Le poisson importé coûte donc en moyenne environ 400 FCFA le kilogramme à l'importation (hors marges de distribution).
(2) Ce constat révèle un écart entre une demande nationale forte et une production locale insuffisante : le pays sort chaque année d'énormes devises pour acheter à l'étranger ce qu'il pourrait produire. Moderniser la pêche (motorisation, conservation) et développer la pisciculture et l'élevage permettraient d'augmenter l'offre locale, de créer des emplois ruraux, de réduire la facture d'importation et de renforcer l'intégration économique nationale en reliant les zones de production aux marchés de consommation.
\end{exoresolu}

\begin{aretenir}[L'essentiel de l'introduction]
\begin{itemize}
\item L'élevage traditionnel camerounais est \textbf{extensif} : troupeaux nombreux, mobilité pastorale (transhumance), faible productivité par tête ; il est concentré au Nord, à l'Extrême-Nord et à l'Adamaoua (7 à 8 millions de bovins au total, ordre de grandeur).
\item La pêche traditionnelle est \textbf{artisanale et rudimentaire} : pirogues non motorisées, filets peu adaptés, conservation précaire ; elle se pratique sur la côte atlantique, les fleuves et les lacs.
\item Paradoxe majeur : le Cameroun importe \textbf{200\,000 à 300\,000 tonnes de poisson par an} pour plus de 100 milliards de FCFA, malgré son potentiel.
\item La réponse passe par la \textbf{modernisation} de l'élevage, de la pêche et par l'essor de la pisciculture, soutenue par l'action de l'État (dossier 1).
\end{itemize}
\end{aretenir}

\begin{exercice}[Pour vérifier ta compréhension]
\begin{enumerate}
\item Définis l'élevage extensif et donne deux de ses limites pour l'économie nationale.
\item Sur un croquis vierge du Cameroun, place et légende : deux zones d'élevage, trois zones de pêche, trois villes de consommation ou de transformation.
\item À l'aide du graphique des importations, décris en quelques lignes la tendance observée et propose une explication.
\end{enumerate}
\end{exercice}

\coursec{Les transformations de l'élevage traditionnel}

Dans l'introduction, nous avons vu que l'élevage et la pêche camerounais, longtemps dominés par des pratiques héritées, connaissent depuis les années 1970 une \cle{mutation} profonde : l'objectif n'est plus seulement de nourrir la famille ou le village, mais de produire davantage, de mieux vendre et d'intégrer les campagnes à l'économie nationale. Cette transformation touche d'abord l'élevage, pilier de l'économie des savanes du Nord et de l'Adamaoua. Comment passe-t-on d'un troupeau de zébus maigres conduit en transhumance à un ranch organisé, à une bergerie moderne ou à un poulailler périurbain ? C'est toute la question de cette section, que nous aborderons en suivant la démarche de l'éleveur lui-même : choisir de bons animaux, les croiser, mieux les nourrir, mieux les soigner, mieux les loger, puis changer d'échelle et d'organisation.

\courssub{La sélection des espèces et des races}

Tout commence par l'animal. L'éleveur traditionnel peul de l'Adamaoua possède des zébus parfaitement adaptés à leur milieu, mais peu productifs. La première transformation consiste donc à \cle{sélectionner} les espèces et les races.

\begin{definition}[Sélection des races]
La \cle{sélection} est le choix délibéré, par l'éleveur ou le technicien, des animaux reproducteurs les plus performants (croissance rapide, bonne production laitière, résistance aux maladies) afin d'améliorer progressivement les qualités du troupeau.
\end{definition}

Deux grandes stratégies coexistent au Cameroun :

\begin{itemize}
\item \textbf{La valorisation des races locales rustiques.} Le zébu \cle{Goudali}, race emblématique de l'Adamaoua et du Nord, est réputé pour sa bonne conformation bouchère et sa résistance à la sécheresse ; le zébu \cle{Bororo} (ou Peul), plus grand, est apprécié pour sa viande et sa relative production laitière. Ces races supportent les pâturages pauvres, la chaleur et les longues marches : c'est leur \emph{rusticité} qui fait leur force.
\item \textbf{L'introduction de races améliorées.} Des races exotiques, plus productives en viande et surtout en lait, ont été introduites : la \cle{Holstein} (grande laitière d'origine européenne, pouvant produire 20 à 30 litres de lait par jour dans de bonnes conditions, contre 2 à 5 litres pour une zébu locale), la \cle{Montbéliarde}, race mixte (viande et lait) originaire de France, plus tolérante aux conditions difficiles.
\end{itemize}

\begin{attention}[Le piège de la race « miracle »]
Une race exotique très productive n'est performante que si elle est bien nourrie, bien soignée et abritée. Une Holstein conduite en transhumance sur des parcours pauvres produira peu et tombera malade. La sélection ne fonctionne donc jamais seule : elle exige l'ensemble du paquet technique (alimentation, soins, logement).
\end{attention}

\courssub{Les croisements : combiner rusticité et productivité}

Comment concilier la résistance des races locales et la productivité des races exotiques ? La réponse technique est le \cle{croisement}.

\begin{definition}[Croisement]
Le \cle{croisement} est l'accouplement d'animaux de races différentes afin d'obtenir des descendants (métis) combinant les qualités des deux parents : rusticité de la race locale, productivité de la race exotique.
\end{definition}

Au Cameroun, la politique d'amélioration génétique s'est appuyée sur des stations publiques, notamment le \textbf{ranch de Wakwa}, près de Ngaoundéré (Adamaoua), créé dans les années 1960-1970 avec l'appui de la SODEPA (Société de développement et d'exploitation des productions animales). On y a croisé des zébus locaux (Goudali, Bororo) avec des taureaux ou des semences de races exotiques (Holstein, Montbéliarde) pour diffuser des métis à meilleure production laitière. L'insémination artificielle permet de multiplier ces croisements sans transporter les animaux.

\begin{exemplebox}[Le métis, un compromis raisonnable]
Une vache métisse Goudali $\times$ Montbéliarde produit typiquement 6 à 10 litres de lait par jour (ordre de grandeur) : moins qu'une Holstein pure, mais deux à trois fois plus qu'une zébu locale, tout en restant capable de pâturer les savanes soudaniennes et de résister mieux aux parasites. C'est ce compromis qui est diffusé chez les éleveurs sédentarisés de l'Adamaoua.
\end{exemplebox}

\courssub{L'amélioration de l'alimentation : nourrir pour produire}

Un animal sélectionné mais mal nourri reste improductif. Dans l'élevage traditionnel, le troupeau dépend entièrement des pâturages naturels : abondants en saison des pluies, rares et secs de novembre à avril, période où les bêtes maigrissent et où le lait tarit. La modernisation passe par la \textbf{sécurisation de l'alimentation} :

\begin{itemize}
\item les \textbf{cultures fourragères} : parcelles d'herbes améliorées (Brachiaria, Panicum) et de légumineuses fourragères ( Stylosanthes, luzerne), foin et paille conservés pour la saison sèche ;
\item les \textbf{compléments alimentaires} : tourteaux de coton (sous-produit de la SODECOTON dans le Nord), sons de maïs, pierres à lécher et compléments minéralo-vitaminés ;
\item les \textbf{parcours aménagés} : pâturages clôturés et divisés en parcelles exploitées en rotation, pour laisser l'herbe se régénérer ;
\item les \textbf{points d'eau} : forages pastoraux, retenues collinaires et mares aménagées, qui évitent aux troupeaux des dizaines de kilomètres de marche vers les rares cours d'eau permanents.
\end{itemize}

\begin{aretenir}[La logique de la transformation alimentaire]
Passer d'une alimentation \emph{subie} (le troupeau cherche seul sa nourriture selon les saisons) à une alimentation \emph{maîtrisée} (fourrages cultivés et stockés, compléments, eau disponible toute l'année) est la condition de toute intensification : c'est ce qui permet de fixer les animaux, de les engraisser plus vite et de produire du lait toute l'année.
\end{aretenir}

\courssub{L'amélioration des soins : la santé animale}

Les maladies animales (\cle{épizooties}) ont longtemps décimé les troupeaux : peste bovine au début du XX\textsuperscript{e} siècle, trypanosomiase (maladie transmise par la mouche tsé-tsé), péripneumonie contagieuse, parasitoses. La modernisation sanitaire repose sur :

\begin{itemize}
\item la \textbf{vaccination} systématique (campagnes annuelles contre la péripneumonie, l'anthrax, la pasteurellose) ;
\item le \textbf{déparasitage} régulier (bains parasiticides, vermifuges contre les tiques et les vers) ;
\item l'action des \textbf{vétérinaires de campagne} et des techniciens du MINEPIA (ministère de l'Élevage, des Pêches et des Industries animales), qui encadrent les éleveurs jusque dans les cantons ;
\item les \textbf{pharmacies vétérinaires rurales}, qui rapprochent les médicaments des zones d'élevage.
\end{itemize}

\begin{savaistu}[La trypanosomiase, frein historique à l'élevage]
La mouche tsé-tsé, présente dans les zones humides et forestières, transmet la trypanosomiase (ou « nagana »), mortelle pour les bovins. Elle a longtemps interdit l'élevage dans une grande partie du Centre et du Sud du Cameroun. La lutte contre la tsé-tsé (pièges imprégnés, assainissement des berges) et l'élevage de races tolérantes (petits ruminants, bovins N'Dama à l'Ouest) ont permis d'étendre les zones d'élevage vers le sud.
\end{savaistu}

\courssub{La construction d'infrastructures : bergeries, poulaillers, abattoirs}

L'animal moderne doit être logé, et la production doit pouvoir être vendue. D'où un effort de construction d'infrastructures :

\begin{definition}[Bergerie]
Une \cle{bergerie} est un bâtiment ou un enclos aménagé pour abriter et protéger les petits ruminants (ovins, caprins) et parfois les bovins : elle met les animaux à l'abri des pluies, des prédateurs et des vols, et facilite l'alimentation, les soins et la reproduction contrôlée.
\end{definition}

\begin{itemize}
\item \textbf{Poulaillers modernes} : bâtiments ventilés, sur caillebotis ou litière, avec mangeoires et abreuvoirs, pouvant abriter de quelques centaines à plusieurs milliers de poulets de chair ou de pondeuses ;
\item \textbf{porcheries} : logement du porc, animal à croissance rapide très prisé dans l'Ouest et le Centre ;
\item \textbf{abattoirs et marchés à bétail} : abattoirs modernes (Douala, Yaoundé, Ngaoundéré, Maroua), marchés à bétail hebdomadaires (Ngaoundéré, Figuil, Meiganga) qui structurent la commercialisation et permettent le contrôle sanitaire de la viande.
\end{itemize}

\courssub{L'évolution des formes d'élevage : de la transhumance au ranch}

\begin{definition}[Transhumance]
La \cle{transhumance} est le déplacement saisonnier des troupeaux et de leurs gardiens entre les zones de pâturage de saison sèche (vallées, plaines inondées comme celles du Logone ou du Mbéré, où l'herbe repousse après la décrue) et les zones de pâturage de saison des pluies (hautes savanes, à l'abri des mouches et des inondations).
\end{definition}

La transhumance n'est pas un errance : elle suit des \textbf{itinéraires codifiés} (couloirs de passage) et un \textbf{calendrier précis} (descente vers les plaines en novembre-décembre, remontée en mai-juin). Mais elle est aujourd'hui \textbf{encadrée et en recul} : textes réglementaires sur les couloirs de transhumance, conflits avec les agriculteurs dont les champs sont dégradés, pression foncière croissante.

\begin{attention}[Erreur fréquente au Bac]
Ne présente jamais la transhumance comme un « nomadisme désorganisé ». C'est un système pastoral \emph{rationnel}, adapté aux saisons, avec des routes et un calendrier. Dire qu'elle est « archaïque » sans nuance est une erreur : elle reste efficace pour valoriser des espaces non cultivables, mais elle est de plus en plus encadrée par l'État et concurrencée par la sédentarisation.
\end{attention}

La tendance lourde est la \textbf{sédentarisation} :

\begin{definition}[Ranch]
Un \cle{ranch} est une grande exploitation d'élevage clôturée et organisée, où les animaux sont conduits de manière rationnelle : pâturages divisés en parcelles exploitées en rotation, points d'eau permanents, soins vétérinaires réguliers, sélection du cheptel.
\end{definition}

On distingue désormais trois systèmes :

\begin{itemize}
\item l'\textbf{élevage extensif} (transhumant ou sédentaire), dominant dans l'Adamaoua, le Nord et l'Extrême-Nord ;
\item l'\textbf{élevage semi-intensif} : animaux parqués la nuit, compléments alimentaires, soins réguliers (bergeries améliorées, petits élevages laitiers de l'Ouest) ;
\item l'\textbf{élevage intensif} : animaux en bâtiments, alimentation industrielle complète — c'est le cas de l'\textbf{aviculture périurbaine} autour de Yaoundé et de Douala, qui fournit l'essentiel des poulets de chair et des œufs du pays.
\end{itemize}

\begin{exemplebox}[L'aviculture périurbaine, vitrine de l'élevage moderne]
La production avicole camerounaise fournit chaque année des \textbf{dizaines de millions de poulets de chair} et plusieurs centaines de millions d'œufs (ordres de grandeur), produits pour l'essentiel dans les périphéries des grandes villes (Yaoundé, Douala, Bafoussam, Bamenda). Cette aviculture intensive emploie des milliers de personnes mais reste fragile : coût élevé des aliments (maïs, soja) et des intrants vétérinaires, concurrence des importations de poulet congelé.
\end{exemplebox}

\begin{popfigure}
\begin{center}
\begin{tikzpicture}[font=\small]
% Colonne traditionnel
\fill[poporangeL,rounded corners=4pt] (0,0) rectangle (5.4,6.2);
\node[font=\small\bfseries,text=popdark] at (2.7,5.8) {Élevage traditionnel (extensif)};
\node[align=left,text width=4.8cm] at (2.7,3.1) {
\textbf{Alimentation :} pâturages naturels, saisonniers\\[2pt]
\textbf{Soins :} pharmacopée locale, peu de vaccins\\[2pt]
\textbf{Logement :} enclos de nuit en branchages\\[2pt]
\textbf{Déplacement :} transhumance saisonnière\\[2pt]
\textbf{Productivité :} faible (2--5 L lait/j, croissance lente)
};
% Colonne moderne
\fill[popgreenL,rounded corners=4pt] (6.2,0) rectangle (11.6,6.2);
\node[font=\small\bfseries,text=popdark] at (8.9,5.8) {Élevage modernisé (semi-intensif / intensif)};
\node[align=left,text width=4.8cm] at (8.9,3.1) {
\textbf{Alimentation :} fourrages cultivés, foin, compléments\\[2pt]
\textbf{Soins :} vaccination, déparasitage, vétérinaire\\[2pt]
\textbf{Logement :} bergerie, poulailler, ranch clôturé\\[2pt]
\textbf{Déplacement :} sédentarisation, parcours en rotation\\[2pt]
\textbf{Productivité :} élevée (6--20 L lait/j, croissance rapide)
};
% Flèche centrale
\draw[-{Stealth[length=4mm]},line width=2.5pt,popblue] (5.55,3.1) -- (6.05,3.1);
\node[font=\scriptsize\itshape,text=popblue] at (5.8,3.45) {mutation};
\end{tikzpicture}
\end{center}
\legende{Schéma comparatif de l'élevage traditionnel et de l'élevage modernisé au Cameroun. La modernisation porte simultanément sur cinq piliers : alimentation, soins, logement, mobilité et productivité.}
\end{popfigure}

\courssub{Étude de document : photographie d'une bergerie moderne dans un ranch de l'Adamaoua}

\textbf{Description de la photographie (document d'étude).} Au premier plan, un bâtiment long, à toit de tôle ondulée, aux murs bas en parpaings surmontés d'une claustra : c'est la bergerie, où des ovins et des caprins sont rentrés pour la nuit. Devant le bâtiment, des auges (mangeoires) en béton et un abreuvoir alimenté par un forage équipé d'une éolienne. À l'arrière-plan, des pâturages divisés par des clôtures de fil de fer, quelques zébus Goudali au pacage, et un hangar à foin. Au loin, la savane arborée de l'Adamaoua.

\begin{methode}[Analyser une photographie d'élevage]
\begin{enumerate}
\item \textbf{Décrire} : inventorier les éléments visibles (animaux, bâtiments, équipements, paysage) sans les interpréter.
\item \textbf{Localiser} : situer le lieu (région, milieu naturel) à partir des indices paysagers (ici, savane soudanienne de l'Adamaoua).
\item \textbf{Interpréter} : relever les indices de modernisation (bâtiment durable, clôtures, forage, stockage du fourrage) et les relier aux transformations étudiées.
\item \textbf{Conclure} : identifier le type d'élevage (ici, élevage semi-intensif sédentarisé de type ranch) et évaluer ses avantages et ses limites.
\end{enumerate}
\end{methode}

\begin{exoresolu}[Interpréter la photographie de la bergerie de l'Adamaoua]
Énoncé : à partir de la photographie décrite ci-dessus, montrez qu'il s'agit d'un élevage en voie de modernisation et précisez le système d'élevage représenté.
\tcblower
Solution détaillée : la photographie montre d'abord des éléments traditionnels : des zébus Goudali (race locale rustique) au pâturage dans une savane arborée, ce qui situe la scène dans l'Adamaoua, grande région d'élevage. Mais plusieurs indices révèlent la modernisation : (1) la \textbf{bergerie} en matériaux durables (parpaings, tôle) remplace l'enclos de branchages : les animaux sont protégés et les soins facilités ; (2) les \textbf{clôtures} divisent les pâturages en parcelles, signe d'une conduite rationnelle en rotation ; (3) le \textbf{forage à éolienne} garantit l'abreuvement permanent, libérant le troupeau de la transhumance ; (4) le \textbf{hangar à foin} témoigne du stockage de fourrage pour la saison sèche. La combinaison de ces éléments (clôtures, eau, soins, logement) caractérise un \textbf{ranch}, c'est-à-dire un élevage semi-intensif sédentarisé, à l'image des ranches de l'Adamaoua (Wakwa). Conclusion : cette photographie illustre la transition de l'élevage extensif transhumant vers un élevage organisé et plus productif, sans rupture totale avec les races locales.
\end{exoresolu}

\courssub{Bilan : des productions en hausse, des contraintes persistantes}

La modernisation porte ses fruits : la production de \textbf{viande}, de \textbf{lait} et d'\textbf{œufs} a fortement augmenté depuis les années 1980, l'autosuffisance en œufs est atteinte, et le Cameroun exporte même des bovins sur pied vers le Nigeria et le Gabon. Le cheptel bovin national est estimé à plusieurs millions de têtes (environ 6 à 7 millions, ordre de grandeur), très inégalement réparti.

\begin{popfigure}
\begin{center}
\begin{tikzpicture}
\begin{axis}[
  width=0.85\linewidth, height=6.5cm,
  ybar, bar width=18pt,
  ymin=0, ymax=40,
  ylabel={Part du cheptel bovin (\%)},
  symbolic x coords={Adamaoua, Nord, Extrême-Nord, Ouest, Autres},
  xtick=data,
  x tick label style={font=\scriptsize},
  y tick label style={font=\scriptsize},
  nodes near coords, nodes near coords style={font=\scriptsize},
  axis lines=left,
  enlarge x limits=0.15
]
\addplot[fill=popblue,draw=popdark] coordinates {(Adamaoua,32) (Nord,22) (Extrême-Nord,20) (Ouest,10) (Autres,16)};
\end{axis}
\end{tikzpicture}
\end{center}
\legende{Répartition approximative du cheptel bovin camerounais par grande région (parts en \%, ordres de grandeur). Les trois régions septentrionales (Adamaoua, Nord, Extrême-Nord) concentrent environ les trois quarts du cheptel national.}
\end{popfigure}

Mais les contraintes demeurent lourdes :

\begin{itemize}
\item le \textbf{coût des intrants} (aliments concentrés, médicaments vétérinaires, matériel), souvent importés, pèse sur la rentabilité, surtout en aviculture ;
\item les \textbf{conflits agriculteurs--éleveurs}, liés aux dégâts des troupeaux dans les champs et à la pression foncière (Nord-Ouest, Adamaoua, vallée de la Bénoué) ;
\item l'\textbf{enclavement} des zones d'élevage : pistes dégradées, coût du transport du bétail vers les métropoles du Sud ;
\item la \textbf{concurrence des importations} (poulet congelé, lait en poudre) qui freine les productions locales.
\end{itemize}

\begin{aretenir}[L'essentiel de la section]
La transformation de l'élevage camerounais repose sur cinq leviers complémentaires : \cle{sélection} des races (Goudali, Bororo, races exotiques), \cle{croisements} (ranches de l'Adamaoua, Wakwa), amélioration de l'\cle{alimentation} (fourrages, forages pastoraux) et des \cle{soins} (vaccination, vétérinaires de campagne), et construction d'\cle{infrastructures} (bergeries, poulaillers, abattoirs). Elle s'accompagne d'un passage progressif de la \cle{transhumance} extensive à la sédentarisation (\cle{ranch}, élevage semi-intensif et intensif, aviculture périurbaine). Les productions augmentent, mais le coût des intrants, les conflits d'usage de l'espace et l'enclavement limitent encore la mutation.
\end{aretenir}

\begin{exercice}[Pour t'entraîner]
Sur un croquis schématique du Cameroun, localise : (a) les trois grandes régions d'élevage bovin extensif ; (b) le ranch de Wakwa ; (c) les deux métropoles autour desquelles se concentre l'aviculture intensive ; (d) un couloir de transhumance entre l'Adamaoua et la vallée de la Bénoué. Légende et orientation obligatoires. Puis réponds en cinq lignes : pourquoi la sédentarisation des éleveurs est-elle à la fois une solution et une source de tensions ?
\end{exercice}

\coursec{La modernisation de la pêche artisanale}

Nous venons de voir comment l'élevage traditionnel se transforme : sélection des races, croisements, bergeries modernes, transhumance organisée. Le même mouvement de modernisation touche un autre pilier des pratiques agro-pastorales camerounaises : la \cle{pêche}. Sur la plage de Youpwè à Douala ou sur les berges du Logone à Kousséri, observe bien : les pirogues ont désormais des moteurs, les glacières remplacent les calebasses, et le poisson part chaque matin vers les marchés des grandes villes. La pêche artisanale, héritée des traditions, est en pleine mutation. C'est cette transformation que nous allons étudier : où pêche-t-on au Cameroun ? Qu'est-ce qui a changé dans les équipements ? Quels sont les effets et les limites de cette modernisation ?

\begin{definition}[La pêche artisanale]
La \cle{pêche artisanale} est une pêche pratiquée à \textbf{petite échelle}, avec des \textbf{embarcations légères} (pirogues, barques) et des \textbf{engins simples} (filets, lignes, nasses), destinée à l'\textbf{autoconsommation} et au \textbf{marché local}. Elle s'oppose à la pêche industrielle, pratiquée par de grands chalutiers équipés de technologies lourdes (radars, sonars, congélation à bord) et visant l'exportation ou les grands marchés urbains.
\end{definition}

\courssub{Les zones de pêche du Cameroun}

Le Cameroun dispose de deux grands ensembles de pêche, qu'il faut savoir distinguer et localiser.

\textbf{La pêche maritime} se pratique sur la côte atlantique, longue d'environ \SI{400}{km}, entre le Rio del Rey (frontière nigériane) et Campo (frontière équato-guinéenne). Les principaux sites de pêche artisanale maritime sont :
\begin{itemize}
\item \textbf{Limbé} et la baie d'Ambas, au pied du mont Cameroun ;
\item \textbf{Kribi}, grand centre de pêche du Sud, où la pêche artisanale côtoie le port en eau profonde ;
\item \textbf{Douala} et l'estuaire du Wouri (Youpwè, Bonabéri, Manoka), zone la plus poissonneuse du pays ;
\item \textbf{Campo}, à l'extrême sud, près de l'embouchure du Ntem.
\end{itemize}

\textbf{La pêche continentale} (ou pêche en eaux intérieures) se pratique sur les lacs, fleuves et retenues de barrages :
\begin{itemize}
\item le \textbf{lac Tchad} et la plaine d'inondation du \textbf{Logone} (Waza, Kousséri), dans l'Extrême-Nord ;
\item le \textbf{lac de barrage de Maga} et le lac de barrage de \textbf{Lagdo} sur la Bénoué ;
\item les grands fleuves : la \textbf{Sanaga}, le \textbf{Wouri}, le Nyong, le Ntem, la Bénoué.
\end{itemize}

\begin{exemplebox}[Une activité vitale en chiffres]
La pêche artisanale maritime fournit l'\textbf{essentiel du poisson consommé localement} au Cameroun et fait vivre, directement ou indirectement, des \textbf{dizaines de milliers de personnes} : pêcheurs, mareyeuses, fumeurs, transporteurs, réparateurs de moteurs. La production nationale totale de pêche (captures + aquaculture) est de l'ordre de \num{200000} à \num{300000} tonnes par an (ordre de grandeur), mais elle reste \textbf{inférieure à la demande} : le Cameroun importe chaque année plusieurs centaines de milliers de tonnes de poisson congelé (maquereaux, sardinelles) pour nourrir sa population.
\end{exemplebox}

\begin{popfigure}
\begin{center}
\includegraphics[width=\linewidth,height=9.5cm,keepaspectratio]{N01/cmr_map.jpg}
\end{center}
\legende{Fond de carte : carte générale du Cameroun sur laquelle se repèrent les principales zones de pêche. Pêche maritime artisanale : la côte du golfe de Guinée, avec Douala, Kribi et Campo (Buea et le mont Cameroun à proximité). Pêche continentale : le lac Tchad, le Logone et le Chari autour de Kousséri, la Bénoué et la retenue de Lagdo, la Sanaga, ainsi que la retenue de Mbakaou. \\ Source : Wikimedia Commons, CIA, domaine public.}
\end{popfigure}

\courssub{La motorisation des pirogues : aller plus loin, plus longtemps}

\textbf{L'intuition d'abord.} Imagine un pêcheur de Kribi il y a cinquante ans : il pagaie pendant des heures, ne s'éloigne guère de la côte, et rentre dès que le vent se lève. Aujourd'hui, le même pêcheur installe un \textbf{moteur hors-bord} à l'arrière de sa pirogue : en une heure, il atteint des zones de pêche autrefois inaccessibles.

\begin{definition}[La motorisation des pirogues]
La \cle{motorisation} consiste à équiper les pirogues traditionnelles de \textbf{moteurs hors-bord} (souvent de 15 à 40 chevaux). Les \textbf{grandes pirogues motorisées}, longues parfois de plus de dix mètres, permettent d'aller \textbf{plus loin en mer}, de rester \textbf{plus longtemps} sur les zones de pêche, de transporter \textbf{plus de poisson} et de rentrer plus vite, ce qui limite la détérioration des captures.
\end{definition}

Les effets sont considérables : élargissement de l'espace de pêche, augmentation du temps de pêche effectif, hausse des captures. Mais la motorisation a un coût : achat et entretien du moteur, carburant, pièces détachées importées. Le pêcheur devient dépendant du prix de l'essence et des devises.

\courssub{L'adaptation des engins de pêche}

La modernisation concerne aussi les \textbf{engins} :
\begin{itemize}
\item les \textbf{filets maillants} adaptés à chaque espèce (le poisson se prend par les ouïes dans les mailles) ;
\item les \textbf{sennes}, grands filets qui encerclent un banc de poissons ;
\item les \textbf{palangres}, longues lignes munies de nombreux hameçons.
\end{itemize}

Ces engins, fabriqués en fibres synthétiques (nylon), sont plus résistants et plus efficaces que les fibres végétales d'autrefois. Mais leur efficacité pose un problème majeur : des \textbf{mailles trop petites} capturent les juvéniles, qui n'ont pas eu le temps de se reproduire. C'est pourquoi l'État fixe des \textbf{tailles de mailles réglementaires} : c'est le cœur du débat sur la \cle{pêche durable}, c'est-à-dire une pêche qui exploite la ressource sans épuiser les stocks pour les générations futures.

\begin{attention}[Erreur fréquente : « le moteur règle tout »]
Ne crois pas que la motorisation résout à elle seule tous les problèmes de la pêche artisanale. \textbf{Sans conservation ni organisation commerciale, une partie importante des captures est perdue après capture} : le poisson se gâte en quelques heures sous les tropiques. Un pêcheur qui capture beaucoup mais ne peut ni conserver ni vendre rapidement perd son revenu. La modernisation est une \textbf{chaîne} : capture, conservation, transport, commercialisation. Un maillon faible fragilise tout le reste.
\end{attention}

\courssub{L'amélioration de la conservation}

Le poisson est un produit \textbf{très périssable}. Toute la modernisation de la pêche artisanale repose donc sur la conservation :
\begin{itemize}
\item les \textbf{glacières} et la \textbf{glace produite localement} (usines à glace à Douala, Kribi, Limbé) permettent de conserver le poisson frais pendant le transport vers les villes ;
\item les \textbf{chambres froides} installées sur les sites de débarquement stockent les excédents ;
\item le \textbf{fumage amélioré} : les \textbf{foyers améliorés} (fours en briques ou en tôle, type Chorkor) remplacent le fumage traditionnel sur bois de chauffe ; ils consomment moins de bois, fument plus uniformément et donnent un produit plus sain et mieux présenté ;
\item le \textbf{séchage} au soleil, pratiqué surtout dans le Nord (poisson du lac Tchad et du Logone séché puis exporté vers tout le pays et le Nigeria).
\end{itemize}

\courssub{L'organisation des pêcheurs et le rôle des femmes}

La modernisation n'est pas seulement technique : elle est aussi \textbf{sociale et commerciale}. Les pêcheurs se regroupent en \textbf{coopératives} et en \textbf{GIC} (Groupes d'Initiative Commune) pour acheter ensemble moteurs et filets, obtenir des crédits et défendre leurs intérêts. Des \textbf{unités de transformation et de commercialisation} s'organisent autour des débarcadères.

\begin{savaistu}[Les mareyeuses, un maillon central]
Au Cameroun comme dans toute l'Afrique de l'Ouest, les \textbf{femmes jouent un rôle central} dans la pêche artisanale : ce sont elles, les \cle{mareyeuses} et les \textbf{fumeuses de poisson}, qui achètent le poisson au débarcadère, le transforment (fumage, séchage) et le commercialisent sur les marchés urbains et ruraux. Sans elles, le poisson pêché le matin n'arriverait jamais dans ta marmite le soir. Moderniser la pêche, c'est donc aussi équiper et former ces femmes.
\end{savaistu}

\courssub{Effets et limites de la modernisation}

\textbf{Les effets positifs} sont nets :
\begin{itemize}
\item \textbf{augmentation des captures} et des \textbf{revenus} des pêcheurs ;
\item meilleur \textbf{approvisionnement des villes} en poisson frais et fumé (Yaoundé, Douala, Bafoussam) ;
\item création d'emplois en amont et en aval (mécaniciens, fabricants de glace, mareyeuses, transporteurs).
\end{itemize}

\textbf{Mais les limites sont réelles} :
\begin{itemize}
\item la \textbf{surpêche} : l'efficacité accrue des engins épuise les stocks côtiers ;
\item la \textbf{concurrence des pêches industrielles} (chalutiers, parfois étrangers) qui raclent les fonds marins, et celle des \textbf{importations} de poisson congelé bon marché qui concurrencent le poisson local ;
\item l'\textbf{insécurité en mer} : piraterie dans le golfe de Guinée, accidents, intempéries ;
\item le \textbf{coût des intrants} (moteurs, carburant, filets) qui endette les pêcheurs.
\end{itemize}

\begin{exoresolu}[Construire le croquis d'un site de pêche artisanale modernisé]
\textbf{Énoncé.} À partir de la visite (réelle ou imaginée) du débarcadère de Youpwè (Douala), construis un croquis géographique légendé d'un site de pêche artisanale modernisé.
\tcblower
\textbf{Solution détaillée.} On applique la méthode du croquis (voir boîte méthode) : esquisse à main levée, titre en majuscules, légende classée, orientation. Le croquis ci-dessous est un modèle attendu.

\begin{popfigure}
\begin{center}
\begin{tikzpicture}[scale=1.0, font=\small]
% Mer
\fill[popblue!25] (-0.5,2.6) rectangle (9.5,5.2);
\node[popblue, font=\itshape] at (4.5,4.85) {Mer (estuaire du Wouri)};
% Plage
\fill[popgold!40] (-0.5,2.1) rectangle (9.5,2.6);
% Terre
\fill[popgreen!15] (-0.5,-0.6) rectangle (9.5,2.1);
% Pirogues motorisées
\draw[popdark, thick, fill=poporange!50] (1.2,3.3) .. controls (1.6,3.0) and (2.6,3.0) .. (3.0,3.3) -- cycle;
\draw[popdark, thick] (3.0,3.3) -- (3.25,3.15);
\node[font=\scriptsize] at (2.1,3.6) {1};
\draw[popdark, thick, fill=poporange!50] (5.2,3.7) .. controls (5.6,3.4) and (6.6,3.4) .. (7.0,3.7) -- cycle;
\draw[popdark, thick] (7.0,3.7) -- (7.25,3.55);
\node[font=\scriptsize] at (6.1,4.0) {1};
% Débarcadère
\draw[popdark, very thick, fill=popmuted!40] (3.6,2.1) rectangle (4.4,3.0);
\node[font=\scriptsize] at (4.0,1.85) {2};
% Glacières / chambre froide
\draw[popdark, thick, fill=popblue!40] (5.2,1.2) rectangle (6.6,2.1);
\node[font=\scriptsize] at (5.9,1.65) {3};
% Fours de fumage
\draw[popdark, thick, fill=poppink!40] (1.0,1.0) rectangle (2.4,2.1);
\draw[popdark, thick] (1.3,2.1) -- (1.3,2.45); \draw[popdark, thick] (2.1,2.1) -- (2.1,2.45);
\node[font=\scriptsize] at (1.7,1.55) {4};
% Hangar mareyeuses
\draw[popdark, thick, fill=popgreen!40] (7.2,1.2) rectangle (8.8,2.1);
\node[font=\scriptsize] at (8.0,1.65) {5};
% Piste vers le marché
\draw[popdark, dashed, very thick] (4.0,2.1) -- (4.0,0.2) -- (8.6,0.2) -- (9.3,-0.3);
\node[font=\scriptsize, anchor=west] at (8.6,-0.45) {6 : piste vers le marché de la ville};
% Nord
\draw[-{Stealth}, thick, popdark] (9.0,4.4) -- (9.0,5.0);
\node[font=\scriptsize] at (9.0,5.15) {N};
\end{tikzpicture}
\end{center}
\legende{CROQUIS : UN SITE DE PÊCHE ARTISANALE MODERNISÉ (ex. Youpwè, Douala). Légende classée : \textbf{Hydrographie} : mer et plage. \textbf{Activités} : 1 = pirogues motorisées (moteur hors-bord) ; 4 = fours de fumage améliorés ; 5 = hangar des mareyeuses (transformation et vente). \textbf{Infrastructures} : 2 = débarcadère ; 3 = chambre froide et glacières ; 6 = piste vers le marché urbain. Orientation : flèche du nord.}
\end{popfigure}

\textbf{Critères de réussite au Bac} : titre en majuscules, légende classée en catégories (hydrographie, infrastructures, activités), flèche du nord, traits simples et propres.
\end{exoresolu}

\begin{methode}[Construire un croquis géographique]
\begin{enumerate}
\item \textbf{Esquisse à main levée} : dessine d'abord les grandes lignes (côte, plan d'eau, relief) au crayon, sans chercher la précision cartographique.
\item \textbf{Titre en majuscules}, placé au-dessus du croquis, qui annonce le lieu et le thème.
\item \textbf{Légende classée} : regroupe les éléments par catégories (hydrographie, relief, infrastructures, activités) avec des symboles simples et constants.
\item \textbf{Orientation} : place toujours la flèche du nord ; ajoute si possible une échelle indicative.
\item \textbf{Sobriété} : peu de couleurs (2 ou 3), pas de détails inutiles ; chaque élément dessiné doit figurer dans la légende, et réciproquement.
\end{enumerate}
\end{methode}

\begin{exercice}[Pour t'entraîner]
Sur la carte schématique des zones de pêche du Cameroun présentée plus haut, reproduis le contour du pays et légende-la en plaçant : deux ports de pêche maritime, deux zones de pêche continentale, un fleuve et un lac de barrage. Précise ensuite, en cinq lignes, pourquoi la conservation est le maillon décisif de la modernisation de la pêche artisanale.
\end{exercice}

\begin{corrige}[Exercice]
\textbf{Carte} : on attend par exemple Kribi et Limbé (pêche maritime), le lac Tchad et le lac Maga (pêche continentale), la Sanaga ou le Logone (fleuve), le lac Lagdo (lac de barrage). Toute localisation cohérente est acceptée.
\textbf{Question} : le poisson est un produit très périssable sous climat tropical. Sans glace, chambre froide, fumage ou séchage, une partie des captures se gâte avant d'atteindre le consommateur : le pêcheur perd son revenu et les villes sont mal approvisionnées. La conservation permet donc de valoriser la capture, d'acheminer le poisson vers les marchés éloignés et de stabiliser les prix. C'est le maillon qui transforme une bonne pêche en revenu réel.
\end{corrige}

\begin{aretenir}[L'essentiel de la section]
\begin{itemize}
\item Le Cameroun pêche en \textbf{mer} (Limbé, Kribi, Douala, Campo) et en \textbf{eaux intérieures} (lac Tchad, lac Maga, Sanaga, Wouri, Logone, Lagdo).
\item La modernisation repose sur quatre piliers : \textbf{motorisation} des pirogues, \textbf{engins adaptés} (filets maillants, sennes, palangres), \textbf{conservation} (glace, chambres froides, fumage amélioré, séchage) et \textbf{organisation} (coopératives, GIC, mareyeuses).
\item Effets : hausse des captures et des revenus, approvisionnement des villes. Limites : \textbf{surpêche}, concurrence industrielle et des importations, insécurité en mer.
\item La pêche durable exige le respect des \textbf{mailles réglementaires} pour protéger les juvéniles.
\end{itemize}
\end{aretenir}

La pêche en mer et en eaux intérieures ne suffisant plus à couvrir la demande nationale, le Cameroun s'est tourné vers une solution complémentaire : produire le poisson plutôt que le capturer. C'est l'essor de la \textbf{pisciculture}, que nous étudions dans la section suivante.

\coursec{L'essor de la pisciculture}

La modernisation des équipements de pêche artisanale, étudiée dans la section précédente, a permis d'augmenter les captures sur les fleuves, les lacs et le littoral. Mais elle se heurte à une limite : les ressources naturelles ne sont pas inépuisables, et la demande en poisson des villes camerounaises croît plus vite que les captures. Face à ce constat, une solution s'impose : au lieu de seulement \emph{capturer} le poisson, pourquoi ne pas l'\emph{élever}, comme on élève des poulets ou des chèvres ? C'est exactement le principe de la \cle{pisciculture}, dont l'essor constitue l'une des transformations majeures des activités halieutiques au Cameroun.

\courssub{Définition et formes de la pisciculture}

\textbf{De l'intuition à la définition.} Imagine un agriculteur de Foumban ou de Bafoussam qui creuse une mare derrière sa concession, la remplit d'eau de source, y introduit de jeunes poissons, les nourrit chaque jour avec des sons de riz et des tourteaux, puis vide la mare six mois plus tard pour récolter des poissons de belle taille. Il ne « pêche » plus : il \emph{produit} du poisson, selon un calendrier qu'il maîtrise. C'est toute la différence entre la pêche et la pisciculture.

\begin{definition}[La pisciculture]
La \cle{pisciculture} est l'\textbf{élevage de poissons en milieu contrôlé} (étangs, cages, bassins, rizières aménagées) en vue d'une production \textbf{commerciale} ou \textbf{familiale}. L'éleveur, appelé \emph{pisciculteur}, contrôle l'alimentation des poissons, la qualité de l'eau, la reproduction (via l'alevinage) et la récolte.
\end{definition}

\begin{definition}[L'étang piscicole]
Un \cle{étang} piscicole est un \textbf{plan d'eau artificiel aménagé et géré pour l'élevage de poissons}. Il est construit avec une digue de retenue, un système d'arrivée et d'évacuation de l'eau (le \emph{moine}), et permet le contrôle de l'alimentation, de la qualité de l'eau et de la récolte (par vidange totale ou partielle).
\end{definition}

\textbf{Les formes de pisciculture pratiquées au Cameroun.} On distingue plusieurs systèmes, selon le milieu d'élevage et l'objectif de production :

\begin{itemize}
\item la \textbf{pisciculture en étangs} : c'est la forme la plus répandue. Des étangs de quelques centaines de mètres carrés sont creusés dans les zones où l'eau est disponible (sources, ruisseaux, nappes peu profondes) ;
\item la \textbf{pisciculture en cages flottantes} : des cages en filet sont ancrées dans des plans d'eau naturels (lacs, retenues de barrages, bras de fleuve). Les poissons y sont nourris et grossissent en milieu confiné, sans creusement d'étang ;
\item la \textbf{rizipisciculture} : on associe la culture du riz et l'élevage de poissons dans la même parcelle inondée. Les poissons mangent les insectes et les herbes, et leurs déjections fertilisent le riz : c'est un système intégré, économique et écologique ;
\item la \textbf{pisciculture en bassins} (bâchés ou en béton), plus récente, pratiquée en zone périurbaine, même sur de petites surfaces.
\end{itemize}

Selon l'échelle, on parle de \textbf{pisciculture familiale} (petit étang destiné d'abord à la consommation du ménage, le surplus étant vendu au marché local) et de \textbf{pisciculture commerciale} (plusieurs étangs, investissements importants en aliments et en alevins, production destinée à la vente régulière en ville).

\begin{attention}[Ne pas confondre pêche continentale et pisciculture]
Erreur très fréquente dans les copies : la \textbf{pêche continentale} consiste à \emph{capturer} des poissons sauvages dans les eaux naturelles (fleuves, lacs, lagunes), tandis que la \textbf{pisciculture} consiste à \emph{élever} des poissons dans un milieu aménagé et contrôlé. L'une est une activité de \emph{cueillette sur ressource naturelle}, l'autre une activité d'\emph{élevage}. Au Baccalauréat, cette distinction doit apparaître clairement dès l'introduction de ta copie.
\end{attention}

\courssub{Les espèces élevées et la question de l'alevinage}

\textbf{Des poissons choisis pour leur rusticité.} Toutes les espèces ne se prêtent pas à l'élevage. Le pisciculteur camerounais privilégie des poissons qui supportent bien le confinement, acceptent une alimentation artificielle et grossissent vite :

\begin{itemize}
\item le \textbf{tilapia} (\emph{Oreochromis niloticus} notamment) : l'espèce reine de la pisciculture africaine. Omnivore, résistant, il se reproduit facilement en étang ;
\item la \textbf{carpe} (carpe commune, carpes chinoises) : introduite au Cameroun, elle valorise bien les végétaux et les déchets agricoles ;
\item le \textbf{clarias} ou poisson-chat (\emph{Clarias gariepinus}) : très rustique, il respire l'air atmosphérique grâce à un organe accessoire, supporte des eaux peu oxygénées et atteint de fortes densités d'élevage ;
\item le \textbf{capitaine} et d'autres espèces locales sont testés ou élevés à plus petite échelle pour diversifier la production.
\end{itemize}

\textbf{L'alevinage : la clé de voûte du système.} Pour peupler ses étangs, le pisciculteur a besoin d'\textbf{alevins}, c'est-à-dire de jeunes poissons de quelques centimètres. Ces alevins sont produits dans des \textbf{stations piscicoles publiques} (stations de recherche et de multiplication gérées par l'État, comme celles de Foumban dans l'Ouest ou de Yaoundé-Mvolyé dans le Centre) et, de plus en plus, par des écloseries privées. La qualité et la disponibilité des alevins conditionnent toute la filière : sans bons alevins, pas de bonne récolte.

\begin{savaistu}[Deux récoltes par an grâce au tilapia]
Dans un étang bien géré (bonne alimentation, densité maîtrisée, eau renouvelée), le tilapia peut atteindre la taille marchande (environ 250 à 400 g) en \textbf{6 à 8 mois}. Un pisciculteur organisé peut donc réaliser \textbf{deux cycles de production par an}, ce qui fait de la pisciculture une source de revenus régulière, contrairement à bien des cultures saisonnières.
\end{savaistu}

\courssub{Comprendre le fonctionnement d'un étang piscicole}

Pour bien saisir la pisciculture en étang, il faut comprendre son infrastructure. Un étang n'est pas une simple mare : c'est un \emph{dispositif technique} pensé pour contrôler l'eau et les poissons.

\begin{popfigure}
\begin{center}
\begin{tikzpicture}[scale=1.0, every node/.style={font=\small}]
% sol
\fill[popgoldL] (-6.5,0) rectangle (6.5,-0.4);
\draw[popdark] (-6.5,0) -- (6.5,0);
% creux de l'étang
\fill[popblueL] (-4.2,0) -- (-4.2,-2.2) -- (4.2,-2.2) -- (4.2,0) -- cycle;
\draw[popdark, thick] (-4.2,0) -- (-4.2,-2.2) -- (4.2,-2.2) -- (4.2,0);
% digues
\fill[popgold] (-5.6,0) rectangle (-4.2,0.9);
\fill[popgold] (4.2,0) rectangle (5.6,0.9);
\draw[popdark, thick] (-5.6,0) rectangle (-4.2,0.9);
\draw[popdark, thick] (4.2,0) rectangle (5.6,0.9);
% moine (tuyau d'évacuation)
\draw[line width=3pt, popmuted] (4.2,-1.6) -- (5.9,-1.6);
\draw[popdark, thick, ->] (5.9,-1.6) -- (6.4,-1.6);
\node[popdark, anchor=west] at (6.4,-1.6) {évacuation};
% arrivée d'eau
\draw[popblue, thick, ->] (-6.4,0.5) -- (-4.6,0.1);
\node[popblue, anchor=east] at (-6.4,0.5) {arrivée d'eau};
% poissons
\foreach \x/\y in {-2.5/-1.2, -0.8/-1.6, 1.0/-1.0, 2.6/-1.5}{
  \draw[popink] (\x,\y) ellipse (0.28 and 0.1);
  \draw[popink] (\x+0.28,\y) -- (\x+0.45,\y+0.09) -- (\x+0.45,\y-0.09) -- cycle;
}
% granulés (zone d'alimentation)
\foreach \p in {(-3.2,-0.25),(-3.0,-0.3),(-2.8,-0.22),(-3.1,-0.15),(-2.9,-0.35)}{
  \fill[poporange] \p circle (0.05);
}
% aérateur
\draw[popgreen, thick] (0.5,0.15) circle (0.18);
\draw[popgreen, thick, ->] (0.5,0.33) -- (0.5,0.8);
\draw[popgreen, thick, ->] (0.32,0.24) -- (-0.1,0.65);
\draw[popgreen, thick, ->] (0.68,0.24) -- (1.1,0.65);
% labels
\node[popdark, font=\small\bfseries] at (-4.9,1.15) {digue};
\node[popdark, font=\small\bfseries] at (4.9,1.15) {digue};
\node[popmuted, font=\small] at (5.0,-1.95) {moine};
\node[poporange, font=\small] at (-3.0,0.35) {zone d'alimentation};
\node[popgreen, font=\small] at (0.5,1.05) {aérateur};
\node[popink, font=\small] at (-0.8,-2.5) {poissons en croissance};
\node[popblue, font=\small] at (2.6,-0.35) {eau (1 à 1,5 m)};
\end{tikzpicture}
\end{center}
\legende{Schéma en coupe d'un étang piscicole. La \textbf{digue} retient l'eau ; le \textbf{moine} (tuyau muni d'une grille) permet la vidange et l'évacuation des eaux de surface sans perte de poissons ; l'\textbf{aérateur} oxygène l'eau lorsque la densité de poissons est forte ; la \textbf{zone d'alimentation} concentre la distribution des granulés. Croquis schématique, sans échelle.}
\end{popfigure}

\begin{methode}[Conduire un cycle de production en étang]
\begin{enumerate}
\item \textbf{Aménager} : construire la digue et le moine, niveler le fond en pente douce vers le moine pour permettre la vidange complète.
\item \textbf{Mettre en eau et fertiliser} : remplir l'étang et, si besoin, le fertiliser pour développer le plancton (nourriture naturelle).
\item \textbf{Empoissonner} : introduire des alevins sains, à densité adaptée (par exemple 2 à 3 alevins de tilapia par mètre carré en système semi-intensif).
\item \textbf{Nourrir et surveiller} : distribuer chaque jour un aliment (granulés, sons, tourteaux), contrôler la couleur et la qualité de l'eau.
\item \textbf{Récolter} : après 6 à 8 mois, pêcher au filet puis vidanger l'étang pour la récolte totale ; trier, vendre, et remettre un nouveau cycle en route.
\end{enumerate}
\end{methode}

\courssub{Localisation de la pisciculture au Cameroun}

La pisciculture camerounaise se concentre là où se combinent \textbf{disponibilité en eau}, \textbf{relief favorable au creusement des étangs} et \textbf{proximité des marchés urbains} :

\begin{itemize}
\item l'\textbf{Ouest} : région pionnière et premier pôle piscicole du pays, autour de Foumban, Bafoussam, Dschang et Bangangté. Le relief accidenté y est traversé de nombreux ruisseaux, et la station piscicole de Foumban a historiquement diffusé alevins et techniques ;
\item le \textbf{Centre} : forte dynamique autour de Yaoundé et de sa périphérie (Mbankomo, Obala, Sa'a), portée par l'énorme demande de la capitale ;
\item le \textbf{Sud} : étangs autour d'Ebolowa, Sangmélima et Ambam ;
\item le \textbf{Littoral} : pisciculture périurbaine près de Douala et essais de cages flottantes, au plus près du plus grand marché de poisson du pays.
\end{itemize}

La logique de localisation est claire : le poisson d'élevage se vend frais, sans chaîne du froid développée ; il faut donc produire \emph{près des villes consommatrices}. C'est pourquoi les zones périurbaines à forte demande (Yaoundé, Douala, Bafoussam, Bamenda) concentrent les investissements.

\courssub{Atouts et contraintes de la pisciculture camerounaise}

\begin{center}
\begin{tabularx}{\linewidth}{|X|X|}
\hline
\rowcolor{popblueL} \textbf{Atouts} & \textbf{Contraintes} \\
\hline
Production \textbf{régulière et contrôlée} (calendrier de récolte maîtrisé, 2 cycles/an possibles) & \textbf{Coût élevé des aliments} pour poissons (granulés souvent importés ou chers : l'aliment peut représenter plus de la moitié des charges) \\
\hline
\rowcolor{popcream} \textbf{Proximité des marchés urbains} : vente fraîche, prix rémunérateurs & \textbf{Qualité de l'eau} difficile à gérer (pollution, manque d'oxygène, sécheresses) \\
\hline
Création d'\textbf{emplois et de revenus} ruraux et périurbains (jeunes, femmes, GIC) & \textbf{Manque d'alevins de qualité} en quantité suffisante (écloseries insuffisantes) \\
\hline
\rowcolor{popcream} \textbf{Réduction de la pression} sur les pêches naturelles surexploitées & \textbf{Faible encadrement technique} : vulgarisation insuffisante, mortalités, rendements décevants \\
\hline
\end{tabularx}
\end{center}

\begin{exemplebox}[Une production encore modeste face à une demande immense]
La production aquacole du Cameroun reste modeste : de l'ordre de \textbf{quelques milliers de tonnes par an} (ordre de grandeur : environ 5 000 à 10 000 tonnes selon les années, chiffres à manier avec prudence). Or la demande nationale en poisson est estimée à \textbf{plusieurs centaines de milliers de tonnes} (environ 400 000 à 500 000 tonnes par an), couverte en grande partie par la pêche et surtout par d'importantes \textbf{importations} de poissons congelés. L'écart entre l'offre d'élevage et la demande montre l'ampleur du \textbf{potentiel} : chaque tonne produite en étang est une tonne de moins à importer.
\end{exemplebox}

\begin{popfigure}
\begin{center}
\begin{tikzpicture}
\begin{axis}[
  ybar, bar width=0.55cm,
  width=0.85\linewidth, height=6cm,
  ymin=0, ymax=12,
  ylabel={Production (milliers de tonnes)},
  xlabel={Année},
  symbolic x coords={2000,2005,2010,2015,2020},
  xtick=data,
  ytick={0,2,4,6,8,10,12},
  nodes near coords,
  every node near coord/.append style={font=\scriptsize},
  axis line style={popdark},
  tick label style={font=\small},
  label style={font=\small},
]
\addplot[fill=popblue, draw=popdark] coordinates {(2000,1) (2005,1.5) (2010,3) (2015,5.5) (2020,8)};
\end{axis}
\end{tikzpicture}
\end{center}
\legende{Évolution \emph{indicative} de la production piscicole du Cameroun (ordres de grandeur, en milliers de tonnes). La croissance est réelle mais lente : la production reste très inférieure à la demande nationale (plusieurs centaines de milliers de tonnes). Source : estimations d'après données FAO et MINEPIA, valeurs arrondies.}
\end{popfigure}

\courssub{Pisciculture, sécurité alimentaire et réduction des importations}

L'enjeu dépasse le seul revenu des pisciculteurs. Le poisson est la \textbf{principale source de protéines animales} d'une large partie des Camerounais : en produire davantage sur place, c'est renforcer la \textbf{sécurité alimentaire} du pays. Chaque année, le Cameroun importe des centaines de milliers de tonnes de poissons congelés (maquereaux, sardinelles, capitaines), ce qui pèse lourdement sur la \textbf{balance commerciale} : les dépenses d'importation de poisson se chiffrent en dizaines de milliards de FCFA par an (ordre de grandeur). Développer la pisciculture, c'est donc à la fois :

\begin{itemize}
\item améliorer l'approvisionnement régulier des villes en poisson frais et accessible ;
\item créer des emplois ruraux et périurbains tout au long de la filière (alevinage, alimentation, élevage, commercialisation) ;
\item \textbf{réduire la facture des importations} et préserver les devises du pays ;
\item soulager les stocks naturels surexploités des fleuves, des lacs et du littoral.
\end{itemize}

C'est précisément pour accompagner cet essor que l'État camerounais a mis en place des programmes d'encadrement et de financement du secteur agro-pastoral et piscicole : ce sera l'objet du \textbf{Dossier 1}.

\begin{exoresolu}[Calculer le potentiel de substitution aux importations]
\textbf{Énoncé.} La demande nationale de poisson est estimée à 450 000 tonnes par an. La pêche (maritime et continentale) en fournit environ 200 000 tonnes, la pisciculture 8 000 tonnes, et le reste est importé. (1) Calcule la quantité de poisson importée. (2) Calcule la part de la pisciculture dans l'approvisionnement national. (3) Si la production piscicole était multipliée par cinq, de combien diminuerait la part des importations dans la demande ?
\tcblower
\textbf{Solution détaillée.}
\begin{enumerate}
\item Quantité importée : $450\,000 - 200\,000 - 8\,000 = 242\,000$ tonnes.
\item Part de la pisciculture : $\dfrac{8\,000}{450\,000} \times 100 \approx \num{1,8}\,\%$. La pisciculture ne couvre donc qu'environ 2 \% de la demande nationale.
\item Production piscicole multipliée par cinq : $8\,000 \times 5 = 40\,000$ tonnes. Les importations tomberaient à $242\,000 - 32\,000 = 210\,000$ tonnes, soit une part de $\dfrac{210\,000}{450\,000} \times 100 \approx \num{46,7}\,\%$ contre $\dfrac{242\,000}{450\,000} \times 100 \approx \num{53,8}\,\%$ auparavant. La part des importations diminuerait donc d'environ \textbf{7 points de pourcentage}. Conclusion : même un essor spectaculaire de la pisciculture ne suffirait pas, seul, à supprimer les importations, mais il les réduirait sensiblement tout en créant des milliers d'emplois.
\end{enumerate}
\end{exoresolu}

\begin{exercice}[À toi de jouer]
Un pisciculteur de Bafoussam possède deux étangs de 400 m$^2$ chacun. Il empoissonne à raison de 3 alevins de tilapia par m$^2$ et récolte, après 7 mois, 80 \% des poissons introduits, d'un poids moyen de 300 g. (1) Calcule le nombre d'alevins introduits. (2) Calcule la masse totale de poisson récoltée. (3) À 2 000 FCFA le kilogramme, calcule le chiffre d'affaires du cycle. (4) Dégage deux contraintes qui pourraient réduire ce résultat en réalité.
\end{exercice}

\begin{corrige}[Exercice : le pisciculteur de Bafoussam]
\begin{enumerate}
\item Surface totale : $2 \times 400 = 800$ m$^2$. Nombre d'alevins : $800 \times 3 = 2\,400$ alevins.
\item Poissons récoltés : $2\,400 \times \num{0,80} = 1\,920$ poissons. Masse totale : $1\,920 \times \num{0,3} = 576$ kg.
\item Chiffre d'affaires : $576 \times 2\,000 = 1\,152\,000$ FCFA pour le cycle.
\item Contraintes possibles : mortalités supérieures dues à une mauvaise qualité de l'eau ou à des alevins fragiles ; coût élevé de l'aliment qui réduit la marge bénéficiaire ; prédateurs (oiseaux, serpents) ; difficultés de vente si le marché local est saturé au moment de la récolte.
\end{enumerate}
\end{corrige}

\begin{aretenir}[L'essentiel de la section]
\begin{itemize}
\item La \cle{pisciculture} est l'élevage de poissons en milieu contrôlé ; l'\cle{étang} est un plan d'eau artificiel aménagé (digue, moine) pour cet élevage. Ne pas la confondre avec la pêche continentale.
\item Trois formes principales : étangs, cages flottantes, rizipisciculture ; pisciculture familiale ou commerciale.
\item Espèces élevées : tilapia, carpe, clarias, capitaine ; les alevins proviennent des stations piscicoles publiques (Foumban, Yaoundé-Mvolyé) et d'écloseries privées.
\item Localisation : Ouest (pôle pionnier), Centre, Sud, Littoral — toujours près des marchés urbains.
\item Atouts : production régulière, emplois, proximité des villes, moindre pression sur les pêches naturelles. Contraintes : coût des aliments, qualité de l'eau, pénurie d'alevins, faible encadrement.
\item Enjeu national : la production (quelques milliers de tonnes) reste très inférieure à la demande (plusieurs centaines de milliers de tonnes) ; développer la pisciculture, c'est lutter contre les importations et renforcer la sécurité alimentaire.
\end{itemize}
\end{aretenir}

\coursec{Dossier 1 : le rôle de l'État dans le développement des activités agro-pastorales et piscicoles}

Nous venons de voir que la pisciculture connaît un essor remarquable au Cameroun : étangs aménagés, alevins distribués, fermes piscicoles autour de Yaoundé, Douala ou Foumban. Mais cet essor n'est pas spontané. Derrière un étang qui produit du tilapia, il y a souvent une station d'alevinage publique, un technicien formé par l'État, une piste rurale qui permet d'écouler le poisson vers le marché. De même, derrière le zébu amélioré d'un éleveur de l'Adamaoua, il y a la recherche menée à Wakwa et les campagnes de vaccination vétérinaire. Ce dossier répond à une question centrale : \cle{quel rôle joue l'État camerounais dans le développement des activités agro-pastorales et piscicoles ?} Nous verrons successivement le cadre institutionnel, les programmes d'encadrement et de financement, la formation des acteurs, la régulation de la commercialisation, les infrastructures publiques, puis nous dresserons un bilan critique avant de relier l'ensemble à la problématique de l'intégration nationale.

\courssub{Le cadre institutionnel : qui encadre l'élevage et la pêche ?}

Pour comprendre l'action de l'État, il faut d'abord identifier ses instruments. Au Cameroun, deux ministères se partagent l'essentiel de l'encadrement agro-pastoral et piscicole.

\begin{definition}[Les principaux acteurs institutionnels]
Le \cle{MINEPIA} (Ministère de l'Élevage, des Pêches et des Industries Animales) conçoit et met en œuvre la politique nationale en matière d'élevage, de pêche et de production animale : santé animale, appui aux éleveurs et pêcheurs, contrôle des produits d'origine animale. Le \cle{MINADER} (Ministère de l'Agriculture et du Développement Rural) s'occupe de l'agriculture et du développement rural en général, notamment des pistes rurales et de la vulgarisation. L'\cle{IRAD} (Institut de Recherche Agricole pour le Développement) mène la recherche scientifique appliquée ; sa station de \cle{Wakwa}, près de Ngaoundéré dans l'Adamaoua, est spécialisée dans l'élevage et les systèmes pastoraux.
\end{definition}

À ces structures centrales s'ajoutent des services déconcentrés présents sur le terrain : les \cle{services vétérinaires} (vaccinations, contrôle sanitaire, prophylaxie), les délégations régionales et départementales du MINEPIA et du MINADER, les centres zootechniques et les stations piscicoles. L'État ne travaille pas seul : il s'appuie sur les \cle{coopératives} et les \cle{GIC} (Groupements d'Initiative Commune) des producteurs, sur les communes, et sur des partenaires extérieurs (Banque mondiale, FAO, BAD, coopérations bilatérales) qui cofinancent de nombreux projets.

\begin{savaistu}[Wakwa, une référence africaine]
La station de recherche de Wakwa (Adamaoua), créée dans les années 1960, est une référence en Afrique centrale pour l'amélioration des races bovines et des systèmes pastoraux. C'est là qu'ont été menés des travaux pionniers sur le croisement entre les zébus locaux (Goudali) et des races améliorées, sur l'adaptation des bovins à la trypanosomose (maladie transmise par la mouche tsé-tsé) et sur la gestion des parcours. Des chercheurs de toute la sous-région y sont formés.
\end{savaistu}

\begin{popfigure}
\begin{center}
\begin{tikzpicture}[
  font=\small,
  box/.style={draw=popblue, fill=popblueL, rounded corners=2pt, align=center, minimum height=0.9cm, text width=3.1cm},
  boxp/.style={draw=poporange, fill=poporangeL, rounded corners=2pt, align=center, minimum height=0.9cm, text width=3.1cm},
  boxg/.style={draw=popgreen, fill=popgreenL, rounded corners=2pt, align=center, minimum height=0.9cm, text width=3.1cm},
  boxm/.style={draw=poppurple, fill=poppurpleL, rounded corners=2pt, align=center, minimum height=0.9cm, text width=3.4cm},
  fl/.style={-{Stealth[length=2.5mm]}, thick, popdark},
  flb/.style={{Stealth[length=2.5mm]}-{Stealth[length=2.5mm]}, thick, popmuted}
]
\node[box, fill=popblue!40, text width=5.2cm] (etat) at (0,4.2) {\textbf{ÉTAT}\\ MINEPIA, MINADER, IRAD (Wakwa), services vétérinaires};
\node[boxp, align=center] (projets) at (-4.6,2.2){\textbf{Programmes et projets}\\ PDMA, projets élevage et pêche, microfinance rurale};
\node[boxp, align=center] (part) at (4.6,2.2){\textbf{Partenaires extérieurs}\\ Banque mondiale, FAO, BAD, coopérations};
\node[boxg, align=center] (coop) at (-4.6,0){\textbf{Coopératives et GIC}\\ d'éleveurs, pêcheurs, pisciculteurs};
\node[boxg, align=center] (prod) at (0,0){\textbf{Producteurs}\\ éleveurs, pêcheurs artisans, pisciculteurs};
\node[boxg, align=center] (form) at (4.6,0){\textbf{Formation et encadrement}\\ vulgarisation, techniciens, vétérinaires};
\node[boxm, align=center] (marche) at (0,-2.2){\textbf{Marchés}\\ marchés à bétail, mareyeurs, filières viande--lait--poisson};
\draw[fl] (etat) -- (projets);
\draw[fl] (etat) -- (part);
\draw[fl] (etat) -- (prod);
\draw[fl] (projets) -- (coop);
\draw[fl] (part) -- (form);
\draw[fl] (coop) -- (marche);
\draw[fl] (prod) -- (marche);
\draw[fl] (form) -- (prod);
\draw[flb] (coop) -- (prod);
\end{tikzpicture}
\end{center}
\legende{Organigramme des acteurs de l'encadrement agro-pastoral et piscicole au Cameroun. L'État (en haut) agit directement sur les producteurs et indirectement via les programmes, les coopératives et la formation ; l'ensemble débouche sur les marchés. (Schéma réalisé par l'auteur.)}
\end{popfigure}

\courssub{Les programmes et projets d'encadrement et de financement}

L'État intervient rarement « à mains nues » : il agit le plus souvent à travers des \cle{programmes} et des \cle{projets} dotés d'un budget, d'un calendrier et de zones d'intervention précises.

\begin{definition}[PDMA]
Le \cle{PDMA} (Programme de Développement des Marchés Agricoles) est un programme public, cofinancé avec des partenaires internationaux, qui vise à améliorer l'\cle{accès des producteurs aux marchés}. Ses actions principales : la réhabilitation et l'ouverture de \cle{pistes rurales}, la construction d'infrastructures de commercialisation (marchés, magasins de stockage, aires de vente) et l'appui au financement des activités des producteurs organisés.
\end{definition}

L'intuition est simple : un éleveur de l'Extrême-Nord peut produire les plus beaux bœufs du pays, si la piste qui mène à son village est impraticable en saison des pluies, ses animaux n'atteindront jamais les marchés de Yaoundé ou de Douala dans de bonnes conditions. Le PDMA attaque précisément ce goulot d'étranglement : le \cle{désenclavement} et la \cle{commercialisation}.

D'autres dispositifs complètent ce paysage :

\begin{itemize}
\item les \cle{projets de développement de l'élevage} (appui à l'aviculture villageoise, à l'engraissement bovin, à la production laitière dans l'Ouest et l'Adamaoua) ;
\item les \cle{projets de développement de la pêche et de l'aquaculture} (stations d'alevinage, appui aux pêcheurs artisans du littoral et du lac Tchad, distribution d'équipements) ;
\item la \cle{microfinance rurale} : les MC2 (Mutuelles Communautaires de Croissance), nées dans les années 1990, accordent de petits crédits aux producteurs qui n'ont pas accès aux banques classiques ;
\item les \cle{subventions d'équipements} : motopompes, filets, moteurs hors-bord, matériel de forage, distribués gratuitement ou à prix réduit dans le cadre de projets.
\end{itemize}

\begin{exemplebox}[Des résultats concrets sur le terrain]
Les programmes publics ont financé des \cle{centaines de kilomètres de pistes rurales} et des \cle{dizaines de marchés et d'infrastructures de conservation} (magasins, chambres froides, aires de séchage du poisson) dans les zones de production. À titre d'ordre de grandeur, le seul PDMA a concerné plusieurs centaines de kilomètres de pistes et plusieurs dizaines d'infrastructures de marché dans de nombreuses communes rurales ; les chiffres exacts varient selon les phases du programme. Résultat observé : le temps de transport des produits vers les marchés diminue, les pertes après récolte ou après pêche baissent, et les prix payés au producteur s'améliorent.
\end{exemplebox}

\courssub{La formation des acteurs : transformer les pratiques}

Moderniser, ce n'est pas seulement donner du matériel : c'est d'abord \cle{changer les savoir-faire}. L'État déploie pour cela un dispositif de formation et d'encadrement.

\begin{itemize}
\item Les \cle{centres de formation} (centres de formation professionnelle agricole, lycées techniques agricoles, centres zootechniques) forment de jeunes éleveurs, pêcheurs et pisciculteurs aux techniques modernes : conduite d'un poulailler, fabrication d'aliments pour poissons, gestion d'un étang.
\item L'\cle{encadrement vétérinaire} assure les campagnes de vaccination (contre la peste bovine, la maladie de Newcastle en aviculture, la pasteurellose), les traitements et le contrôle sanitaire des troupeaux et des produits.
\item L'\cle{encadrement halieutique} (relatif à la pêche) conseille les pêcheurs sur les tailles de mailles réglementaires, les périodes de reproduction à respecter et les techniques de conservation.
\item La \cle{vulgarisation agricole} désigne l'ensemble des actions qui diffusent les innovations vers les producteurs : démonstrations en plein champ, émissions rurales à la radio, visites-conseils des techniciens.
\item L'appui aux \cle{coopératives et GIC} permet aux producteurs de se regrouper pour acheter les intrants moins cher, mutualiser le matériel et négocier les prix de vente.
\end{itemize}

\begin{attention}[Erreur fréquente : une vision simpliste de la modernisation]
Dans les copies, on lit souvent que « l'État a tout fait » ou, à l'inverse, que « les producteurs se sont modernisés seuls ». C'est faux dans les deux cas. La modernisation de l'élevage et de la pêche résulte de l'\cle{interaction} entre quatre catégories d'acteurs : l'État (lois, programmes, infrastructures), les acteurs privés (commerçants, mareyeurs, fabricants d'aliments), les coopératives et GIC (organisation des producteurs) et les partenaires extérieurs (financements, expertise). Une bonne copie de Baccalauréat montre toujours cette complémentarité.
\end{attention}

\courssub{La régulation de la commercialisation : organiser les filières}

\begin{definition}[Filière]
Une \cle{filière} est l'ensemble des étapes et des acteurs qui mènent un produit de la production à la consommation : production, collecte, transformation, transport, distribution et vente. Parler de « filière viande », « filière lait » ou « filière poisson », c'est raisonner sur toute la chaîne, du pâturage ou de l'étang jusqu'à l'assiette du consommateur.
\end{definition}

L'État intervient à chaque maillon pour organiser et sécuriser ces filières :

\begin{enumerate}
\item \textbf{Organisation des marchés à bétail} : les grands marchés à bétail (par exemple ceux du Nord et de l'Adamaoua, comme le marché de Figuil ou les points de collecte autour de Ngaoundéré) sont aménagés et contrôlés ; les taxes et les règles de transaction y sont fixées.
\item \textbf{Normes sanitaires} : inspection vétérinaire des animaux avant abattage, contrôle de l'hygiène dans les abattoirs et les étals, lutte contre les fraudes. Objectif : protéger la santé du consommateur et valoriser le produit local.
\item \textbf{Gestion des importations de poisson} : le Cameroun importe chaque année d'importantes quantités de poisson congelé (de l'ordre de plusieurs centaines de milliers de tonnes, pour un coût de l'ordre de centaines de milliards de FCFA — ordres de grandeur à manier avec prudence) afin de combler le déficit entre la production nationale et la demande. L'État régule ces importations (quotas, autorisations) tout en cherchant à développer la production locale, notamment piscicole, pour les réduire.
\item \textbf{Stabilisation des prix et de l'approvisionnement} : encadrement des circuits de distribution, appui aux mareyeuses et aux commerçants, information sur les prix.
\end{enumerate}

\courssub{Les infrastructures publiques : l'armature du développement}

Aucune modernisation n'est possible sans infrastructures. L'État et ses partenaires investissent dans :

\begin{itemize}
\item les \cle{forages pastoraux} : points d'eau modernes creusés dans les zones sèches du Nord et de l'Extrême-Nord, qui permettent aux troupeaux de boire en saison sèche sans parcourir des dizaines de kilomètres ; ils structurent aussi les itinéraires de transhumance ;
\item les \cle{pistes rurales} : elles désenclavent les bassins de production et relient les villages aux marchés (action centrale du PDMA) ;
\item les \cle{abattoirs} modernes ou réhabilités (abattoirs de Yaoundé, Douala, Ngaoundéré, Maroua), qui garantissent une viande saine et contrôlée ;
\item les \cle{chambres froides} et unités de conservation, qui réduisent les pertes de poisson et de viande, particulièrement élevées sous climat chaud ;
\item les \cle{stations d'alevinage} publiques, qui produisent et distribuent les alevins (jeunes poissons) indispensables au démarrage des fermes piscicoles.
\end{itemize}

\courssub{Bilan critique : des progrès réels, des défis persistants}

\begin{methode}[Construire un tableau d'analyse d'un programme de développement]
Pour analyser rigoureusement un programme public (au Baccalauréat comme dans un dossier), suis ces étapes :
\begin{enumerate}
\item Identifier l'\cle{objectif} du programme (que cherche-t-il à changer ?).
\item Préciser les \cle{bénéficiaires} (quels producteurs, quelles régions ?).
\item Inventorier les \cle{moyens} mobilisés (financement, infrastructures, formation).
\item Évaluer les \cle{résultats} obtenus (chiffres, changements observés).
\item Discuter les \cle{limites} (ce qui n'a pas marché, ce qui reste à faire).
\end{enumerate}
C'est exactement la structure du tableau de synthèse ci-dessous, que tu dois savoir compléter.
\end{methode}

\begin{center}
\begin{tabularx}{\linewidth}{|l|X|X|X|}
\hline
\rowcolor{popblueL} \textbf{Programme / action} & \textbf{Objectifs} & \textbf{Zones d'intervention} & \textbf{Résultats et limites} \\
\hline
PDMA & Améliorer l'accès des producteurs aux marchés & Communes rurales de plusieurs régions & Pistes réhabilitées, marchés construits ; entretien des pistes parfois insuffisant \\
\hline
Projets élevage (aviculture, engraissement, lait) & Intensifier et moderniser la production animale & Ouest, Adamaoua, Nord & Chevillères, poulaillers et races améliorées diffusés ; coût des intrants élevé \\
\hline
Projets pêche et aquaculture & Augmenter la production halieutique et piscicole & Littoral, lac Tchad, zones intérieures & Stations d'alevinage, étangs aménagés ; production encore inférieure à la demande \\
\hline
Microfinance rurale (MC2) & Financer les petits producteurs exclus des banques & Milieu rural, surtout Nord-Ouest et Ouest à l'origine & Petits crédits accessibles ; montants limités, taux parfois élevés \\
\hline
Forages pastoraux et vétérinaire & Sécuriser l'eau et la santé des troupeaux & Nord, Extrême-Nord, Adamaoua & Mortalité animale réduite ; conflits d'accès aux points d'eau possibles \\
\hline
\end{tabularx}
\end{center}

Le bilan est donc \cle{nuancé}. Les progrès sont réels : cheptel vacciné, races améliorées, pisciculture en expansion, zones désenclavées, produits locaux mieux valorisés. Mais les défis persistent :

\begin{itemize}
\item le \cle{financement} reste insuffisant : les budgets publics sont limités et le crédit bancaire classique reste inaccessible à la plupart des petits producteurs ;
\item l'\cle{enclavement} de nombreuses zones de production n'est pas totalement levé, surtout en saison des pluies ;
\item la \cle{concurrence des importations} (poisson congelé, poulet importé) tire les prix vers le bas et décourage certains producteurs locaux ;
\item le \cle{changement climatique} (avancée de la sécheresse au Nord, irrégularité des pluies) menace les parcours et les points d'eau, et attise parfois les conflits entre agriculteurs et éleveurs.
\end{itemize}

\begin{exoresolu}[Analyser l'action de l'État à partir d'un document]
Énoncé. Un document indique : « Dans une commune de l'Adamaoua, un programme public a financé 45 km de pistes rurales, un marché à bétail modernisé et la formation de 120 éleveurs à l'engraissement bovin. Trois ans après, le nombre de bœufs vendus par les éleveurs de la commune a augmenté de 40 \% et le prix moyen au kilo de viande a progressé. » Question : à l'aide de la méthode d'analyse d'un programme, dégage l'objectif, les moyens, les résultats et une limite possible de cette action.
\tcblower
Solution détaillée.
\begin{enumerate}
\item \textbf{Objectif} : améliorer la commercialisation de la viande et les revenus des éleveurs (accès au marché et montée en gamme de la production).
\item \textbf{Bénéficiaires} : les éleveurs de la commune (120 formés directement, l'ensemble profitant du marché et des pistes).
\item \textbf{Moyens} : 45 km de pistes rurales (désenclavement), un marché à bétail modernisé (infrastructure de vente), une formation à l'engraissement (transfert de savoir-faire). On reconnaît les trois leviers classiques de l'action publique : infrastructure, équipement, formation.
\item \textbf{Résultats} : hausse de 40 \% des ventes de bœufs et amélioration du prix au kilo ; le désenclavement réduit les coûts de transport et la formation améliore la qualité des animaux, donc leur valeur.
\item \textbf{Limite possible} : la pérennité dépend de l'entretien des pistes (qui se dégradent vite sans maintenance) et de la capacité des éleveurs à acheter les intrants d'engraissement ; de plus, la concurrence de la viande ou du bétail venant d'autres zones peut limiter la hausse des prix.
\end{enumerate}
Conclusion rédigée : l'action de l'État est efficace parce qu'elle combine plusieurs moyens complémentaires, mais ses effets durables exigent un entretien continu et un financement pérenne.
\end{exoresolu}

\courssub{Lien avec l'intégration nationale : pourquoi ce dossier compte}

Ce dossier s'inscrit dans la famille de situations « l'intégration nationale » : l'action de l'État dans l'élevage et la pêche n'est pas seulement économique, elle est \cle{politique et territoriale}.

\begin{itemize}
\item \textbf{Valoriser les produits locaux} : en améliorant la qualité de la viande, du lait et du poisson camerounais, l'État encourage les consommateurs à préférer les produits nationaux aux importations — c'est le sens de la campagne « consommons local ».
\item \textbf{Désenclaver les régions du Nord} : grenier de l'élevage camerounais, l'Adamaoua, le Nord et l'Extrême-Nord sont loin des grandes villes de consommation du Sud. Pistes, forages et marchés à bétail intègrent ces régions à l'économie nationale et réduisent les inégalités territoriales.
\item \textbf{Réduire le déficit alimentaire} : chaque tonne de poisson ou de viande produite localement est une tonne de moins importée, donc des devises économisées et des emplois ruraux créés. Le développement de la pisciculture, encouragé par les stations d'alevinage publiques, répond directement à cet enjeu.
\item \textbf{Fixer les populations rurales} : en rendant l'élevage et la pêche rentables, l'État freine l'exode rural vers Yaoundé et Douala et équilibre l'occupation du territoire.
\end{itemize}

\begin{aretenir}[L'essentiel du dossier]
\begin{itemize}
\item L'encadrement repose sur un \cle{cadre institutionnel} : MINEPIA, MINADER, IRAD (station de Wakwa), services vétérinaires.
\item L'État agit par des \cle{programmes} (PDMA, projets élevage-pêche, microfinance rurale, subventions) qui combinent financement, infrastructures et formation.
\item Il \cle{forme} les acteurs (vulgarisation, centres de formation, coopératives) et \cle{régule} les filières viande, lait et poisson (marchés à bétail, normes sanitaires, gestion des importations).
\item Il construit les \cle{infrastructures} indispensables : forages pastoraux, pistes rurales, abattoirs, chambres froides, stations d'alevinage.
\item Le bilan est positif mais inachevé : financement insuffisant, enclavement, concurrence des importations, changement climatique.
\item Enjeu majeur : l'\cle{intégration nationale}, par la valorisation des produits locaux et le désenclavement des régions de production.
\end{itemize}
\end{aretenir}

\begin{exercice}[Pour t'entraîner]
\begin{enumerate}
\item Définis les termes suivants : filière, PDMA, vulgarisation agricole, forage pastoral.
\item Construis un tableau d'analyse complet (objectif, bénéficiaires, moyens, résultats, limites) d'un projet de station d'alevinage publique implantée près de ta localité.
\item Rédige un paragraphe argumenté (10 à 15 lignes) montrant comment l'action de l'État dans l'élevage contribue à l'intégration nationale du Cameroun.
\end{enumerate}
\end{exercice}

\begin{corrige}[Exercice]
\begin{enumerate}
\item \textbf{Filière} : ensemble des étapes et des acteurs qui mènent un produit de la production à la consommation (production, transformation, transport, vente). \textbf{PDMA} : Programme de Développement des Marchés Agricoles, visant à améliorer l'accès des producteurs aux marchés (pistes rurales, infrastructures de commercialisation, financement). \textbf{Vulgarisation agricole} : ensemble des actions de diffusion des innovations techniques vers les producteurs (démonstrations, conseils, radios rurales). \textbf{Forage pastoral} : point d'eau moderne creusé en zone d'élevage sèche pour abreuver les troupeaux, notamment en saison sèche.
\item Exemple attendu : \emph{objectif} : fournir des alevins de qualité aux pisciculteurs et réduire les importations de poisson ; \emph{bénéficiaires} : pisciculteurs et jeunes entrepreneurs de la zone ; \emph{moyens} : bassins, bâtiments, techniciens, budget public et partenaires ; \emph{résultats} : alevins distribués, étangs créés, production locale en hausse ; \emph{limites} : coût de fonctionnement, qualité variable des alevins, accès difficile si les pistes sont dégradées.
\item Éléments de réponse : l'État désenclave les régions d'élevage du Nord (pistes, forages, marchés à bétail), ce qui les relie aux marchés du Sud ; il améliore la qualité des produits locaux (vaccination, races améliorées, normes), ce qui encourage leur consommation face aux importations ; il crée des revenus ruraux qui fixent les populations et réduisent les inégalités régionales. Ainsi, l'élevage, activité traditionnelle du Nord, devient un maillon de l'économie nationale unifiée : c'est là le sens de l'intégration nationale.
\end{enumerate}
\end{corrige}

\coursec{TP guidé : cartographier et représenter les transformations}

Dans le dossier précédent, tu as découvert comment l'État, à travers des programmes comme le PDMA ou les actions du MINEPIA, accompagne la modernisation de l'élevage et de la pêche. Mais un géographe ne se contente pas de lire des rapports : il \cle{représente} l'espace. Savoir légender une carte, construire un croquis à partir d'une photographie ou transformer un tableau statistique en diagramme sont des savoir-faire exigés au Baccalauréat. Ces travaux pratiques (TP) te conduisent pas à pas, de l'observation à l'interprétation. Chaque TP suit la même logique : \textbf{observer} les documents, \textbf{représenter} selon les règles cartographiques, puis \textbf{interpréter} pour dégager les contrastes régionaux du Cameroun.

\courssub{TP 1 : légender une carte muette du Cameroun}

\textbf{Intuition.} Imagine que tu reçoives une carte muette du Cameroun — juste le contour du pays — et une liste d'informations : « les zébus sont surtout à l'Adamaoua et à l'Extrême-Nord, les pêcheurs de crevettes travaillent au large de Kribi et de Limbe, les étangs piscicoles se concentrent autour de Yaoundé et de Foumban ». Comment rendre tout cela visible d'un seul coup d'œil ? C'est exactement le travail du cartographe : choisir des \cle{symboles} simples et les placer au bon endroit.

\begin{methode}[Légender une carte en cinq étapes]
\begin{enumerate}
\item \textbf{Lire la consigne et repérer les informations à placer.} Souligne dans le sujet les faits géographiques (zones, axes, points) et classe-les par nature : surfaces (zones d'élevage), lignes (axes de transhumance), points (ports de pêche, étangs, ranchs).
\item \textbf{Tracer le fond de carte.} Reproduis le contour du Cameroun au crayon, place la \cle{flèche du nord} et une échelle indicative. Reporte quelques repères : Yaoundé, Douala, Garoua, le lac Tchad, l'océan Atlantique.
\item \textbf{Choisir les symboles.} Une couleur ou un tramage par zone, des flèches pour les flux (transhumance), des figurés ponctuels (ancre pour un port de pêche, petit rectangle bleu pour un étang, carré noir pour un ranch). Règle d'or : un symbole = une information.
\item \textbf{Placer les informations sur la carte.} Les zones d'élevage en aplats de couleur, les axes de transhumance en flèches orientées du nord vers le sud (mouvement de la saison sèche vers les pâturages et points d'eau), les ports sur le littoral, les étangs dans l'Ouest et le Centre, les ranchs dans l'Adamaoua.
\item \textbf{Rédiger le titre et la légende.} Titre en \textbf{majuscules}, précis (quoi, où, quand). Légende \textbf{classée et hiérarchisée} : d'abord les surfaces, puis les lignes, puis les points. Termine par la source des données.
\end{enumerate}
\end{methode}

\begin{popfigure}
\begin{center}
\begin{tikzpicture}[scale=1]
% Titre du modèle
\node[font=\bfseries, popdark] at (3.2,5.6) {MODÈLE DE LÉGENDE — ÉLEVAGE ET PÊCHE AU CAMEROUN};
% Encadré légende
\draw[popline, thick, fill=popcream, rounded corners] (0,0) rectangle (6.4,5.2);
\node[font=\bfseries\small, popblue, anchor=west] at (0.3,4.8) {Légende (classée : surfaces, lignes, points)};
% Surfaces
\fill[poporangeL] (0.4,3.9) rectangle (1.1,4.4);
\node[anchor=west, font=\small] at (1.3,4.15) {Zone d'élevage extensif (bovins, zébus)};
\fill[popgreenL] (0.4,3.2) rectangle (1.1,3.7);
\node[anchor=west, font=\small] at (1.3,3.45) {Zone d'élevage périurbain (volailles, porcs)};
% Lignes
\draw[->, line width=2pt, poporange] (0.4,2.75) -- (1.1,2.75);
\node[anchor=west, font=\small] at (1.3,2.75) {Axe de transhumance (saison sèche)};
% Points
\node[font=\small, popblue] at (0.75,2.1) {\textbf{\textbullet}};
\node[anchor=west, font=\small] at (1.3,2.1) {Port de pêche artisanale (Kribi, Limbe, Douala)};
\fill[popblue!50] (0.45,1.35) rectangle (1.05,1.65);
\node[anchor=west, font=\small] at (1.3,1.5) {Étang piscicole / station de pisciculture};
\draw[popdark, thick] (0.4,0.65) rectangle (1.1,1.0);
\node[font=\scriptsize] at (0.75,0.82) {R};
\node[anchor=west, font=\small] at (1.3,0.82) {Ranch / bergerie modernisée};
\end{tikzpicture}
\end{center}
\legende{Modèle de légende hiérarchisée pour le TP 1 : les surfaces d'abord, puis les lignes (flux), puis les points. Source : élaboration à partir des données du MINEPIA.}
\end{popfigure}

\begin{attention}[Ne surcharge pas ta carte]
L'erreur la plus fréquente est de vouloir \textbf{tout} représenter : dix couleurs, des symboles illisibles, des noms de villes partout. Une bonne carte est \cle{sélective} : elle ne montre que ce qui sert la démonstration demandée par le sujet. Quatre à six rubriques de légende suffisent. Si le correcteur ne comprend pas ta carte en dix secondes, elle est ratée, même si elle est « complète ».
\end{attention}

\courssub{TP 2 : construire un croquis d'un site de pêche modernisé}

\textbf{Intuition.} Une photographie du débarcadère de Kribi montre des dizaines de détails : pirogues motorisées, moteurs hors-bord, glacières, cases de fumage, filets étendus, acheteuses... Le croquis n'est pas un dessin artistique : c'est une \textbf{sélection organisée} de ce qui est géographiquement significatif.

\begin{methode}[De la photographie au croquis]
\begin{enumerate}
\item \textbf{Observer la photographie} et y repérer les objets qui traduisent la modernisation : moteurs hors-bord, grandes pirogues, filets en nylon, chambre froide ou glacières, hangars de conservation.
\item \textbf{Découper l'espace en plans} : l'avant-plan (la plage, les pirogues), le plan intermédiaire (les installations : hangars, fumoirs), l'arrière-plan (la mer, l'horizon).
\item \textbf{Schématiser} : chaque objet devient une forme simple (rectangle, triangle, ligne) placée au bon endroit. On ne dessine pas, on \emph{symbolise}.
\item \textbf{Légender par des renvois numérotés} (1, 2, 3...) renvoyant à une légende placée sous le croquis, ou par des flèches courtes.
\item \textbf{Donner un titre en majuscules} indiquant le lieu et le thème.
\end{enumerate}
\end{methode}

\begin{popfigure}
\begin{center}
\begin{tikzpicture}[scale=1]
% Mer
\fill[popblue!25] (0,3.2) rectangle (10,5);
\node[font=\itshape\small, popblue] at (8.6,4.6) {Océan Atlantique};
% Vagues
\draw[popblue, decorate, decoration={snake, amplitude=0.6mm}] (0.3,3.6) -- (3,3.6);
\draw[popblue, decorate, decoration={snake, amplitude=0.6mm}] (4,4.1) -- (6.5,4.1);
% Plage
\fill[popgoldL] (0,0) rectangle (10,3.2);
\node[font=\itshape\small, popdark] at (1.2,0.4) {Plage de Kribi};
% Pirogues motorisées
\draw[fill=poporangeL, popdark] (1.5,3.1) .. controls (2,2.7) and (3.5,2.7) .. (4,3.1) -- cycle;
\draw[fill=popdark] (4,2.95) rectangle (4.3,3.25);
\node[font=\scriptsize] at (2.7,2.95) {1};
\draw[fill=poporangeL, popdark] (5.5,3.3) .. controls (6,3.0) and (7.5,3.0) .. (8,3.3) -- cycle;
\draw[fill=popdark] (8,3.15) rectangle (8.3,3.45);
\node[font=\scriptsize] at (6.7,3.15) {1};
% Hangar de conservation
\draw[fill=poppurpleL, popdark] (0.5,1.2) rectangle (2.3,2.3);
\draw[popdark, fill=poppurple!40] (0.4,2.3) -- (1.4,2.9) -- (2.4,2.3) -- cycle;
\node[font=\scriptsize] at (1.4,1.7) {2};
% Fumoir
\draw[fill=popgreenL, popdark] (3,1.1) rectangle (4.4,2.0);
\draw[popdark] (3.2,2.0) -- (3.2,2.6); \draw[popdark] (4.2,2.0) -- (4.2,2.6);
\draw[popdark] (3.1,2.6) -- (4.3,2.6);
\node[font=\scriptsize] at (3.7,1.5) {3};
% Filets étendus
\draw[popdark, thick] (5,1.0) -- (5,2.4); \draw[popdark, thick] (6.8,1.0) -- (6.8,2.4);
\draw[popline, thick] (5,2.3) .. controls (5.9,1.8) .. (6.8,2.3);
\draw[popline] (5,2.0) .. controls (5.9,1.55) .. (6.8,2.0);
\draw[popline] (5,1.7) .. controls (5.9,1.3) .. (6.8,1.7);
\node[font=\scriptsize] at (5.9,2.55) {4};
% Glacières / caisses
\draw[fill=popblueL, popdark] (7.4,1.0) rectangle (8.2,1.5);
\draw[fill=popblueL, popdark] (8.3,1.0) rectangle (9.1,1.5);
\draw[fill=popblue!40, popdark] (7.4,1.5) rectangle (8.2,1.9);
\node[font=\scriptsize] at (8.25,0.75) {5};
% Nord
\draw[->, thick, popdark] (9.5,4.0) -- (9.5,4.8);
\node[font=\small\bfseries] at (9.5,4.95) {N};
\end{tikzpicture}
\end{center}
\legende{Croquis corrigé : le site de pêche artisanale modernisé de Kribi. 1 : grandes pirogues à moteur hors-bord ; 2 : hangar de conservation (chambre froide) ; 3 : fumoir amélioré ; 4 : filets en nylon étendus au séchage ; 5 : glacières et caisses isothermes pour le transport vers Yaoundé.}
\end{popfigure}

\courssub{TP 3 : construire un diagramme à partir d'un tableau statistique}

\textbf{Intuition.} Un tableau de chiffres « parle » peu. Transformé en diagramme, il révèle immédiatement les hiérarchies : quelle région possède le plus gros cheptel ? quelle zone fournit le plus de poisson ? Le diagramme est un \textbf{outil de démonstration}, pas une décoration.

Voici un tableau de données (ordres de grandeur, à titre pédagogique) :

\begin{center}
\begin{tabularx}{\linewidth}{|l|X|X|}
\hline
\rowcolor{popblueL} \textbf{Région} & \textbf{Cheptel bovin (milliers de têtes, ordre de grandeur)} & \textbf{Part estimée (\%)} \\
\hline
Adamaoua & \num{2500} & 41 \\
\hline
Extrême-Nord & \num{1500} & 25 \\
\hline
Nord & \num{900} & 15 \\
\hline
Nord-Ouest & \num{500} & 8 \\
\hline
Autres régions & \num{700} & 11 \\
\hline
\textbf{Total} & \textbf{\num{6100}} & \textbf{100} \\
\hline
\end{tabularx}
\end{center}

\begin{methode}[Construire un diagramme en bâtons]
\begin{enumerate}
\item \textbf{Tracer deux axes perpendiculaires} : en abscisse, les catégories (les régions) ; en ordonnée, la grandeur (effectif du cheptel), avec une \textbf{échelle graduée régulièrement} (par exemple 1 cm = 500 milliers de têtes).
\item \textbf{Titrer chaque axe} avec la grandeur et l'unité : « Cheptel bovin (milliers de têtes) ».
\item \textbf{Élever des bâtons de hauteur proportionnelle} aux valeurs, de largeur égale et régulièrement espacés.
\item \textbf{Écrire le titre du diagramme} (quoi, où, quand) et la \textbf{source des données}.
\item \textbf{Vérifier la proportionnalité} : un bâton deux fois plus haut doit correspondre à une valeur deux fois plus grande. Ne jamais couper l'axe des ordonnées sans le signaler.
\end{enumerate}
\end{methode}

\begin{popfigure}
\begin{center}
\begin{tikzpicture}
\begin{axis}[
    width=0.85\linewidth, height=6.5cm,
    ybar, bar width=18pt,
    ymin=0, ymax=2800,
    ylabel={Cheptel bovin (milliers de têtes)},
    symbolic x coords={Adamaoua, Extrême-Nord, Nord, Nord-Ouest, Autres},
    xtick=data,
    x tick label style={rotate=20, anchor=east, font=\small},
    ytick={0,500,1000,1500,2000,2500},
    axis lines=left,
    ymajorgrids=true, grid style={popline!40},
    every axis plot/.append style={fill=poporange, draw=popdark},
    nodes near coords, nodes near coords style={font=\scriptsize, popdark},
]
\addplot[popcurve] coordinates {(Adamaoua,2500) (Extrême-Nord,1500) (Nord,900) (Nord-Ouest,500) (Autres,700)};
\end{axis}
\end{tikzpicture}
\end{center}
\legende{Diagramme en bâtons : répartition estimée du cheptel bovin camerounais par région (ordres de grandeur, environ 6,1 millions de têtes au total). Source : élaboration à partir des données du MINEPIA. La lecture est immédiate : l'Adamaoua et l'Extrême-Nord concentrent à eux deux environ les deux tiers du cheptel national.}
\end{popfigure}

Pour un \textbf{diagramme circulaire}, le principe est identique, mais on convertit chaque valeur en \textbf{angle} : l'effectif total correspond à $360^{\circ}$. Ainsi, pour l'Adamaoua :
\[ \text{angle} = \frac{\num{2500}}{\num{6100}} \times 360 \approx 147^{\circ}. \]
Chaque secteur est ensuite colorié et renvoyé à une légende.

\begin{attention}[Un diagramme sans interprétation ne vaut rien]
Recopier les données dans un beau graphique ne suffit pas : au Baccalauréat, le diagramme doit \textbf{servir une démonstration géographique}. Après l'avoir construit, rédige deux ou trois phrases : quel est le fait dominant (la concentration du cheptel au nord) ? Quelle explication (climat soudanien favorable au pâturage, présence des ranchs et des programmes d'appui) ? Quelle conséquence (approvisionnement des villes du sud, flux de transhumance) ? C'est cette analyse qui est notée.
\end{attention}

\courssub{TP 4 : interpréter les documents produits}

Les trois productions (carte, croquis, diagramme) doivent maintenant être \textbf{confrontées} pour dégager une synthèse. La démarche d'interprétation suit trois questions :

\begin{enumerate}
\item \textbf{Que voit-on ?} (constats) : le cheptel et les axes de transhumance se concentrent dans les régions septentrionales ; la pêche maritime artisanale s'organise autour de quelques ports (Kribi, Limbe, Douala) ; la pisciculture se développe près des villes de l'Ouest et du Centre.
\item \textbf{Pourquoi ?} (explications) : au nord, le climat soudano-sahélien offre de vastes pâturages mais impose la mobilité (transhumance) ; sur la côte, la proximité du marché urbain et l'équipement (moteurs, glacières, chambres froides) expliquent la modernisation ; autour de Yaoundé et de Foumban, la demande urbaine en poisson frais stimule les étangs.
\item \textbf{Quels contrastes ?} (synthèse) : on oppose un \textbf{nord pastoral} extensif et mobile à un \textbf{sud} où l'élevage est périurbain et intensif (poulaillers, bergeries améliorées) ; une pêche maritime en voie de motorisation à une pêche continentale (fleuves, lacs) encore peu équipée ; des régions bien encadrées par les programmes de l'État à des zones encore marginalisées.
\end{enumerate}

\begin{exoresolu}[Interpréter un diagramme du cheptel]
\textbf{Énoncé.} À partir du diagramme en bâtons construit au TP 3, rédige un commentaire de cinq à huit lignes dégageant les contrastes régionaux de l'élevage bovin au Cameroun.
\tcblower
\textbf{Solution détaillée.} Le diagramme montre une répartition très inégale du cheptel bovin camerounais. L'Adamaoua domine nettement avec environ 2,5 millions de têtes, soit près de 41 \% du total ; avec l'Extrême-Nord (environ 1,5 million) et le Nord (environ \num{900000}), les trois régions septentrionales concentrent plus de 80 \% du cheptel national. Ce contraste s'explique d'abord par le milieu : les savanes soudaniennes du nord offrent de vastes pâturages, alors que le sud forestier et humide est défavorable aux bovins (maladie du sommeil, manque d'étendues herbeuses). Il s'explique aussi par l'histoire et l'action de l'État : tradition pastorale des éleveurs peuls, création de ranchs et de bergeries, appui des programmes de développement. En conséquence, le sud du pays, plus peuplé et plus urbanisé, dépend du nord pour sa viande, ce qui alimente les flux de transhumance et le commerce de bétail vers Douala et Yaoundé. L'élevage bovin illustre ainsi à la fois la complémentarité des régions et le déséquilibre spatial de la production animale au Cameroun.
\end{exoresolu}

\courssub{Restitution : présentation orale et correction collective}

Chaque groupe présente ses productions au tableau en trois minutes : la carte légendée, le croquis et le diagramme sont affichés côte à côte, et un élève commente les contrastes dégagés. La classe corrige ensuite collectivement selon un barème explicite :

\begin{center}
\begin{tabularx}{\linewidth}{|l|X|c|}
\hline
\rowcolor{popblueL} \textbf{Critère} & \textbf{Ce qui est attendu} & \textbf{Barème} \\
\hline
Exactitude & Informations placées au bon endroit, données fidèles au tableau & /5 \\
\hline
Légende & Classée, hiérarchisée, symboles simples et cohérents & /5 \\
\hline
Propreté & Tracés soignés, titre en majuscules, orientation, écriture lisible & /4 \\
\hline
Interprétation & Contrastes régionaux dégagés, explications proposées & /6 \\
\hline
\end{tabularx}
\end{center}

\begin{aretenir}[L'essentiel des TP]
\begin{itemize}
\item Une carte se construit en cinq étapes : lire, tracer le fond, choisir les symboles, placer, légender. Elle doit être \cle{sélective} et lisible.
\item Un croquis n'est pas un dessin : c'est une sélection schématisée de ce qui est géographiquement significatif, avec renvois numérotés.
\item Un diagramme exige des axes gradués et titrés, des bâtons proportionnels, un titre et une source ; il doit être \textbf{interprété}.
\item La synthèse finale dégage les \cle{contrastes régionaux} : nord pastoral extensif, sud périurbain intensif, littoral de pêche en modernisation.
\end{itemize}
\end{aretenir}

\begin{exercice}[Pour t'entraîner]
À partir du tableau suivant (ordres de grandeur) : pêche maritime artisanale 120 milliers de tonnes, pêche continentale 80 milliers de tonnes, aquaculture 5 milliers de tonnes — construis un diagramme circulaire (calcule les angles) et rédige trois phrases d'interprétation montrant le poids encore faible mais prometteur de la pisciculture.
\end{exercice}

\begin{corrige}[Exercice 1]
Total : $120 + 80 + 5 = 205$ milliers de tonnes. Angles : pêche maritime $\frac{120}{205}\times 360 \approx 211^{\circ}$ ; pêche continentale $\frac{80}{205}\times 360 \approx 140^{\circ}$ ; aquaculture $\frac{5}{205}\times 360 \approx 9^{\circ}$. Interprétation : la pêche maritime artisanale fournit près de 60 \% des captures, devant la pêche continentale (environ 39 \%). L'aquaculture ne représente encore qu'environ 2,5 \% de la production, mais son essor récent (étangs de l'Ouest et du Centre, appuis de l'État) en fait une piste majeure pour réduire les importations de poisson.
\end{corrige}

\newpage

\coursec{Activité d'intégration}

\coursec{Leçon 5 : Les transformations de l'élevage et de la pêche traditionnels}

\courssub{Objectifs de l'activité}

À l'issue de cette activité d'intégration, tu seras capable de :

\begin{itemize}
\item \cle{observer}, \cle{localiser} et \cle{décrire} les grandes zones d'élevage, de pêche et de pisciculture du Cameroun ;
\item \cle{légender} une carte muette et construire un \cle{croquis} simple des axes de transhumance ;
\item analyser des \cle{photographies} selon la méthode en quatre étapes (description, localisation, interprétation, conclusion) ;
\item inventorier et interpréter un \cle{tableau statistique}, puis construire un \cle{diagramme en bâtons} et un \cle{diagramme circulaire} ;
\item rédiger une \cle{synthèse argumentée} évaluant le rôle de l'État dans la modernisation des activités agro-pastorales et piscicoles et leur contribution à l'\cle{intégration nationale}.
\end{itemize}

\courssub{Situation}

\textbf{Du troupeau au marché : moderniser l'élevage et la pêche pour nourrir le Cameroun.}

En mars 2024, le lycée de Ngaoundéré organise une semaine du « consommons local ». Le proviseur confie à ta classe de Terminale la préparation d'une exposition intitulée : \emph{« Comment les transformations de l'élevage et de la pêche, soutenues par l'État, contribuent-elles à l'intégration nationale et à la valorisation des produits locaux ? »}. Tu disposes, avec ton groupe, du dossier documentaire suivant.

\textbf{Document 1 — Carte muette du Cameroun.} Une carte muette avec le contour du pays, les dix chefs-lieux de région, la ligne de chemin de fer Transcamerounais (Douala--Yaoundé--Ngaoundéré) et l'axe routier Douala--Bafoussam--Bamenda. Tu dois y placer : les zones d'élevage extensif (Adamaoua, Nord, Extrême-Nord), les axes de \cle{transhumance} vers les pâturages de saison sèche (vers le Mbéré, le Mbam, la plaine du Logone), les zones de pêche maritime artisanale (Douala, Limbe, Kribi, Campo), les zones de pêche continentale (fleuve Logone, lac Tchad, barrage de Lagdo, barrage de Mbakaou), et les principaux bassins piscicoles (Ouest, Centre, Littoral).

\textbf{Document 2 — Deux photographies.}
\begin{itemize}
\item \emph{Photo A} : une \cle{bergerie} moderne de l'Adamaoua (enclos en dur, abreuvoirs, cases de stockage d'aliments, zébus de race améliorée par \cle{croisement} Goudali $\times$ Zébu Peul).
\item \emph{Photo B} : le débarcadère de pêche artisanale de Kribi : grandes \cle{pirogues motorisées} (moteurs hors-bord de 25 à 40 CV), filets en nylon, glacières et chambre froide de \cle{conservation}, mareyeuses chargeant des caisses isothermes.
\end{itemize}

\textbf{Document 3 — Tableau statistique (ordres de grandeur, sources : MINEPIA, INS ; valeurs arrondies).}

\begin{center}
\begin{tabularx}{\linewidth}{|l|X|X|}\hline
\rowcolor{popblueL} \textbf{Rubrique} & \textbf{Valeur (ordre de grandeur)} & \textbf{Observation} \\ \hline
Cheptel bovin national & 7 à 8 millions de têtes & concentré dans les trois régions septentrionales \\ \hline
Petits ruminants (ovins, caprins) & 10 à 12 millions de têtes & élevage villageois dominant \\ \hline
Volailles & 60 à 70 millions de sujets & forte croissance des poulaillers modernes \\ \hline
Captures de pêche (mer + eaux continentales) & environ \num{250000} à \num{300000} tonnes/an & pêche artisanale : plus de 80 \% des captures \\ \hline
Production piscicole & environ \num{5000} à \num{10000} tonnes/an & en forte progression mais encore faible \\ \hline
Importations de poisson & environ \num{200000} tonnes/an & coût : plus de 100 milliards de FCFA/an \\ \hline
\end{tabularx}
\end{center}

\textbf{Document 4 — Extrait de présentation du PDMA.} Le \cle{Programme de Développement des Marchés Agricoles} (PDMA), cofinancé par l'État et ses partenaires, appuie les producteurs par la formation, le financement d'équipements (bergeries, poulaillers, étangs, pirogues motorisées, chaînes du froid) et l'organisation de la commercialisation. Il s'inscrit dans la stratégie de promotion des produits locaux et de réduction des importations.

\courssub{Consignes}

Travail en groupes de quatre à six élèves. Chaque groupe rend : la carte légendée, l'analyse des deux photographies, les deux diagrammes, puis une synthèse écrite d'une page. La restitution orale (5 minutes par groupe) est suivie d'un débat.

\begin{enumerate}
\item \textbf{Localiser et cartographier.} Sur la carte muette, place et légende : (a) les trois grandes régions d'élevage extensif ; (b) deux axes de transhumance (flèches) ; (c) deux zones de pêche maritime et deux zones de pêche continentale ; (d) deux bassins piscicoles. Rédige une légende numérotée et ordonnée (du plan de fond aux éléments ponctuels).
\item \textbf{Observer et décrire.} Analyse les photographies A et B selon la méthode en quatre étapes. Pour chacune, dégage deux indices de \cle{modernisation} (équipements, bâtiments, organisation du travail).
\item \textbf{Traiter les données.} À partir du document 3 :
\begin{enumerate}
\item calcule la part de la pêche artisanale dans les captures totales (en prenant \num{275000} tonnes de captures) ;
\item calcule le rapport entre les importations de poisson et la production piscicole (en prenant \num{7500} tonnes de production piscicole) ;
\item construis un \cle{diagramme en bâtons} comparant captures, production piscicole et importations ;
\item construis un \cle{diagramme circulaire} de la couverture de la consommation nationale de poisson (production locale : pêche + pisciculture ; importations).
\end{enumerate}
\item \textbf{Interpréter.} Explique en cinq à huit lignes pourquoi, malgré un cheptel important et des captures significatives, le Cameroun importe massivement du poisson et reste dépendant pour sa consommation de protéines animales.
\item \textbf{Évaluer le rôle de l'État.} À partir du document 4 et de ton cours, cite trois types d'actions de l'État (encadrement, financement, formation, régulation de la commercialisation) et montre comment chacune transforme l'élevage ou la pêche.
\item \textbf{Synthétiser.} Rédige une page répondant à la question : \emph{« Comment les transformations de l'élevage et de la pêche, soutenues par l'État, contribuent-elles à l'intégration nationale et à la valorisation des produits locaux ? »}
\item \textbf{Débattre.} Faut-il limiter les importations de poisson pour encourager la production locale ? Prépare deux arguments « pour » et deux arguments « contre ».
\end{enumerate}

\courssub{Ressources à mobiliser}

\begin{aretenir}[Boîte à outils de l'activité]
\begin{itemize}
\item \textbf{Vocabulaire :} \cle{transhumance} (déplacement saisonnier des troupeaux vers les pâturages et points d'eau), \cle{bergerie} (bâtiment d'élevage moderne pour petits ruminants ou bovins), \cle{ranch} (exploitation d'élevage extensive clôturée et gérée de façon rationnelle), \cle{étang} (plan d'eau aménagé pour la pisciculture), \cle{pisciculture} (élevage de poissons : tilapias, clarias, carpes).
\item \textbf{Transformations de l'élevage :} sélection des espèces, croisements, amélioration de l'alimentation et des soins vétérinaires, construction de bergeries et de poulaillers.
\item \textbf{Modernisation de la pêche :} grandes pirogues à moteur hors-bord, filets adaptés (nylon), conservation (glace, chambre froide, fumage amélioré).
\item \textbf{Calculs :} part en pourcentage : $\text{part} = \dfrac{\text{valeur}}{\text{total}} \times 100$ ; angle d'un secteur circulaire : $\text{angle} = \dfrac{\text{valeur}}{\text{total}} \times 360$ (en degrés).
\item \textbf{Méthode d'analyse de photographie :} 1) décrire objectivement ; 2) localiser ; 3) interpréter (relier aux notions du cours) ; 4) conclure sur le phénomène géographique.
\item \textbf{Rôle de l'État :} programmes et projets (PDMA), formation des acteurs, financement, régulation de la commercialisation.
\end{itemize}
\end{aretenir}

\courssub{Démarche et corrigé commenté}

\begin{exoresolu}[Corrigé guidé de l'activité]
\textbf{Énoncé.} Réaliser la carte légendée, analyser les photographies, traiter les données statistiques, puis rédiger la synthèse répondant à la question : comment les transformations de l'élevage et de la pêche, soutenues par l'État, contribuent-elles à l'intégration nationale et à la valorisation des produits locaux ?

\tcblower

\textbf{Étape 1 — Carte légendée.} La légende s'ordonne du plan de fond vers les éléments ponctuels (surfaces, puis linéaires, puis ponctuels) :
\begin{enumerate}
\item Zones d'élevage extensif (Adamaoua, Nord, Extrême-Nord) — aplat de couleur claire (ex. jaune pâture) ;
\item Axes de transhumance — flèches pleines orientées vers le Sud (vers la vallée du Mbéré, le Mbam, la plaine du Logone) ;
\item Zones de pêche maritime artisanale (Douala, Limbe, Kribi, Campo) — hachures bleues le long du littoral ;
\item Zones de pêche continentale (fleuve Logone, lac Tchad, barrage de Lagdo, barrage de Mbakaou) — symbole « poisson » bleu sur les cours d'eau et plans d'eau ;
\item Bassins piscicoles (Ouest, Centre, Littoral) — points de couleur verte (étangs) ;
\item Villes et axes de commercialisation (chemin de fer Transcamerounais, route nationale N1/N6) — points nommés et traits noirs/rouges.
\end{enumerate}
\emph{Justification :} une légende hiérarchisée respecte les règles de la cartographie scolaire : d'abord les surfaces (fond), puis les linéaires (flux, axes), enfin les ponctuels (villes, étangs). Le titre, l'orientation (flèche Nord) et l'échelle indicative sont obligatoires.

\begin{popfigure}
\begin{center}
\includegraphics[width=\linewidth,height=9.5cm,keepaspectratio]{N01/cmr_map.jpg}
\end{center}
\legende{Fond de carte : carte générale du Cameroun (relief ombré, régions, villes, routes, chemin de fer, fleuves, lac Tchad et retenues de Lagdo et de Mbakaou). Le croquis attendu situe : 1. les zones d'élevage extensif du Grand Nord et de l'Adamaoua (Ngaoundéré) ; 2. les axes de transhumance vers le sud ; 3. la pêche maritime sur le golfe de Guinée (Douala, Kribi, Campo) ; 4. la pêche continentale (Logone, lac Tchad, Bénoué, retenues de Lagdo et de Mbakaou) ; 5. les villes et axes de commercialisation (chemin de fer Ngaoundéré--Yaoundé--Douala, routes). Les symboles et la légende hiérarchisée sont à ajouter par l'élève. \\ Source : Wikimedia Commons, CIA, domaine public.}
\end{popfigure}

\textbf{Étape 2 — Analyse des photographies.}
\emph{Photo A (bergerie de l'Adamaoua) :} \textbf{Description :} bâtiment en dur (parpaings/tôles), enclos délimités, abreuvoirs fixes, cases de stockage d'aliments (son, tourteaux, foin), présence de zébus à robe claire (type Goudali croisé). \textbf{Localisation :} plateau de l'Adamaoua (région de l'Adamaoua), principale zone d'élevage bovin du pays. \textbf{Interprétation :} la construction de bergeries traduit le passage d'un élevage extensif pur (parcours libre) à un élevage \emph{semi-intensif} (parcours + alimentation complémentaire, abri). Les croisements (Goudali $\times$ Zébu Peul) visent à cumuler la rusticité locale et la productivité bouchère/laitière. \textbf{Conclusion :} l'élevage se modernise pour augmenter la productivité par tête et sécuriser la production (moins de pertes par vol, maladies, intempéries).

\emph{Photo B (débarcadère de Kribi) :} \textbf{Description :} grandes pirogues en bois motorisées (hors-bord 25--40 CV), filets en nylon synthétique empilés, glacières plastiques, chambre froide en conteneur, mareyeuses (femmes commerçantes) chargeant des caisses isothermes sur véhicules. \textbf{Localisation :} Kribi, département de l'Océan, région du Sud, façade atlantique. \textbf{Interprétation :} la motorisation élargit le rayon d'action (pêche au large, zones plus poissonneuses) et réduit la pénibilité. Les filets en nylon (mailles calibrées) sont plus sélectifs et résistants que les fibres naturelles. La chaîne du froid (glace, chambre froide, caisses isothermes) réduit drastiquement les pertes post-capture (pourriture) et permet l'acheminement vers les grands centres urbains (Douala, Yaoundé). \textbf{Conclusion :} la pêche artisanale s'équipe, se professionnalise et s'intègre aux circuits commerciaux nationaux.

\textbf{Étape 3 — Traitement statistique.}
\begin{align*}
\text{Part de la pêche artisanale} &= 80\ \% \times \num{275000} = \num{220000} \text{ tonnes} \\
\text{Rapport importations / pisciculture} &= \frac{\num{200000}}{\num{7500}} \approx 26,7 \approx \textbf{27}
\end{align*}
Le Cameroun importe donc environ \textbf{27 fois plus} de poisson qu'il n'en produit en pisciculture : un écart majeur.

Pour le diagramme circulaire, l'approvisionnement total en poisson (hypothèse : pas d'exportations nettes, stocks constants) est estimé à :
\[ \text{Total} = \num{275000} + \num{7500} + \num{200000} = \num{482500} \text{ tonnes} \]
\begin{align*}
\text{Pêche locale (captures)} &= \frac{\num{275000}}{\num{482500}} \times 100 \approx 57,0\ \% \quad (\text{angle} \approx 205^\circ) \\
\text{Pisciculture} &= \frac{\num{7500}}{\num{482500}} \times 100 \approx 1,6\ \% \quad (\text{angle} \approx 6^\circ) \\
\text{Importations} &= \frac{\num{200000}}{\num{482500}} \times 100 \approx 41,4\ \% \quad (\text{angle} \approx 149^\circ)
\end{align*}

\begin{popfigure}
\centering
\begin{tikzpicture}
\begin{axis}[
    ybar,
    width=\linewidth,
    height=7cm,
    bar width=20pt,
    symbolic x coords={Captures (pêche), Importations, Pisciculture},
    xtick=data,
    x tick label style={rotate=30, anchor=east, font=\small},
    ylabel={Quantité (milliers de tonnes)},
    ymin=0, ymax=300,
    ymajorgrids=true,
    grid style=dashed,
    nodes near coords,
    nodes near coords align={vertical},
    every node near coord/.append style={font=\tiny, rotate=90, anchor=west, yshift=2pt},
    popbar1/.style={fill=popblue},
    popbar2/.style={fill=poporange},
    popbar3/.style={fill=popgreen},
]
\addplot[popbar1] coordinates {(Captures (pêche), 275)};
\addplot[popbar2] coordinates {(Importations, 200)};
\addplot[popbar3] coordinates {(Pisciculture, 7.5)};
\legend{Pêche, Importations, Pisciculture}
\end{axis}
\end{tikzpicture}
\legende{Diagramme en bâtons : comparaison des captures de pêche, des importations et de la production piscicole (en milliers de tonnes). La pisciculture est à peine visible.}
\end{popfigure}

\begin{popfigure}
\centering
\begin{tikzpicture}
\begin{scope}[scale=1.0]
\fill[popblue] (0,0) -- (0:2.3) arc (0:205.2:2.3) -- cycle;
\fill[poporange] (0,0) -- (205.2:2.3) arc (205.2:354.24:2.3) -- cycle;
\fill[popgreen] (0,0) -- (354.24:2.3) arc (354.24:360:2.3) -- cycle;
\draw[white, thick] (0,0) circle (2.3);
\node[white, font=\bfseries] at (100:1.3) {57\,\%};
\node[white, font=\bfseries] at (280:1.3) {41,4\,\%};
\fill[popblue] (3.2,1.0) rectangle ++(0.3,0.3); \node[anchor=west] at (3.6,1.15) {\small Pêche locale (captures)};
\fill[poporange] (3.2,0.4) rectangle ++(0.3,0.3); \node[anchor=west] at (3.6,0.55) {\small Importations};
\fill[popgreen] (3.2,-0.2) rectangle ++(0.3,0.3); \node[anchor=west] at (3.6,-0.05) {\small Pisciculture (1,6\,\%)};
\end{scope}
\end{tikzpicture}
\legende{Diagramme circulaire : structure de l'approvisionnement en poisson au Cameroun (estimation). Les importations représentent près de 42 \% de l'offre.}
\end{popfigure}

\textbf{Étape 4 — Interprétation.} Le Cameroun possède un cheptel important (7--8 millions de bovins) et des captures halieutiques significatives (\num{275000} t), mais la production locale ne suit pas la croissance démographique (environ 2,6 \%/an) et l'urbanisation rapide qui boostent la demande en protéines animales. L'élevage extensif a une faible productivité par tête (pas d'alimentation complémentaire systématique, races locales peu productives en lait/viande). La pêche artisanale, malgré sa modernisation, subit la pression sur les stocks côtiers (surpêche, destruction d'habitats) et la concurrence de la pêche industrielle (souvent illicite). La pisciculture reste embryonnaire (environ 1,6 \% de l'approvisionnement) faute d'aliments de qualité accessibles, de financement et de maîtrise technique généralisée. D'où des importations massives de poisson (environ \num{200000} t/an), coûteuses en devises (plus de 100 milliards de FCFA/an), creusant le déficit de la balance commerciale agroalimentaire.

\textbf{Étape 5 — Rôle de l'État.} Trois actions majeures transforment les filières :
\begin{enumerate}
\item \textbf{Encadrement et financement (PDMA, PEA-Jeunes, PADFA) :} subventions pour la construction de bergeries, poulaillers, étangs ; acquisition de pirogues motorisées, de chaînes du froid, d'aliments composés. \emph{Effet :} baisse du coût d'investissement pour les producteurs, passage à l'intensif.
\item \textbf{Formation des acteurs :} centres de formation (ex. CFP de Mbalmayo, centres régionaux MINEPIA), appui-conseil par les techniciens de terrain, vulgarisation des techniques d'insémination artificielle, de fabrication d'aliments flottants pour poissons, de gestion d'étang. \emph{Effet :} appropriation des innovations, meilleure gestion économique.
\item \textbf{Régulation de la commercialisation :} organisation des marchés ruraux et urbains, promotion de la marque « Produit du Cameroun », lutte contre la fraude sur les importations (poissons avariés), appui aux groupements (coopératives, GIC) pour la vente collective. \emph{Effet :} meilleure rémunération des producteurs, flux organisés entre zones de production (Nord, Littoral) et zones de consommation (villes).
\end{enumerate}

\textbf{Étape 6 — Synthèse (trame attendue).}
\emph{Introduction :} L'élevage et la pêche traditionnels constituent le socle de l'approvisionnement en protéines animales au Cameroun. Pourtant, face à la demande croissante, ils engagent depuis deux décennies une mutation profonde, soutenue par les pouvoirs publics.
\emph{Développement :} Cette mutation se traduit par \textbf{la modernisation des équipements et des pratiques} : en élevage, sélection, croisements (Goudali $\times$ Peul), bergeries, poulaillers, alimentation améliorée ; en pêche, motorisation des pirogues, filets en nylon, conservation par le froid ; en pisciculture, essor des étangs (tilapia, clarias) dans l'Ouest, le Centre et le Littoral. \textbf{L'État est le moteur principal} via le PDMA (financement, équipements), la formation technique et la régulation des marchés. \textbf{Ces transformations favorisent l'intégration nationale} : elles structurent des flux Nord-Sud (viande, bétail sur pied vers Douala/Yaoundé via le Transcamerounais et la route) et Littoral-Intérieur (poisson frais/congelé vers l'Ouest et le Centre), créent des revenus ruraux, fixent les populations et valorisent le « consommons local ».
\emph{Conclusion :} La modernisation agro-pastorale et piscicole est un levier puissant d'intégration territoriale et de souveraineté alimentaire. Son aboutissement suppose de combler le retard de la pisciculture et de réduire la dépendance aux importations par un investissement public-privé soutenu.

\textbf{Étape 7 — Débat.}
\emph{Arguments POUR une limitation des importations :} (1) Protection des revenus des pêcheurs et pisciculteurs locaux face au dumping (poissons importés souvent moins chers car subventionnés à l'origine ou de basse qualité) ; (2) Économie de devises (100+ milliards FCFA/an) réinvestissables dans la filière locale.
\emph{Arguments CONTRE une limitation brutale :} (1) Risque de pénurie et d'envolée des prix pour les consommateurs urbains pauvres, la production locale (pêche + pisciculture $\approx$ 58 \% de l'offre) ne couvrant pas la demande ; (2) Nécessité de respecter les accords commerciaux (CEMAC, OMC, APE UE) qui interdisent les barrières quantitatives non justifiées sanitairement. \emph{Position nuancée :} mise en place de mesures de sauvegarde temporaires (taxes à l'importation modulées, normes sanitaires strictes) couplées à un plan massif d'investissement dans la pisciculture et la pêche durable pour atteindre l'autosuffisance à horizon 10--15 ans.
\end{exoresolu}

\begin{attention}[Erreurs fréquentes à éviter]
\begin{itemize}
\item Confondre \cle{transhumance} (déplacement saisonnier aller-retour entre pâturages secs et humides) et \cle{nomadisme} (déplacement permanent sans base fixe) ;
\item Oublier d'ordonner la légende (surfaces $\to$ linéaires $\to$ ponctuels) ou d'indiquer le titre, l'orientation (flèche Nord) et l'échelle indicative de la carte ;
\item Additionner pêche et importations sans préciser qu'il s'agit d'une \emph{estimation de l'approvisionnement apparent} (hypothèse : pas d'exportations, stocks constants), pas d'un chiffre officiel de consommation ;
\item Affirmer que la pisciculture peut remplacer \emph{immédiatement} les importations : elle ne représente qu'environ 1,6 \% de l'approvisionnement actuel ;
\item Décrire une photographie sans l'interpréter : l'essentiel est de relier ce qu'on voit (bâtiment, moteur, glace) aux notions du cours (intensification, réduction des pertes, intégration au marché).
\end{itemize}
\end{attention}

\courssub{Bilan de l'activité}

Cette activité t'a permis de mobiliser l'ensemble des savoir-faire de la leçon : localiser et cartographier les zones d'élevage, de pêche et de pisciculture ; analyser des photographies avec méthode ; calculer des pourcentages et construire des diagrammes ; rédiger une synthèse argumentée. Elle a mis en évidence trois idées forces. Premièrement, l'élevage et la pêche traditionnels camerounais sont en pleine mutation : bergeries, poulaillers, croisements, pirogues motorisées, conservation du poisson et essor de la pisciculture en témoignent. Deuxièmement, ces transformations restent insuffisantes face à la demande nationale, comme le montre le poids des importations de poisson (environ 41 \% de l'approvisionnement apparent). Troisièmement, l'État joue un rôle moteur — encadrement, financement, formation, régulation de la commercialisation — pour faire de la modernisation agro-pastorale et piscicole un instrument d'\cle{intégration nationale} et de valorisation des produits locaux. Tu es maintenant prêt pour le dossier 1 et pour l'épreuve de géographie du Baccalauréat.

\newpage

\coursec{Méthodes et exercices résolus}

\coursec{Leçon 5 : Les transformations de l'élevage et de la pêche traditionnels}

\courssub{Fiche méthode 1 : Observer, localiser et décrire un phénomène géographique (photographie, paysage d'élevage ou de pêche)}

\begin{methode}[Décrire une photographie ou un paysage en 5 étapes]
\begin{enumerate}
\item \textbf{Identifier le document} : nature (photographie aérienne ou au sol, carte, tableau), lieu, date, auteur si connu. Annonce le thème : « Cette photographie montre un élevage de zébus dans la région de l'Adamaoua ».
\item \textbf{Localiser} : situe précisément (région, département, ville, milieu naturel). Utilise les points cardinaux et les repères visibles (fleuve, route, relief).
\item \textbf{Décrire du général au particulier} : d'abord l'ensemble (premier plan, arrière-plan), puis les détails significatifs (animaux, bergerie, pirogues à moteur, étangs). Décris ce que tu \emph{vois}, pas ce que tu \emph{sais}.
\item \textbf{Interpréter} : relie chaque élément observé à une notion du cours (transhumance, modernisation, pisciculture). C'est ici que tu mobilises tes connaissances.
\item \textbf{Conclure} : une phrase bilan qui répond au problème posé par la consigne (ex. : « Ce paysage illustre la mutation d'un élevage traditionnel vers un élevage amélioré »).
\end{enumerate}
\textbf{Réflexes} : vocabulaire géographique précis (troupeau, pâturage, plan d'eau, diguette d'étang) ; jamais de « on voit des trucs » ; toujours justifier l'interprétation par un indice visible sur le document.
\end{methode}

\begin{exoresolu}[Description d'une photographie : un ranch dans l'Adamaoua]
\textbf{Énoncé.} Une photographie montre, dans la région de l'Adamaoua, un vaste enclos clôturé où pâturent des zébus ; on distingue une mare artificielle, un hangar de stockage de fourrage et une case de gardien. À l'arrière-plan, des savanes arborées s'étendent à perte de vue. Décrivez cette photographie et montrez qu'elle illustre les transformations de l'élevage traditionnel au Cameroun. (4 points)
\tcblower
\textbf{Solution détaillée.}

\emph{Étape 1 — Identification.} Il s'agit d'une photographie prise au sol, dans la région de l'Adamaoua (Nord-Cameroun), montrant une exploitation d'élevage bovin de type ranch.

\emph{Étape 2 — Localisation.} L'Adamaoua, région de hauts plateaux (environ \SI{1000}{m} d'altitude), est la première zone d'élevage bovin du Cameroun : le climat soudanien, plus frais et moins infesté par la mouche tsé-tsé que le Sud forestier, y est favorable.

\emph{Étape 3 — Description.} Au premier plan, un troupeau de zébus pâture dans un enclos délimité par une clôture. Au centre, une mare artificielle sert de point d'eau permanent. À droite, un hangar abrite du fourrage stocké (foin, résidus de récolte). Une case de gardien témoigne d'une présence humaine permanente. L'arrière-plan est une savane arborée.

\emph{Étape 4 — Interprétation.} Chaque indice visible correspond à une transformation étudiée : la \textbf{clôture} marque la sédentarisation de l'élevage, qui rompt avec la transhumance extensive ; la \textbf{mare} et le \textbf{hangar de fourrage} montrent l'amélioration de l'alimentation et de l'abreuvement ; la \textbf{case de gardien} traduit l'amélioration des soins et de la surveillance. L'ensemble forme un \cle{ranch}, c'est-à-dire une exploitation d'élevage moderne, clôturée et gérée de façon rationnelle.

\emph{Conclusion.} Cette photographie illustre le passage d'un élevage traditionnel extensif et mobile à un élevage sédentarisé, encadré et productif, en Adamaoua, principal bassin d'élevage du Cameroun.
\end{exoresolu}

\courssub{Fiche méthode 2 : Légender une carte}

\begin{methode}[Construire et légender une carte en 6 étapes]
\begin{enumerate}
\item \textbf{Lire la consigne} : quels phénomènes faut-il représenter (zones d'élevage, zones de pêche, axes de transhumance) ?
\item \textbf{Tracer le fond de carte} : contour simplifié du Cameroun, flèche du nord, titre, échelle indicative.
\item \textbf{Choisir les figurés} : une couleur ou un tramage par catégorie de zone ; des flèches pour les flux (transhumance) ; des points nommés pour les villes et ports.
\item \textbf{Placer les éléments} avec exactitude relative : Adamaoua au centre-nord, ports de pêche (Douala, Limbé, Kribi) sur la côte, lac Tchad à l'extrême nord.
\item \textbf{Rédiger la légende} : chaque figuré est expliqué, dans l'ordre logique (d'abord les zones, puis les axes, puis les points). La légende est numérotée ou fléchée.
\item \textbf{Vérifier} : titre, orientation, légende complète, aucun figuré orphelin (utilisé sur la carte mais absent de la légende, ou inversement).
\end{enumerate}
\textbf{Réflexe} : la légende est \emph{analysée} dans le texte qui accompagne la carte : ne te contente pas de dessiner, commente ce que la carte montre.
\end{methode}

\begin{exoresolu}[Légender une carte : les grandes zones d'élevage et de pêche du Cameroun]
\textbf{Énoncé.} Sur le fond de carte du Cameroun fourni, placez et légendez : (a) les deux principales zones d'élevage bovin ; (b) les trois principales zones de pêche (maritime, lacustre, continentale) ; (c) un axe de transhumance. (6 points)
\tcblower
\textbf{Solution détaillée.}

\emph{Raisonnement préalable.} Les savoirs du cours permettent de localiser : l'élevage bovin est concentré dans les savanes du Nord (Adamaoua, Nord, Extrême-Nord) ; la pêche maritime s'exerce sur la côte atlantique (Douala, Limbé, Kribi) ; la pêche lacustre autour du lac Tchad (et du lac de Maga) ; la pêche continentale dans les grands fleuves (Sanaga, Logone, Bénoué). La transhumance conduit les troupeaux de l'Extrême-Nord vers l'Adamaoua et le Nord en saison sèche, à la recherche de pâturages et de points d'eau.

\begin{popfigure}
\begin{center}
\includegraphics[width=\linewidth,height=9.5cm,keepaspectratio]{N01/cmr_map.jpg}
\end{center}
\legende{Fond de carte : carte générale du Cameroun (régions, villes, fleuves, lac Tchad, retenues de Lagdo et de Mbakaou). Le croquis attendu situe : 1. la zone d'élevage de l'Adamaoua (autour de Ngaoundéré) ; 2. la zone d'élevage du Grand Nord (Nord et Extrême-Nord, de Garoua à Maroua) ; 3. la pêche maritime sur la côte (Douala, Kribi) ; 4. la pêche lacustre au lac Tchad, à l'extrême nord ; 5. un axe de transhumance de saison sèche, tracé en flèche de l'Extrême-Nord vers le Nord et l'Adamaoua. \\ Source : Wikimedia Commons, CIA, domaine public.}
\end{popfigure}

\emph{Légende rédigée.} 1 : zone d'élevage bovin de l'Adamaoua (premier bassin d'élevage national) ; 2 : zone d'élevage du Nord et de l'Extrême-Nord ; 3 : zone de pêche maritime artisanale (Douala, Limbé, Kribi) ; 4 : zone de pêche lacustre (lac Tchad, lac de Maga) ; 5 : axe de transhumance vers les pâturages de saison sèche.

\emph{Conclusion.} La carte montre une complémentarité spatiale : l'élevage domine dans les savanes septentrionales, la pêche sur la côte et les plans d'eau, et la transhumance relie les deux ensembles, illustrant la mobilité traditionnelle des troupeaux.
\end{exoresolu}

\courssub{Fiche méthode 3 : Construire un croquis de géographie}

\begin{methode}[Réaliser un croquis (différence avec la carte)]
\begin{enumerate}
\item \textbf{Comprendre la différence} : une \emph{carte} reporte des données localisées ; un \emph{croquis} est un schéma \emph{explicatif} qui met en évidence un fonctionnement (ex. : le fonctionnement d'un ranch, d'un étang piscicole).
\item \textbf{Dégager l'idée directrice} : le croquis doit répondre à une problématique (« comment fonctionne un élevage sédentarisé ? »).
\item \textbf{Tracer les grandes masses} : 3 à 5 zones maximum (pâturages, bâtiments, point d'eau, cultures fourragères).
\item \textbf{Figurer les dynamiques} : flèches (déplacements des animaux, apports d'aliments, vente des produits).
\item \textbf{Légender par numéros ou flèches}, avec des verbes d'action dans la légende (« abreuve », « stocke », « exporte »).
\item \textbf{Titrer} le croquis de façon explicative : « Le fonctionnement d'un ranch sédentarisé dans l'Adamaoua ».
\end{enumerate}
\textbf{Réflexe} : un bon croquis tient sur une demi-page, utilise 4 à 6 figurés, et sa légende peut être lue comme une phrase d'analyse.
\end{methode}

\begin{exoresolu}[Construire un croquis : le fonctionnement d'un étang piscicole]
\textbf{Énoncé.} À l'aide d'un croquis légendé, montrez le fonctionnement d'un étang piscicole aménagé et son rôle dans l'essor de la pisciculture au Cameroun. (5 points)
\tcblower
\textbf{Solution détaillée.}

\emph{Idée directrice.} L'étang piscicole est un plan d'eau artificiel, maîtrisé (alimentation en eau, fertilisation, alimentation des poissons, récolte), qui permet une production régulière de poissons (tilapias, carpes, clarias) indépendante des captures en milieu naturel.

\begin{popfigure}
\begin{tikzpicture}[scale=1]
\draw[thick, popdark, fill=popgreenL] (0,0) rectangle (7,4);
\draw[thick, popblue, fill=popblueL] (2,0.8) rectangle (5.5,3.2);
\draw[->, popblue, very thick] (0.4,2) -- (2,2);
\draw[->, popblue, very thick] (5.5,2) -- (6.6,2);
\fill[popgoldL, draw=popdark] (0.6,3.4) rectangle (1.6,3.9);
\node[font=\scriptsize] at (1.1,3.65) {Aleviniers};
\fill[poporangeL, draw=popdark] (5.9,3.3) rectangle (6.9,3.8);
\node[font=\scriptsize] at (6.4,3.55) {Magasin};
\draw[->, poporange, thick] (1.6,3.4) -- (2.6,3.15);
\draw[->, poporange, thick] (4.8,3.2) -- (6.1,3.35);
\node[font=\scriptsize, popdark] at (3.75,2) {Étang de grossissement};
\node[font=\scriptsize, popdark] at (3.75,1.6) {(tilapias, clarias)};
\end{tikzpicture}
\legende{Croquis de fonctionnement d'un étang piscicole : 1. arrivée d'eau maîtrisée ; 2. mise en charge des alevins (depuis l'alevinier) ; 3. grossissement avec alimentation ; 4. vidange et récolte ; 5. stockage et vente. Échelle indicative.}
\end{popfigure}

\emph{Analyse de la légende.} L'eau entre et sort de façon contrôlée (figurés 1 et 4) : c'est la maîtrise du milieu qui distingue la pisciculture de la pêche de capture. Les alevins produits dans l'alevinier sont transférés dans l'étang de grossissement (2), où ils sont nourris (3) ; la récolte se fait par vidange (4), puis le poisson est stocké et vendu (5), alimentant les marchés locaux (Yaoundé, Douala).

\emph{Conclusion.} Ce croquis montre que la pisciculture transforme le poisson en produit d'élevage : production planifiée, régulière et proche des villes, ce qui réduit la pression sur les pêcheries naturelles et les importations.
\end{exoresolu}

\courssub{Fiche méthode 4 : Construire un diagramme (barres, courbe, circulaire) à partir d'un tableau}

\begin{methode}[Passer du tableau statistique au diagramme]
\begin{enumerate}
\item \textbf{Identifier les variables} : quelle variable en abscisse (temps, catégories), quelle variable en ordonnée (effectif, tonnage) ?
\item \textbf{Choisir le type de diagramme} : barres pour comparer des catégories ; courbe pour une évolution dans le temps ; circulaire pour des parts (en \%).
\item \textbf{Calculer les grandeurs utiles} : pourcentages (part $= \dfrac{\text{valeur}}{\text{total}} \times 100$), angles du diagramme circulaire (angle $= \dfrac{\text{valeur}}{\text{total}} \times 360$), taux de variation $\left(\dfrac{V_f - V_i}{V_i} \times 100\right)$.
\item \textbf{Tracer avec soin} : axes gradués et légendés avec les unités, titre complet (quoi, où, quand).
\item \textbf{Exploiter le diagramme} : dégager la tendance générale, les ruptures, les contrastes ; rédiger 2 ou 3 phrases d'analyse.
\end{enumerate}
\textbf{Formules à retenir} :
\[ \text{Part en \%} = \frac{\text{effectif de la catégorie}}{\text{effectif total}} \times 100 \qquad ; \qquad \text{Taux de variation} = \frac{V_f - V_i}{V_i} \times 100. \]
\end{methode}

\begin{exoresolu}[Construire et analyser un diagramme en barres : la production halieutique]
\textbf{Énoncé.} Le tableau ci-dessous donne des ordres de grandeur de la production de poisson au Cameroun (en milliers de tonnes, données indicatives).

\begin{center}
\begin{tabularx}{\linewidth}{|l|X|X|X|}\hline
\rowcolor{popblueL} \textbf{Origine} & \textbf{Pêche maritime} & \textbf{Pêche continentale} & \textbf{Pisciculture} \\ \hline
Production (milliers de t) & 150 & 80 & 5 \\ \hline
\end{tabularx}
\end{center}

(a) Calculez la part de chaque origine dans la production totale. (b) Représentez ces données par un diagramme en barres. (c) Dégagez deux enseignements. (6 points)
\tcblower
\textbf{Solution détaillée.}

\emph{(a) Calcul des parts.} Le total est : $150 + 80 + 5 = 235$ milliers de tonnes.
\begin{align*}
\text{Pêche maritime} &: \frac{150}{235} \times 100 \approx 63,8\,\% \\
\text{Pêche continentale} &: \frac{80}{235} \times 100 \approx 34,0\,\% \\
\text{Pisciculture} &: \frac{5}{235} \times 100 \approx 2,1\,\%
\end{align*}
Vérification : $63,8 + 34,0 + 2,1 = 99,9\,\% \approx 100\,\%$ (écart dû aux arrondis).

\emph{(b) Diagramme en barres.}

\begin{popfigure}
\begin{tikzpicture}
\begin{axis}[
  ybar, bar width=0.9cm, width=0.85\linewidth, height=5.5cm,
  symbolic x coords={Maritime, Continentale, Pisciculture},
  xtick=data, ymin=0, ymax=170,
  ylabel={Production (milliers de t)},
  nodes near coords, every node near coord/.append style={font=\scriptsize},
  axis lines=left, tick label style={font=\scriptsize}
]
\addplot[fill=popblue, draw=popdark] coordinates {(Maritime,150) (Continentale,80) (Pisciculture,5)};
\end{axis}
\end{tikzpicture}
\legende{Production halieutique du Cameroun par origine (ordres de grandeur indicatifs, en milliers de tonnes).}
\end{popfigure}

\emph{(c) Enseignements.} Premier enseignement : la pêche maritime domine nettement (près des deux tiers de la production), devant la pêche continentale (environ un tiers). Deuxième enseignement : la pisciculture reste marginale (environ 2\,\%) malgré son essor récent, ce qui montre un potentiel de développement considérable pour réduire les importations de poisson.

\emph{Conclusion.} Le Cameroun tire l'essentiel de son poisson des pêcheries de capture ; la pisciculture, encore faible, constitue la voie d'avenir pour accroître l'offre nationale.
\end{exoresolu}

\courssub{Fiche méthode 5 : Inventorier, classer et interpréter des documents}

\begin{methode}[Exploiter un dossier documentaire (méthode de l'inventaire raisonné)]
\begin{enumerate}
\item \textbf{Inventorier} : liste tous les documents (nature, source, date) sans en oublier ; chaque document doit être cité au moins une fois dans la copie.
\item \textbf{Classer} : regroupe les documents selon des critères explicites. Exemple de grille pour ce chapitre : documents sur l'\emph{élevage} / la \emph{pêche} / la \emph{pisciculture} ; ou bien documents sur les \emph{acteurs} / les \emph{techniques} / les \emph{résultats}.
\item \textbf{Croiser} : confronte les documents entre eux (une carte confirme-t-elle un tableau ? une photographie illustre-t-elle un texte ?).
\item \textbf{Interpréter} : dégage les informations qui répondent à la problématique, en citant le document entre parenthèses : « (doc. 2) ».
\item \textbf{Problématiser et conclure} : formule une réponse argumentée à la question posée, en hiérarchisant les facteurs (rôle de l'État, initiatives privées, contraintes naturelles).
\end{enumerate}
\textbf{Réflexes} : annonce ton plan de classement ; cite systématiquement les documents ; ne recopie jamais un document, exploite-le.
\end{methode}

\begin{exoresolu}[Étude de cas documentaire : le rôle de l'État dans le développement agro-pastoral et piscicole]
\textbf{Énoncé.} Dossier de trois documents. \textbf{Doc. 1} : texte présentant le Programme de développement des marchés agricoles (PDMA), qui finance des infrastructures rurales (pistes, marchés, points d'eau). \textbf{Doc. 2} : tableau montrant la hausse du nombre de pisciculteurs formés dans les régions du Centre et de l'Ouest entre 2015 et 2022. \textbf{Doc. 3} : photographie d'un marché à bétail encadré par des services vétérinaires dans le Nord. À l'aide des documents et de vos connaissances, montrez que l'État camerounais joue un rôle moteur dans le développement des activités agro-pastorales et piscicoles. (8 points)
\tcblower
\textbf{Solution détaillée.}

\emph{Inventaire et classement.} Le dossier comporte trois documents complémentaires : un texte sur un programme de financement (doc. 1), un tableau statistique sur la formation (doc. 2), une photographie sur l'encadrement sanitaire et commercial (doc. 3). On peut les classer selon trois types d'intervention de l'État : le \emph{financement} (doc. 1), la \emph{formation} (doc. 2), l'\emph{encadrement et la régulation} (doc. 3).

\emph{Analyse croisée.} Le doc. 1 montre que l'État, à travers des programmes comme le PDMA, finance les infrastructures qui désenclavent les zones de production et rapprochent producteurs et consommateurs : pistes rurales, marchés aménagés, forages et points d'eau pour le bétail. Le doc. 2 atteste l'effort de formation des acteurs : la hausse du nombre de pisciculteurs formés dans le Centre et l'Ouest entre 2015 et 2022 traduit la volonté de diffuser les techniques modernes (alevinage, alimentation, gestion des étangs) ; cela explique l'essor de la pisciculture observé dans le cours. Le doc. 3 illustre la régulation de la commercialisation : la présence de services vétérinaires sur les marchés à bétail garantit la santé des troupeaux (vaccinations, contrôles) et sécurise les transactions.

\emph{Mobilisation des connaissances.} Ces interventions s'inscrivent dans la continuité des actions historiques de l'État (sociétés de développement comme la SODEPA pour l'élevage, le MINEPIA pour l'encadrement de l'élevage, des pêches et des industries animales) et répondent à un objectif d'intégration nationale : valoriser les produits locaux et réduire la dépendance aux importations de viande et de poisson.

\emph{Conclusion.} Par le financement des infrastructures (doc. 1), la formation des acteurs (doc. 2) et la régulation sanitaire et commerciale (doc. 3), l'État camerounais apparaît comme l'acteur central de la modernisation des activités agro-pastorales et piscicoles, même si des limites subsistent (financement insuffisant, enclavement persistant de certaines zones).
\end{exoresolu}

\courssub{Fiche méthode 6 : Calculer et interpréter un taux de variation (production, effectifs)}

\begin{methode}[Taux de variation : calcul et lecture géographique]
\begin{enumerate}
\item \textbf{Identifier} la valeur initiale $V_i$ (année de départ) et la valeur finale $V_f$ (année d'arrivée). Attention à ne pas les inverser.
\item \textbf{Appliquer la formule} : $t = \dfrac{V_f - V_i}{V_i} \times 100$.
\item \textbf{Vérifier le signe} : $t > 0$ = augmentation ; $t < 0$ = diminution.
\item \textbf{Rédiger l'interprétation} : « Entre ... et ..., la production de ... a augmenté/diminué de ... \%, passant de ... à ... ». Relie ensuite cette évolution à une cause géographique (programme de l'État, modernisation des équipements, essor de la pisciculture).
\end{enumerate}
\textbf{Piège classique} : diviser par $V_f$ au lieu de $V_i$, ou confondre le taux de variation avec la simple différence $V_f - V_i$ (qui est une variation absolue, pas un pourcentage).
\end{methode}

\begin{exoresolu}[Calculer un taux de variation : l'essor de la pisciculture]
\textbf{Énoncé.} La production piscicole du Cameroun est passée, en ordres de grandeur indicatifs, d'environ \num{1000} tonnes au début des années 2000 à environ \num{5000} tonnes au début des années 2020. (a) Calculez le taux de variation de la production entre ces deux dates. (b) Rédigez une phrase d'analyse. (c) Proposez deux explications à cette évolution. (5 points)
\tcblower
\textbf{Solution détaillée.}

\emph{(a) Calcul.} On a $V_i = \num{1000}$ tonnes et $V_f = \num{5000}$ tonnes.
\[ t = \frac{V_f - V_i}{V_i} \times 100 = \frac{\num{5000} - \num{1000}}{\num{1000}} \times 100 = \frac{\num{4000}}{\num{1000}} \times 100 = 400\,\%. \]

\emph{(b) Phrase d'analyse.} Entre le début des années 2000 et le début des années 2020, la production piscicole du Cameroun a été multipliée par cinq, soit une augmentation d'environ 400\,\%, passant d'environ \num{1000} à environ \num{5000} tonnes (ordres de grandeur).

\emph{(c) Explications.} Première explication : l'action de l'État, à travers le MINEPIA et divers programmes, a formé des pisciculteurs, distribué des alevins améliorés et encadré la construction d'étangs. Deuxième explication : la forte demande urbaine en poisson (Yaoundé, Douala) et le coût élevé des importations ont rendu la pisciculture rentable, attirant de nouveaux producteurs, y compris de jeunes entrepreneurs.

\emph{Conclusion.} L'augmentation spectaculaire (environ +400\,\%) confirme l'essor de la pisciculture, même si sa part dans la production totale de poisson reste faible (environ 2\,\%) : le potentiel de croissance demeure important.
\end{exoresolu}

\begin{attention}[Les 5 erreurs qui coûtent des points au Bac]
\begin{enumerate}
\item \textbf{Confondre carte et croquis.} La carte localise des données ; le croquis explique un fonctionnement avec des flèches et une légende « active ». Présenter un croquis sans flèches ni dynamique, c'est perdre la moitié des points.
\item \textbf{Oublier les éléments obligatoires de la carte} : titre, flèche du nord, légende complète. Un figuré utilisé sur la carte mais absent de la légende (ou l'inverse) est systématiquement sanctionné.
\item \textbf{Décrire ce qu'on sait au lieu de ce qu'on voit.} Dans la description d'une photographie, chaque affirmation doit s'appuyer sur un indice visible ; les connaissances du cours n'interviennent qu'à l'étape de l'interprétation.
\item \textbf{Se tromper de dénominateur dans le taux de variation} : on divise toujours par la valeur \emph{initiale} $V_i$. Diviser par $V_f$, ou donner la simple différence $V_f - V_i$ en guise de pourcentage, fausse tout le calcul.
\item \textbf{Ne pas citer les documents dans une étude de cas.} Une copie qui ne mentionne jamais « (doc. 1) », « (doc. 2) »... ne peut pas valider la compétence « inventorier, classer et interpréter des documents ». Cite chaque document au moins une fois, sans le recopier.
\end{enumerate}
\end{attention}

\newpage

\coursec{Exercices}

\courssub{Je vérifie mes connaissances}

\begin{exercice}[QCM : les notions clés de la leçon]
Pour chaque question, une seule réponse est exacte. Recopie la lettre correspondante sur ta copie.

\begin{enumerate}
\item La \cle{transhumance} est :
\begin{enumerate}
\item[a)] l'élevage de poissons en étang ;
\item[b)] le déplacement saisonnier des troupeaux entre zones de pâturage ;
\item[c)] la construction de bergeries modernes ;
\item[d)] la pêche au filet maillant.
\end{enumerate}
\item Un \cle{ranch} est :
\begin{enumerate}
\item[a)] un petit étang de pisciculture familiale ;
\item[b)] une vaste exploitation d'élevage clôturée, avec pâturages aménagés et soins vétérinaires ;
\item[c)] un marché à bétail hebdomadaire ;
\item[d)] une pirogue à moteur hors-bord.
\end{enumerate}
\item Le \cle{PDMA} (Programme de développement des marchés agricoles) illustre :
\begin{enumerate}
\item[a)] une initiative privée étrangère sans lien avec l'État ;
\item[b)] l'action de l'État dans le financement et l'encadrement des activités agro-pastorales ;
\item[c)] une taxe sur les importations de poisson ;
\item[d)] une organisation de pêcheurs artisans.
\end{enumerate}
\item La \cle{pisciculture} désigne :
\begin{enumerate}
\item[a)] la pêche en haute mer ;
\item[b)] l'élevage de poissons en milieu contrôlé (étangs, bassins, cages) ;
\item[c)] la conservation du poisson par fumage ;
\item[d)] la vente du poisson au marché.
\end{enumerate}
\end{enumerate}
\end{exercice}

\begin{exercice}[Vrai ou faux, avec justification]
Indique si chaque affirmation est vraie ou fausse, puis justifie ta réponse en une ou deux phrases.

\begin{enumerate}
\item La modernisation de la pêche artisanale au Cameroun passe notamment par l'usage de grandes pirogues à moteur et de filets adaptés.
\item La transhumance a totalement disparu du Cameroun depuis les années 1980.
\item L'amélioration de l'alimentation et des soins vétérinaires fait partie des transformations de l'élevage traditionnel.
\item Le Cameroun n'importe aucun poisson : sa production couvre largement ses besoins.
\item L'État camerounais intervient dans la formation des acteurs et la régulation de la commercialisation des produits agro-pastoraux.
\end{enumerate}
\end{exercice}

\begin{exercice}[Définitions express]
Définis en deux ou trois phrases chacun des termes suivants, en donnant pour chacun un exemple camerounais :
\begin{enumerate}
\item bergerie ;
\item étang ;
\item pêche artisanale ;
\item croisement (en élevage).
\end{enumerate}
\end{exercice}

\begin{exercice}[Calculs directs sur des données halieutiques]
Le tableau ci-dessous donne des ordres de grandeur de la production de poisson au Cameroun (en milliers de tonnes, données arrondies).

\begin{center}
\begin{tabularx}{\linewidth}{|l|X|X|}\hline
\rowcolor{popblueL} \textbf{Source} & \textbf{Année 1} & \textbf{Année 2} \\ \hline
Pêche maritime artisanale & 100 & 120 \\ \hline
Pêche continentale (fleuves, lacs) & 60 & 55 \\ \hline
Pisciculture & 2 & 8 \\ \hline
\end{tabularx}
\end{center}

\begin{enumerate}
\item Calcule la production totale de poisson pour chaque année.
\item Calcule le taux de variation de la production piscicole entre l'année 1 et l'année 2 (en \%).
\item Calcule la part (en \%) de la pêche maritime artisanale dans la production totale de l'année 2.
\item Que montre la comparaison entre l'évolution de la pêche continentale et celle de la pisciculture ? (2 lignes)
\end{enumerate}
\end{exercice}

\courssub{Je m'entraîne}

\begin{exercice}[Associer transformation et secteur]
Recopie le tableau ci-dessous et complète-le en plaçant chaque transformation dans la colonne qui convient : grandes pirogues à moteur ; bergeries et poulaillers ; filets adaptés ; sélection et croisements des espèces ; étangs aménagés ; glace et chambres froides ; aliments composés pour poissons ; vaccination et soins vétérinaires.

\begin{center}
\begin{tabularx}{\linewidth}{|X|X|X|}\hline
\rowcolor{popblueL} \textbf{Élevage} & \textbf{Pêche artisanale} & \textbf{Pisciculture} \\ \hline
 & & \\ \hline
 & & \\ \hline
 & & \\ \hline
\end{tabularx}
\end{center}

Choisis ensuite une transformation par colonne et explique en trois lignes en quoi elle améliore la production ou les revenus des producteurs.
\end{exercice}

\begin{exercice}[Commentaire court d'une photographie décrite]
Photographie (décrite) : « Sur les hauts plateaux de l'Adamaoua, un troupeau de zébus est parqué la nuit dans un enclos en bois ; à proximité, un hangar abrite du fourrage et un point d'eau aménagé ; un agent vétérinaire vaccine un veau. »

\begin{enumerate}
\item Identifie et localise la forme d'élevage représentée.
\item Relève dans la description trois indices de modernisation de l'élevage traditionnel.
\item Explique en cinq lignes comment ces transformations renforcent l'intégration nationale (approvisionnement des villes du Sud en viande).
\end{enumerate}
\end{exercice}

\begin{exercice}[Construire un diagramme en bâtons]
On donne les captures de pêche maritime artisanale par façade (ordres de grandeur, en milliers de tonnes) :

\begin{center}
\begin{tabularx}{\linewidth}{|l|X|}\hline
\rowcolor{popblueL} \textbf{Zone de pêche} & \textbf{Captures (milliers de tonnes)} \\ \hline
Littoral de Douala (Wouri) & 55 \\ \hline
Littoral de Kribi (Océan) & 30 \\ \hline
Littoral de Limbé--Idenau (Sud-Ouest) & 25 \\ \hline
Lagune Ebrié du Cameroun (Mouanko, Sanaga-Maritime) & 10 \\ \hline
\end{tabularx}
\end{center}

\begin{enumerate}
\item Construis un diagramme en bâtons représentant ces captures (axe vertical gradué de 5 en 5).
\item Calcule la part de la façade de Douala dans le total (en \%).
\item Rédige un commentaire de cinq lignes : hiérarchie des zones, explication possible (proximité des marchés urbains, équipements portuaires).
\end{enumerate}
\end{exercice}

\begin{exercice}[Question à réponse construite]
« Les transformations de l'élevage traditionnel au Cameroun. »

Rédige une réponse construite de 15 à 20 lignes, en suivant la démarche : (1) rappelle les traits de l'élevage traditionnel (extensif, transhumant, faible productivité) ; (2) présente au moins quatre transformations (sélection et croisements, alimentation, soins, bâtiments d'élevage, ranchs) ; (3) montre les limites qui persistent (coût des intrants, conflits agriculteurs--éleveurs, concurrence des importations). Appuie ta réponse sur deux exemples régionaux précis (Adamaoua, Extrême-Nord, Ouest\ldots).
\end{exercice}

\courssub{Je me prépare au Bac}

\begin{exercice}[Dissertation type Bac]
Sujet : \emph{« Les transformations de l'élevage et de la pêche au Cameroun : progrès et limites. »}

Consignes : rédige une dissertation comprenant une introduction problématisée (accroche située au Cameroun, définition des termes, problématique, annonce du plan), un développement en deux parties (I. les progrès permis par la modernisation de l'élevage, de la pêche et l'essor de la pisciculture ; II. les limites et difficultés qui freinent cette mutation), et une conclusion ouverte (par exemple sur la place du « consommer local »). Soigne la langue, illustre chaque argument d'un exemple camerounais précis (région, ville, projet).
\end{exercice}

\begin{exercice}[Étude de documents : carte et tableau des pêches]
\textbf{Document 1 (carte décrite)} : carte des principales zones de pêche du Cameroun : côte atlantique (Douala, Kribi, Limbé), fleuves Sanaga, Wouri et Logone, lac Tchad, lac de Lagdo, retenue de Mbakaou, étangs piscicoles de la région de l'Ouest (Foumban, Bafoussam) et du Centre (Yaoundé, Obala).

\textbf{Document 2 (tableau)} : production par origine (ordres de grandeur, milliers de tonnes) :

\begin{center}
\begin{tabularx}{\linewidth}{|l|X|X|}\hline
\rowcolor{popblueL} \textbf{Origine} & \textbf{2010} & \textbf{2020} \\ \hline
Pêche maritime & 95 & 115 \\ \hline
Pêche continentale & 70 & 60 \\ \hline
Pisciculture & 1 & 9 \\ \hline
\end{tabularx}
\end{center}

\begin{enumerate}
\item À partir du document 2, construis un diagramme en bâtons groupés comparant les trois origines pour les deux années.
\item Calcule le taux de variation de chaque origine entre 2010 et 2020.
\item Décris la répartition spatiale des zones de pêche à partir du document 1 (4 lignes).
\item Rédige un commentaire organisé de 15 lignes montrant : la stagnation ou le recul des pêches de capture, le dynamisme de la pisciculture, et les conséquences sur les importations de poisson.
\end{enumerate}
\end{exercice}

\begin{exercice}[Sujet de synthèse : le rôle de l'État]
Dossier de trois documents :

\textbf{Document 3 (texte)} : « L'État camerounais multiplie les programmes d'encadrement et de financement des activités agro-pastorales et piscicoles : appui aux marchés agricoles (PDMA), projets de développement de l'élevage, distribution d'alevins et d'aliments, formation des producteurs, organisation des marchés à bétail et régulation de la commercialisation. »

\textbf{Document 4 (photographie décrite)} : un marché à bétail hebdomadaire dans l'Extrême-Nord : centaines de zébus, négociants venus de Yaoundé et de Douala, camions bétaillères, agents municipaux contrôlant les transactions.

\textbf{Document 5 (tableau des financements, valeurs fictives d'entraînement, milliards de FCFA)} :

\begin{center}
\begin{tabularx}{\linewidth}{|l|X|X|}\hline
\rowcolor{popblueL} \textbf{Programme} & \textbf{Montant} & \textbf{Bénéficiaires} \\ \hline
PDMA & 30 & 60\,000 producteurs \\ \hline
Projets élevage et pastoralisme & 20 & 25\,000 éleveurs \\ \hline
Appui à la pisciculture & 5 & 8\,000 pisciculteurs \\ \hline
\end{tabularx}
\end{center}

\begin{enumerate}
\item Présente chaque document en une phrase (nature, thème, source supposée).
\item À partir du document 5, calcule la part du PDMA dans le financement total (en \%) et représente les trois montants par un diagramme circulaire ou en bâtons.
\item Dégage les trois grandes formes d'intervention de l'État (financement, formation, régulation) en citant les documents.
\item Rédige une synthèse de 20 lignes répondant à la question : « Comment l'État camerounais soutient-il le développement des activités agro-pastorales et piscicoles, et avec quelles limites ? »
\end{enumerate}
\end{exercice}

\newpage

\coursec{Corrigés des exercices}

\begin{corrige}[Exercice 1 — QCM : les notions clés de la leçon]
\textbf{Question 1 : réponse b).} La transhumance est le déplacement saisonnier des troupeaux entre zones de pâturage : au Cameroun, les éleveurs peuls conduisent leurs zébus de l'Extrême-Nord et de l'Adamaoua vers les pâturages du Nord-Ouest ou du Centre pendant la saison sèche, puis reviennent à la saison des pluies. Les réponses a) et d) relèvent de la pêche, et c) d'une transformation de l'élevage, pas d'un déplacement.

\textbf{Question 2 : réponse b).} Un ranch est une vaste exploitation d'élevage clôturée, avec pâturages aménagés, points d'eau et soins vétérinaires ; le ranch de Wakwa, près de Ngaoundéré (Adamaoua), en est l'exemple camerounais classique.

\textbf{Question 3 : réponse b).} Le PDMA (Programme de développement des marchés agricoles) illustre l'action de l'État dans le financement et l'encadrement des activités agro-pastorales : il appuie les producteurs, améliore l'accès aux marchés et finance des infrastructures rurales.

\textbf{Question 4 : réponse b).} La pisciculture est l'élevage de poissons en milieu contrôlé (étangs, bassins, cages) ; elle se développe notamment dans l'Ouest (Foumban, Bafoussam) et le Centre (Obala). Le fumage (c) est une technique de conservation, pas un élevage.
\end{corrige}

\begin{corrige}[Exercice 2 — Vrai ou faux, avec justification]
\textbf{1. Vrai.} La modernisation de la pêche artisanale repose sur les grandes pirogues à moteur hors-bord (qui élargissent le rayon d'action), les filets adaptés (filets maillants, sennes) et la conservation (glace, chambres froides), visibles à Youpwè (Douala), Kribi ou Limbé.

\textbf{2. Faux.} La transhumance existe toujours : chaque saison sèche, des troupeaux quittent l'Extrême-Nord et l'Adamaoua vers le Nord-Ouest, le Centre ou le Tchad voisin. Elle est toutefois encadrée (couloirs de transhumance) et concurrencée par l'élevage sédentaire en ranchs.

\textbf{3. Vrai.} L'amélioration de l'alimentation (fourrages conservés, aliments composés) et des soins (vaccination, déparasitage, vétérinaires) fait partie, avec la sélection, les croisements et les bâtiments d'élevage, des transformations de l'élevage traditionnel.

\textbf{4. Faux.} La production nationale (de l'ordre de 150\,000 à 200\,000 tonnes par an, ordre de grandeur) ne couvre pas la demande : le Cameroun importe massivement du poisson congelé (plus de 200\,000 tonnes par an, ordre de grandeur), surtout du maquereau, pour alimenter les villes.

\textbf{5. Vrai.} L'État intervient par la formation des acteurs (centres de formation, encadrement technique du MINEPIA), le financement (PDMA, projets d'élevage) et la régulation de la commercialisation (marchés à bétail contrôlés, normes sanitaires).
\end{corrige}

\begin{corrige}[Exercice 3 — Définitions express]
\textbf{1. Bergerie.} Bâtiment ou enclos couvert destiné à abriter le petit bétail (ovins, caprins) et parfois les bovins, pour les protéger des intempéries, des prédateurs et faciliter les soins. \emph{Exemple :} bergeries construites dans les projets d'élevage de l'Adamaoua (région de Ngaoundéré).

\textbf{2. Étang.} Plan d'eau peu profond, naturel ou artificiel, aménagé et contrôlé pour l'élevage de poissons (pisciculture) : on y introduit des alevins, on les nourrit et on les récolte. \emph{Exemple :} étangs piscicoles de Foumban et de Bafoussam (Ouest) ou d'Obala (Centre).

\textbf{3. Pêche artisanale.} Pêche pratiquée à petite échelle, avec des embarcations simples (pirogues, parfois motorisées) et des engins peu mécanisés (filets, lignes, nasses), par des pêcheurs individuels ou des groupements, pour l'autoconsommation et la vente locale. \emph{Exemple :} pêcheurs de Youpwè à Douala ou de Kribi.

\textbf{4. Croisement (en élevage).} Accouplement d'animaux de races différentes pour obtenir des descendants combinant les qualités des deux races (rusticité locale et productivité des races améliorées). \emph{Exemple :} croisement de zébus locaux avec des races laitières ou à viande importées au ranch de Wakwa (Adamaoua).
\end{corrige}

\begin{corrige}[Exercice 4 — Calculs directs sur des données halieutiques]
\textbf{1. Production totale.}
\begin{align*}
\text{Année 1} &: 100 + 60 + 2 = 162 \text{ milliers de tonnes}\\
\text{Année 2} &: 120 + 55 + 8 = 183 \text{ milliers de tonnes}
\end{align*}

\textbf{2. Taux de variation de la pisciculture.}
\[
t = \frac{8 - 2}{2}\times 100 = \frac{6}{2}\times 100 = 300\ \%
\]
La production piscicole a été multipliée par 4 (soit une hausse de 300 \%).

\textbf{3. Part de la pêche maritime artisanale en année 2.}
\[
p = \frac{120}{183}\times 100 \approx 65,6\ \%
\]
Soit environ 66 \% de la production totale.

\textbf{4. Comparaison.} La pêche continentale recule (de 60 à 55 milliers de tonnes, soit $-8,3\ \%$), sous l'effet de la surpêche et de la dégradation des milieux, tandis que la pisciculture explose ($+300\ \%$) : l'élevage de poissons compense progressivement l'essoufflement des pêches de capture.

\begin{attention}[Points de vigilance]
Ne pas confondre taux de variation (en \%) et coefficient multiplicateur : multiplier par 4 correspond à une hausse de 300 \%, pas de 400 \%. Pour la part en \%, diviser toujours la partie par le total de la \emph{même} année (183, pas 162).
\end{attention}
\end{corrige}

\begin{corrige}[Exercice 5 — Associer transformation et secteur]
Tableau complété :

\begin{center}
\begin{tabularx}{\linewidth}{|X|X|X|}\hline
\rowcolor{popblueL} \textbf{Élevage} & \textbf{Pêche artisanale} & \textbf{Pisciculture} \\ \hline
Bergeries et poulaillers & Grandes pirogues à moteur & Étangs aménagés \\ \hline
Sélection et croisements des espèces & Filets adaptés & Aliments composés pour poissons \\ \hline
Vaccination et soins vétérinaires & Glace et chambres froides & (Alevins sélectionnés — complément possible) \\ \hline
\end{tabularx}
\end{center}

\textbf{Explications (une par colonne) :}
\begin{itemize}
\item \emph{Élevage — vaccination et soins vétérinaires :} ils réduisent la mortalité des troupeaux (peste bovine, péripneumonie), ce qui augmente le nombre d'animaux commercialisables et donc les revenus de l'éleveur.
\item \emph{Pêche artisanale — grandes pirogues à moteur :} elles permettent d'aller plus loin et plus longtemps en mer, d'atteindre des zones de pêche moins exploitées et d'accroître les captures journalières.
\item \emph{Pisciculture — aliments composés :} ils accélèrent la croissance des poissons (tilapias, clarias), raccourcissent le cycle de production et augmentent le rendement à l'hectare d'étang, donc le chiffre d'affaires du pisciculteur.
\end{itemize}
\end{corrige}

\begin{corrige}[Exercice 6 — Commentaire court d'une photographie décrite]
\textbf{1. Identification et localisation.} Il s'agit d'un élevage bovin (zébus) en voie de sédentarisation, sur les hauts plateaux de l'Adamaoua (région de Ngaoundéré), grande région d'élevage du Cameroun.

\textbf{2. Trois indices de modernisation.} (a) L'enclos en bois où le troupeau est parqué la nuit (bâtiment d'élevage, sécurisation) ; (b) le hangar à fourrage et le point d'eau aménagé (amélioration de l'alimentation et de l'abreuvement) ; (c) la vaccination d'un veau par un agent vétérinaire (amélioration des soins).

\textbf{3. Lien avec l'intégration nationale (5 lignes).} Ces transformations augmentent la production de viande de l'Adamaoua, qui approvisionne par la route et le chemin de fer (Transcamerounais) les grandes villes du Sud (Yaoundé, Douala). L'élevage modernisé crée ainsi des flux commerciaux Nord--Sud, fait vivre négociants, transporteurs et bouchers, et intègre les régions d'élevage à l'économie nationale : le Nord pastoral fournit la viande, le Sud urbain fournit les marchés et les revenus.

\begin{attention}[Points de vigilance]
Dans un commentaire de photographie, toujours procéder en trois temps : identifier/localiser, décrire (en citant des éléments précis du document), interpréter (relier au problème posé, ici l'intégration nationale). Ne jamais se contenter de réciter la leçon sans mobiliser le document.
\end{attention}
\end{corrige}

\begin{corrige}[Exercice 7 — Construire un diagramme en bâtons]
\textbf{1. Diagramme en bâtons.} Total des captures : $55 + 30 + 25 + 10 = 120$ milliers de tonnes.

\begin{popfigure}
\begin{tikzpicture}
\begin{axis}[
  width=0.95\linewidth, height=6.5cm,
  ybar, bar width=0.9cm,
  ymin=0, ymax=60,
  ytick={0,5,10,15,20,25,30,35,40,45,50,55,60},
  ylabel={Captures (milliers de tonnes)},
  symbolic x coords={Douala, Kribi, Limbé--Idenau, Mouanko},
  xtick=data,
  x tick label style={rotate=15, anchor=east, font=\small},
  nodes near coords, every node near coord/.append style={font=\small},
  axis lines=left, grid=major, grid style={popline!40}
]
\addplot[fill=popblue!70, draw=popblue] coordinates {(Douala,55) (Kribi,30) (Limbé--Idenau,25) (Mouanko,10)};
\end{axis}
\end{tikzpicture}
\legende{Captures de pêche maritime artisanale par façade (ordres de grandeur, milliers de tonnes ; données arrondies d'entraînement).}
\end{popfigure}

\textbf{2. Part de la façade de Douala.}
\[
p = \frac{55}{120}\times 100 \approx 45,8\ \%
\]
Soit environ 46 \% du total.

\textbf{3. Commentaire (5 lignes).} Le littoral de Douala (Wouri) domine nettement avec près de la moitié des captures, devant Kribi (30), Limbé--Idenau (25) et la lagune de Mouanko (10). Cette hiérarchie s'explique par la proximité du plus grand marché urbain du pays (Douala), par les équipements portuaires et de conservation (glace, chambres froides) et par l'ancienneté des communautés de pêcheurs (Youpwè). Kribi et Limbé bénéficient aussi de ports et de marchés régionaux, tandis que Mouanko, en milieu lagunaire enclavé, reste une pêche de proximité.

\begin{attention}[Points de vigilance]
Graduer l'axe vertical de 5 en 5 comme demandé, nommer les axes, donner un titre et une source au diagramme : un graphique sans titre ni unité est pénalisé au Bac.
\end{attention}
\end{corrige}

\begin{corrige}[Exercice 8 — Question à réponse construite]
\textbf{Corrigé rédigé (15 à 20 lignes).}

L'élevage traditionnel camerounais, pratiqué surtout dans l'Adamaoua, le Nord et l'Extrême-Nord, est un élevage extensif : de grands troupeaux de zébus pâturent librement sur de vastes espaces, souvent au gré de la transhumance saisonnière, avec une faible productivité (viande et lait par tête limitées, forte mortalité).

Cet élevage connaît de profondes transformations. D'abord, la sélection et les croisements améliorent les races : au ranch de Wakwa (Adamaoua), des zébus locaux sont croisés avec des races améliorées pour accroître le rendement en viande et en lait. Ensuite, l'alimentation progresse : conservation du fourrage, compléments alimentaires, points d'eau aménagés. Les soins se développent également : campagnes de vaccination, déparasitage, présence de vétérinaires, ce qui réduit la mortalité. Enfin, les bâtiments d'élevage (bergeries, poulaillers, enclos) et les ranchs clôturés sédentarisent une partie des éleveurs ; dans l'Ouest, les poulaillers modernes ont fait du pays un producteur majeur de volailles et d'œufs.

Mais des limites persistent : le coût élevé des intrants (aliments, médicaments vétérinaires) exclut les petits éleveurs ; les conflits agriculteurs--éleveurs se multiplient (Nord-Ouest, Adamaoua) quand les troupeaux détruisent les champs ; enfin, la concurrence des importations (viande congelée, volailles) pèse sur les prix. La mutation de l'élevage reste donc inachevée et inégale.

\textbf{Barème indicatif (sur 10).} Traits de l'élevage traditionnel : 2 pts ; quatre transformations correctement présentées : 4 pts ; limites : 2 pts ; deux exemples régionaux précis : 1 pt ; clarté et langue : 1 pt.

\begin{attention}[Points de vigilance]
Respecter la démarche en trois temps imposée par la consigne ; citer au moins deux régions nommément ; ne pas présenter la transhumance comme une « transformation » : c'est un trait de l'élevage traditionnel.
\end{attention}
\end{corrige}

\begin{corrige}[Exercice 9 — Dissertation type Bac]
\textbf{Corrigé rédigé.}

\emph{Introduction.} Sur le marché de Douala, le poisson congelé importé côtoie le capitaine fumé de Kribi, tandis que la viande de zébu de l'Adamaoua arrive par camions entiers : l'élevage et la pêche camerounais sont en pleine mutation. Longtemps dominés par des pratiques traditionnelles — élevage extensif et transhumant, pêche artisanale à la pirogue —, ces secteurs se modernisent sous l'effet de la sélection, des équipements et de l'essor de la pisciculture. On peut alors se demander : dans quelle mesure les transformations de l'élevage et de la pêche constituent-elles un progrès, et quelles limites freinent encore cette mutation ? Nous montrerons d'abord les progrès permis par la modernisation, puis les limites qui persistent.

\emph{I. Les progrès permis par la modernisation.} L'élevage se transforme : sélection et croisements des races (ranch de Wakwa, Adamaoua), amélioration de l'alimentation (fourrages conservés, aliments composés), soins vétérinaires (vaccinations) et construction de bergeries et de poulaillers, notamment dans l'Ouest, premier bassin avicole du pays. La pêche artisanale se modernise : grandes pirogues à moteur hors-bord à Youpwè et Kribi, filets adaptés, glace et chambres froides qui réduisent les pertes après capture. Enfin, la pisciculture connaît un essor spectaculaire : étangs de Foumban et de Bafoussam (Ouest), d'Obala (Centre), production multipliée par plusieurs fois en une décennie. Ces progrès augmentent les quantités produites, les revenus des acteurs et l'approvisionnement des villes, renforçant l'intégration nationale.

\emph{II. Les limites et difficultés.} La mutation reste inégale : le coût des intrants (aliments, médicaments, carburant) exclut les petits producteurs ; les conflits agriculteurs--éleveurs éclatent lors de la transhumance (Nord-Ouest, Adamaoua) ; les pêches de capture stagnent ou reculent par surpêche, notamment en pêche continentale (Logone, lac Tchad). Surtout, la production nationale reste insuffisante : le Cameroun importe chaque année des centaines de milliers de tonnes de poisson congelé (ordre de grandeur), ce qui pèse sur la balance commerciale et concurrence les producteurs locaux. L'encadrement de l'État (PDMA, projets d'élevage) demeure limité face à l'ampleur des besoins.

\emph{Conclusion.} Les transformations de l'élevage et de la pêche ont indéniablement modernisé la production camerounaise, mais elles n'ont pas encore assuré l'autosuffisance. L'essor de la pisciculture et la promotion du « consommer local » pourraient demain réduire la dépendance aux importations : encore faut-il que les consommateurs préfèrent le poisson et la viande du terroir.

\textbf{Barème indicatif (sur 20).} Introduction (accroche située, définitions, problématique, plan) : 4 pts ; partie I argumentée et illustrée : 6 pts ; partie II argumentée et illustrée : 6 pts ; conclusion ouverte : 2 pts ; langue et organisation : 2 pts.

\begin{attention}[Points de vigilance]
La problématique doit être une vraie question ; chaque argument doit être suivi d'un exemple camerounais précis (région, ville, projet) ; éviter le catalogue de notions sans articulation ; annoncer le plan en fin d'introduction.
\end{attention}
\end{corrige}

\begin{corrige}[Exercice 10 — Étude de documents : carte et tableau des pêches]
\textbf{1. Diagramme en bâtons groupés.}

\begin{popfigure}
\begin{tikzpicture}
\begin{axis}[
  width=0.95\linewidth, height=7cm,
  ybar, bar width=0.45cm,
  ymin=0, ymax=130,
  ylabel={Production (milliers de tonnes)},
  symbolic x coords={Maritime, Continentale, Pisciculture},
  xtick=data,
  legend style={at={(0.02,0.98)}, anchor=north west, font=\small},
  axis lines=left, grid=major, grid style={popline!40}
]
\addplot[fill=popblue!70, draw=popblue] coordinates {(Maritime,95) (Continentale,70) (Pisciculture,1)};
\addplot[fill=poporange!80, draw=poporange] coordinates {(Maritime,115) (Continentale,60) (Pisciculture,9)};
\legend{2010, 2020}
\end{axis}
\end{tikzpicture}
\legende{Production de poisson par origine au Cameroun, 2010 et 2020 (ordres de grandeur, milliers de tonnes ; données arrondies d'entraînement).}
\end{popfigure}

\textbf{2. Taux de variation entre 2010 et 2020.}
\begin{align*}
\text{Maritime} &: \frac{115-95}{95}\times 100 \approx +21,1\ \%\\
\text{Continentale} &: \frac{60-70}{70}\times 100 \approx -14,3\ \%\\
\text{Pisciculture} &: \frac{9-1}{1}\times 100 = +800\ \%
\end{align*}

\textbf{3. Répartition spatiale (4 lignes).} Les zones de pêche se répartissent en trois ensembles : le littoral atlantique (Douala, Kribi, Limbé), domaine de la pêche maritime ; les eaux continentales du Nord (lac Tchad, Logone, retenue de Lagdo, Mbakaou) et du Centre-Sud (Sanaga, Wouri) ; enfin les étangs piscicoles, concentrés dans l'Ouest (Foumban, Bafoussam) et le Centre (Yaoundé, Obala), près des marchés urbains.

\textbf{4. Commentaire organisé (15 lignes).} Le document 2 révèle deux dynamiques contrastées. Les pêches de capture stagnent ou reculent : la pêche maritime ne progresse que d'environ 21 \% en dix ans, malgré la motorisation des pirogues, car les ressources côtières sont surexploitées ; la pêche continentale recule même de 14 \% environ, sous l'effet conjugué de la surpêche, de la réduction des plans d'eau (retrait du lac Tchad) et de la dégradation des milieux. À l'inverse, la pisciculture affiche un dynamisme spectaculaire : multipliée par neuf ($+800\ \%$), elle passe d'une activité marginale à un secteur en plein essor, porté par l'aménagement d'étangs dans l'Ouest et le Centre, la distribution d'alevins et la demande urbaine. Toutefois, ce dynamisme ne suffit pas à combler le déficit : la production totale (environ 184 milliers de tonnes en 2020 d'après le tableau) reste inférieure à la consommation nationale, si bien que le Cameroun importe massivement du poisson congelé — plusieurs centaines de milliers de tonnes par an (ordre de grandeur) —, ce qui alourdit la facture des importations et concurrence les pêcheurs locaux. L'enjeu est donc de poursuivre le développement de la pisciculture tout en gérant durablement les pêcheries naturelles.

\textbf{Barème indicatif (sur 20).} Diagramme correct (axes, légende, titre) : 4 pts ; trois taux de variation exacts : 3 pts ; description spatiale : 3 pts ; commentaire organisé (trois idées, chiffres mobilisés) : 8 pts ; langue : 2 pts.

\begin{attention}[Points de vigilance]
Un taux de variation négatif doit être signalé comme un recul ; multiplier par 9 équivaut à $+800\ \%$, pas à $+900\ \%$ ; dans le commentaire, citer les chiffres calculés et les lieux du document 1.
\end{attention}
\end{corrige}

\begin{corrige}[Exercice 11 — Sujet de synthèse : le rôle de l'État]
\textbf{1. Présentation des documents.}
\begin{itemize}
\item \emph{Document 3 :} texte d'orientation (type rapport ministériel ou manuel) présentant les programmes étatiques d'encadrement et de financement des activités agro-pastorales et piscicoles.
\item \emph{Document 4 :} photographie de terrain montrant un marché à bétail hebdomadaire dans l'Extrême-Nord, illustrant la commercialisation du bétail et son contrôle.
\item \emph{Document 5 :} tableau statistique (valeurs fictives d'entraînement) des financements publics par programme, avec montants et bénéficiaires.
\end{itemize}

\textbf{2. Calculs et représentation.} Financement total : $30 + 20 + 5 = 55$ milliards de FCFA.
\[
\text{Part du PDMA} = \frac{30}{55}\times 100 \approx 54,5\ \%
\]
(Si l'on représentait ces parts en diagramme circulaire, les angles seraient : PDMA $\approx 196^{\circ}$, Élevage $\approx 131^{\circ}$, Pisciculture $\approx 33^{\circ}$.)

\begin{popfigure}
\begin{tikzpicture}
\begin{axis}[
  width=0.9\linewidth, height=5.5cm,
  ybar, bar width=1cm,
  ymin=0, ymax=35,
  ylabel={Montant (milliards de FCFA)},
  symbolic x coords={PDMA, Élevage, Pisciculture},
  xtick=data,
  nodes near coords, every node near coord/.append style={font=\small},
  axis lines=left, grid=major, grid style={popline!40}
]
\addplot[fill=popgreen!70, draw=popgreen!60!black] coordinates {(PDMA,30) (Élevage,20) (Pisciculture,5)};
\end{axis}
\end{tikzpicture}
\legende{Financements publics par programme (valeurs fictives d'entraînement, milliards de FCFA).}
\end{popfigure}

\textbf{3. Les trois formes d'intervention de l'État.}
\begin{itemize}
\item \emph{Financement :} le PDMA (30 milliards, 60\,000 bénéficiaires) et les projets d'élevage et de pisciculture (documents 3 et 5) ;
\item \emph{Formation et encadrement technique :} formation des producteurs, distribution d'alevins et d'aliments (document 3) ;
\item \emph{Régulation de la commercialisation :} organisation et contrôle des marchés à bétail, agents municipaux vérifiant les transactions (documents 3 et 4).
\end{itemize}

\textbf{4. Synthèse rédigée (20 lignes).}

L'État camerounais joue un rôle moteur dans le développement des activités agro-pastorales et piscicoles, en agissant à tous les maillons des filières. Il intervient d'abord par le financement : le Programme de développement des marchés agricoles (PDMA), qui concentre à lui seul plus de la moitié des financements du document 5 (environ 54,5 \%), appuie 60\,000 producteurs en facilitant l'accès aux intrants, aux équipements et aux marchés ; des projets spécifiques soutiennent 25\,000 éleveurs et 8\,000 pisciculteurs. Il intervient ensuite par l'encadrement et la formation : distribution d'alevins et d'aliments composés, appui technique du MINEPIA, formation des éleveurs aux techniques modernes (vaccination, construction de bergeries) et des pisciculteurs à la gestion des étangs. Il intervient enfin par la régulation de la commercialisation : les marchés à bétail hebdomadaires de l'Extrême-Nord, comme le montre le document 4, sont organisés et contrôlés par des agents municipaux, ce qui sécurise les transactions entre éleveurs peuls et négociants venus de Yaoundé et de Douala, et alimente les villes du Sud en viande — contribuant ainsi à l'intégration nationale.

Cependant, cette action connaît des limites. Les montants restent modestes au regard des besoins : la pisciculture, pourtant en plein essor, ne reçoit que 5 milliards, soit moins de 10 \% du total. L'encadrement technique est insuffisant dans les zones enclavées, les intrants demeurent chers pour les petits producteurs, et les conflits agriculteurs--éleveurs ne sont pas toujours prévenus. Enfin, malgré ces efforts, le Cameroun continue d'importer massivement du poisson, signe que l'appui de l'État n'a pas encore permis l'autosuffisance. L'action publique, indispensable, doit donc être renforcée et mieux ciblée, en faveur notamment de la pisciculture et des petits producteurs.

\textbf{Barème indicatif (sur 20).} Présentation des documents : 3 pts ; calculs et diagramme : 4 pts ; dégagement des trois formes d'intervention avec citations : 4 pts ; synthèse (structure, mobilisation des trois documents, limites) : 7 pts ; langue : 2 pts.

\begin{attention}[Points de vigilance]
Une synthèse de documents ne doit ni recopier les documents ni faire un exposé personnel : il faut croiser les informations et citer les documents entre parenthèses. Signaler que les valeurs du document 5 sont fictives (données d'entraînement) ; vérifier que la somme des parts fait bien 100 \% ($54,5 + 36,4 + 9,1 = 100\ \%$).
\end{attention}
\end{corrige}

\newpage

\coursec{Fiche bilan}

\begin{popfigure}
\begin{tikzpicture}[mindmap, grow cyclic, every node/.style={concept, text=white, font=\small},
  root concept/.append style={concept color=popdark, font=\bfseries\small, minimum size=2.6cm},
  level 1/.append style={level distance=3.4cm, sibling angle=60},
  level 1 concept/.append style={font=\scriptsize\bfseries, minimum size=1.9cm},
  level 2/.append style={level distance=2.2cm, sibling angle=45},
  level 2 concept/.append style={font=\tiny, minimum size=1.5cm, text=popdark}]
\node[root concept]{Élevage et pêche traditionnels en mutation}
  child[concept color=popblue]{ node {Sélection et croisements}
    child[concept color=popblueL]{ node {Zébu $\times$ races améliorées} }
    child[concept color=popblueL]{ node {Poules pondeuses, porcs améliorés} }
  }
  child[concept color=popgreen]{ node {Alimentation et soins}
    child[concept color=popgreenL]{ node {Provendes, fourrages} }
    child[concept color=popgreenL]{ node {Vaccination, vétérinaires} }
  }
  child[concept color=poporange]{ node {Bâtiments d'élevage}
    child[concept color=poporangeL]{ node {Bergeries, poulaillers} }
    child[concept color=poporangeL]{ node {Ranches, transhumance encadrée} }
  }
  child[concept color=poppurple]{ node {Pêche modernisée}
    child[concept color=poppurpleL]{ node {Pirogues à moteur} }
    child[concept color=poppurpleL]{ node {Filets adaptés, glace, fumoirs} }
  }
  child[concept color=poppink]{ node {Pisciculture}
    child[concept color=poppinkL]{ node {Étangs, tilapias, carpes} }
    child[concept color=poppinkL]{ node {Alevinage, aliments composés} }
  }
  child[concept color=popgold]{ node {Rôle de l'État}
    child[concept color=popgoldL]{ node {PDMA, MINEPIA, projets} }
    child[concept color=popgoldL]{ node {Formation, financement, régulation} }
  };
\end{tikzpicture}
\legende{Carte mentale de la leçon : les six branches à maîtriser (transformations de l'élevage en bleu, vert et orange ; pêche en violet ; pisciculture en rose ; action de l'État en doré).}
\end{popfigure}

\begin{bilan}
\textbf{Définitions clés à connaître par cœur}
\begin{itemize}
  \item \cle{Transhumance} : déplacement saisonnier des troupeaux et de leurs bergers entre les pâturages de saison sèche (vers le sud, zones de bas-fonds) et ceux de saison des pluies (vers le nord, pour fuir la mouche tsé-tsé et les inondations).
  \item \cle{Bergerie} : bâtiment couvert où les animaux (ovins, caprins, bovins) sont parqués, nourris et soignés ; signe du passage de l'élevage en divagation à l'élevage parqué ou semi-parqué.
  \item \cle{Ranch} : vaste exploitation d'élevage extensif mais organisée, clôturée, avec points d'eau, pistes et soins vétérinaires (ex. : ranchs de l'Adamaoua).
  \item \cle{Étang} : plan d'eau artificiel aménagé pour l'élevage de poissons ; unité de base de la pisciculture.
  \item \cle{Pisciculture} : élevage contrôlé de poissons (tilapias, carpes, clarias) en étangs ou en cages.
  \item \cle{Pêche artisanale} : pêche pratiquée avec des pirogues et des engins simples, continentale (fleuves, lacs) ou maritime (côte de Kribi à Limbé).
\end{itemize}

\textbf{Les transformations à retenir}
\begin{center}
\begin{tabularx}{\linewidth}{|l|X|X|}\hline
\rowcolor{popblueL} \textbf{Domaine} & \textbf{Avant (traditionnel)} & \textbf{Après (transformations)} \\ \hline
Espèces et races & Races locales rustiques mais peu productives (zébu, poules bicyclettes) & Sélection, croisements, introduction de races améliorées (pondeuses, porcs Large White, bovins métissés) \\ \hline
Alimentation & Divagation, pâturage libre, déchets & Provendes, fourrages cultivés, compléments minéraux, abreuvoirs \\ \hline
Soins & Absence de suivi sanitaire & Vaccination, déparasitage, encadrement vétérinaire \\ \hline
Habitat des animaux & Animaux en liberté, parcs de fortune & Bergeries, poulaillers, porcheries, ranches clôturés \\ \hline
Pêche & Pirogues à pagaie, filets rudimentaires, vente immédiate & Grandes pirogues à moteur hors-bord, filets adaptés (maillage réglementaire), glace, fumoirs améliorés, chambres froides \\ \hline
Production halieutique & Capture seule, ressource surexploitée & Pisciculture en étangs et cages, alevinage, aliments composés \\ \hline
\end{tabularx}
\end{center}

\textbf{L'action de l'État (Dossier 1)}
\begin{itemize}
  \item Encadrement : MINEPIA (élevage, pêche, industries animales), SODEPA (développement de l'élevage dans le Grand Nord), stations de recherche (Wakwa, Foumbot).
  \item Financement et programmes : \cle{PDMA} (Projet de développement des marchés agricoles), projets d'appui à la pisciculture, crédit agricole, subventions d'équipements.
  \item Formation des acteurs : centres de formation, vulgarisation agricole, organisations de producteurs (GIC, coopératives).
  \item Régulation de la commercialisation : organisation des filières, contrôle des prix et de la qualité, lutte contre les importations frauduleuses (poulets congelés, poissons importés).
\end{itemize}

\textbf{Méthodes express}
\begin{itemize}
  \item Croquis : localiser les grandes zones (élevage au Grand Nord et Adamaoua ; pêche sur la côte, le lac Tchad, le Logone, le Nyong ; pisciculture dans l'Ouest et le Centre), légende numérotée, titre, nord, échelle indicative.
  \item Commentaire de document : décrire (ce qu'on voit) $\rightarrow$ analyser (ce que cela montre : modernisation, contraintes) $\rightarrow$ conclure (bilan et limites).
  \item Calculs utiles : part en \% $= \dfrac{\text{partie}}{\text{total}} \times 100$ ; taux de variation $= \dfrac{V_f - V_i}{V_i} \times 100$.
\end{itemize}
\end{bilan}

\begin{aretenir}[Checklist avant le Bac]
\begin{itemize}
  \item $\square$ Je sais définir transhumance, bergerie, ranch et étang sans hésiter.
  \item $\square$ Je cite au moins quatre transformations de l'élevage : sélection et croisements, alimentation améliorée, soins vétérinaires, construction de bergeries et poulaillers.
  \item $\square$ Je localise les grandes zones d'élevage : Adamaoua, Nord, Extrême-Nord (et je sais pourquoi : climat soudanien, savanes, moins de glossines).
  \item $\square$ Je décris l'itinéraire de la transhumance (nord en saison des pluies, sud et bas-fonds en saison sèche) et ses problèmes (conflits agriculteurs-éleveurs).
  \item $\square$ Je cite les équipements modernes de la pêche : pirogues à moteur, filets adaptés, glace et fumoirs pour la conservation.
  \item $\square$ Je situe les zones de pêche : côte atlantique (Kribi, Limbé, Douala), lac Tchad, fleuves Logone, Bénoué, Sanaga.
  \item $\square$ J'explique l'essor de la pisciculture : étangs, espèces élevées (tilapia, carpe, clarias), intérêt face à la baisse des captures.
  \item $\square$ Je connais le rôle de l'État : PDMA, MINEPIA, SODEPA, formation des acteurs, régulation de la commercialisation.
  \item $\square$ Je sais construire un croquis légendé avec titre, nord et échelle indicative.
  \item $\square$ Je calcule correctement un pourcentage et un taux de variation à partir d'un tableau statistique.
  \item $\square$ Je relie la leçon à la famille de situations : consommer les produits locaux pour l'intégration nationale.
  \item $\square$ Je reste prudent avec les chiffres : je donne des ordres de grandeur et je cite ma source.
\end{itemize}
\end{aretenir}

\findecours

\end{document}
